% Options for packages loaded elsewhere
\PassOptionsToPackage{unicode}{hyperref}
\PassOptionsToPackage{hyphens}{url}
\PassOptionsToPackage{dvipsnames,svgnames,x11names,table}{xcolor}
\PassOptionsToPackage{space}{xeCJK}
\documentclass[
  chinese,
  a4paper,
  ,captions=tableheading
]{scrartcl}
\usepackage{xcolor}
\usepackage[margin=2.5cm,includehead=true,includefoot=true,centering,]{geometry}
\usepackage{amsmath,amssymb}


% add backlinks to footnote references, cf. https://tex.stackexchange.com/questions/302266/make-footnote-clickable-both-ways
\usepackage{footnotebackref}
\setcounter{secnumdepth}{5}
\usepackage{iftex}
\ifPDFTeX
  \usepackage[T1]{fontenc}
  \usepackage[utf8]{inputenc}
  \usepackage{textcomp} % provide euro and other symbols
\else % if luatex or xetex
  \usepackage{unicode-math} % this also loads fontspec
  \defaultfontfeatures{Scale=MatchLowercase}
  \defaultfontfeatures[\rmfamily]{Ligatures=TeX,Scale=1}
\fi
\usepackage{lmodern}
\ifPDFTeX\else
  % xetex/luatex font selection
  \setmainfont[]{Libertinus Serif}
  \setsansfont[]{Libertinus Sans}
  \setmonofont[]{Libertinus Mono}
  \setmathfont[]{Libertinus Math}
  \ifXeTeX
    \usepackage{xeCJK}
    \setCJKmainfont[]{Noto Serif CJK SC}
  \fi
  \ifLuaTeX
    \usepackage[]{luatexja-fontspec}
    \setmainjfont[]{Noto Serif CJK SC}
  \fi
\fi
% Use upquote if available, for straight quotes in verbatim environments
\IfFileExists{upquote.sty}{\usepackage{upquote}}{}
\IfFileExists{microtype.sty}{% use microtype if available
  \usepackage[]{microtype}
  \UseMicrotypeSet[protrusion]{basicmath} % disable protrusion for tt fonts
}{}

\usepackage{setspace}
\makeatletter
\@ifundefined{KOMAClassName}{% if non-KOMA class
  \IfFileExists{parskip.sty}{%
    \usepackage{parskip}
  }{% else
    \setlength{\parindent}{0pt}
    \setlength{\parskip}{6pt plus 2pt minus 1pt}}
}{% if KOMA class
  \KOMAoptions{parskip=half}}
\makeatother
\usepackage{longtable,booktabs,array}
\usepackage{calc} % for calculating minipage widths
% Correct order of tables after \paragraph or \subparagraph
\usepackage{etoolbox}
\makeatletter
\patchcmd\longtable{\par}{\if@noskipsec\mbox{}\fi\par}{}{}
\makeatother
% Allow footnotes in longtable head/foot
\IfFileExists{footnotehyper.sty}{\usepackage{footnotehyper}}{\usepackage{footnote}}
\makesavenoteenv{longtable}
\ifLuaTeX
\usepackage[bidi=basic,shorthands=off,,provide=*]{babel}
\else
\usepackage[bidi=default,shorthands=off,,provide=*]{babel}
\fi
\ifPDFTeX
\else
\babelfont{rm}[]{Libertinus Serif}
\fi
\ifLuaTeX
  \usepackage{selnolig} % disable illegal ligatures
\fi
\setlength{\emergencystretch}{3em} % prevent overfull lines
\providecommand{\tightlist}{%
  \setlength{\itemsep}{0pt}\setlength{\parskip}{0pt}}
% 中文段落排版：eisvogel 在 preamble 末尾无条件设置 \parindent=0pt，
% 因此必须在 \AtBeginDocument 中覆盖，否则首行缩进失效。
\usepackage{indentfirst}
\AtBeginDocument{%
  \setlength{\parindent}{2em}%
  \setlength{\parskip}{0pt}%
}
% 表格：略增行高，避免中文与数学符号贴边
\renewcommand{\arraystretch}{1.25}
% 中文断行与标点压缩（luatexja 分支生效后的细化设置）
\ltjsetparameter{jacharrange={-2,-3}}
\ltjsetparameter{xkanjiskip=0pt plus .4pt minus .4pt}
\usepackage{bookmark}
\IfFileExists{xurl.sty}{\usepackage{xurl}}{} % add URL line breaks if available
\urlstyle{same}
\definecolor{default-linkcolor}{HTML}{A50000}
\definecolor{default-filecolor}{HTML}{A50000}
\definecolor{default-citecolor}{HTML}{4077C0}
\definecolor{default-urlcolor}{HTML}{4077C0}

\hypersetup{
  pdftitle={信号与线性系统 · 期末复习资料},
  pdfauthor={课程：信号与线性系统（西北工业大学软件学院）},
  pdflang={zh-CN},
  pdfsubject={信号与线性系统},
  pdfkeywords={信号与线性系统, 傅里叶变换, 拉普拉斯变换, 系统函数, 复习},
  colorlinks=true,
  linkcolor={default-linkcolor},
  filecolor={default-filecolor},
  citecolor={default-citecolor},
  urlcolor={default-urlcolor},
  breaklinks=true,
  pdfcreator={LaTeX via pandoc with the Eisvogel template}}

\title{信号与线性系统 · 期末复习资料}
\usepackage{etoolbox}
\makeatletter
\providecommand{\subtitle}[1]{% add subtitle to \maketitle
  \apptocmd{\@title}{\par {\large #1 \par}}{}{}
}
\makeatother
\subtitle{考点清单 · 作业题型精解 · 公式速查}
\author{课程：信号与线性系统（西北工业大学软件学院）}
\date{2026 年秋}


%
% for the background color of the title page
%
\usepackage{pagecolor}
\usepackage{afterpage}
\usepackage[margin=2.5cm,includehead=true,includefoot=true,centering]{geometry}

%
% break urls
%
\PassOptionsToPackage{hyphens}{url}

%
% When using babel or polyglossia with biblatex, loading csquotes is recommended
% to ensure that quoted texts are typeset according to the rules of your main language.
%
\usepackage{csquotes}

%
% captions
%
\definecolor{caption-color}{HTML}{777777}
\usepackage{caption}
\captionsetup{font={stretch=1.2}, textfont={color=caption-color}, position=top, skip=4mm, labelfont=bf, singlelinecheck=false, justification=raggedright}
\setcapindent{0em}

%
% blockquote
%
\definecolor{blockquote-border}{RGB}{221,221,221}
\definecolor{blockquote-text}{RGB}{119,119,119}
\usepackage{mdframed}
\newmdenv[rightline=false,bottomline=false,topline=false,linewidth=3pt,linecolor=blockquote-border,skipabove=\parskip]{customblockquote}
\renewenvironment{quote}{\begin{customblockquote}\list{}{\rightmargin=0em\leftmargin=0em}%
\item\relax\color{blockquote-text}\ignorespaces}{\unskip\unskip\endlist\end{customblockquote}}

%
% Source Sans Pro as the default font family
% Source Code Pro for monospace text
%
% 'default' option sets the default
% font family to Source Sans Pro, not \sfdefault.
%
\ifnum 0\ifxetex 1\fi\ifluatex 1\fi=0 % if pdftex
    \usepackage[default]{sourcesanspro}
  \usepackage{sourcecodepro}
  \else % if not pdftex
    \fi

%
% heading color
%
\definecolor{heading-color}{RGB}{40,40,40}
\addtokomafont{section}{\color{heading-color}}
% When using the classes report, scrreprt, book,
% scrbook or memoir, uncomment the following line.
%\addtokomafont{chapter}{\color{heading-color}}

%
% variables for title, author and date
%
\usepackage{titling}
\title{信号与线性系统 · 期末复习资料}
\author{课程：信号与线性系统（西北工业大学软件学院）}
\date{2026 年秋}

%
% tables
%

\definecolor{table-row-color}{HTML}{F5F5F5}
\definecolor{table-rule-color}{HTML}{999999}

%\arrayrulecolor{black!40}
\arrayrulecolor{table-rule-color}     % color of \toprule, \midrule, \bottomrule
\setlength\heavyrulewidth{0.3ex}      % thickness of \toprule, \bottomrule
\renewcommand{\arraystretch}{1.5}     % spacing (padding)
\setlength{\tabcolsep}{10pt}          % spacing between columns


%
% remove paragraph indentation
%
\setlength{\parindent}{0pt}
\setlength{\parskip}{6pt plus 2pt minus 1pt}
\setlength{\emergencystretch}{3em}  % prevent overfull lines

%
%
% Listings
%
%


%
% header and footer
%
\usepackage[headsepline,footsepline]{scrlayer-scrpage}

\newpairofpagestyles{eisvogel-header-footer}{
  \clearpairofpagestyles
  \ihead*{信号与线性系统 · 期末复习资料}
  \chead*{}
  \ohead*{2026 年秋}
  \ifoot*{课程：信号与线性系统（西北工业大学软件学院）}
  \cfoot*{}
  \ofoot*{\thepage}
  \addtokomafont{pageheadfoot}{\upshape}
}
\pagestyle{eisvogel-header-footer}



%
% Define watermark
%

\begin{document}

\begin{titlepage}
		\newgeometry{left=6cm}
				\newcommand{\colorRule}[3][black]{\textcolor[HTML]{#1}{\rule{#2}{#3}}}
	\begin{flushleft}
		\noindent
		\\[-1em]
		\color[HTML]{5F5F5F}
		\makebox[0pt][l]{\colorRule[435488]{1.3\textwidth}{4pt}}
		\par
		\noindent

				{
				\setstretch{1.4}
				\vfill
				\noindent {\huge \textbf{\textsf{信号与线性系统 · 期末复习资料}}}
								\vskip 1em
					{\Large \textsf{考点清单 · 作业题型精解 · 公式速查}}
								\vskip 2em
				\noindent {\Large \textsf{课程：信号与线性系统（西北工业大学软件学院）}}
				\vfill
			}
		
		
				\textsf{2026 年秋}
			\end{flushleft}
\end{titlepage}
\restoregeometry
\pagenumbering{arabic}

% don't generate the default title
% \maketitle


{
				\hypersetup{linkcolor=}
				\setcounter{tocdepth}{3}
		\tableofcontents
				\newpage
			}
\setstretch{1.4}
\section*{前言：这条课程主线只有一句话}\label{ux524dux8a00ux8fd9ux6761ux8bfeux7a0bux4e3bux7ebfux53eaux6709ux4e00ux53e5ux8bdd}
\addcontentsline{toc}{section}{前言：这条课程主线只有一句话}

给定激励 \(f(t)\) 与线性时不变（LTI）系统，求响应 \(y(t)\)。

六个章节都在回答这同一个问题，区别只在于\textbf{用什么坐标去描述信号}。把这句话想通，全书的公式就不是一百个孤立结论，而是三套坐标下的同一件事：

\begin{itemize}
\tightlist
\item
  \textbf{时域坐标}：把信号拆成冲激的叠加，系统对 \(\delta(t)\) 的响应是
  \(h(t)\)，于是响应是叠加积分------卷积。
\item
  \textbf{频域坐标}：把信号拆成复正弦
  \(\mathrm{e}^{\mathrm{j}\omega t}\)
  的叠加（傅里叶级数对周期信号、傅里叶变换对非周期信号），系统对每个频率只做''乘一个复数''的加权，加权系数就是
  \(H(\mathrm{j}\omega)\)。
\item
  \textbf{复频域坐标}：把信号拆成 \(\mathrm{e}^{st}\)
  的叠加（拉普拉斯变换），系统对 \(\mathrm{e}^{st}\) 的响应仍是
  \(\mathrm{e}^{st}\) 乘一个数 \(H(s)\)，于是微分方程变成代数方程。
\end{itemize}

三套坐标通过''卷积定理''和''系统函数''缝合在一起：

\[
y(t)=f(t)*h(t)
\qquad\Longleftrightarrow\qquad
Y(\mathrm{j}\omega)=F(\mathrm{j}\omega)\,H(\mathrm{j}\omega)
\qquad\Longleftrightarrow\qquad
Y(s)=F(s)\,H(s)
\]

这三条等号是全部章节的骨架。考试里任何一道''求响应''的题，本质上都是在这三条里挑一条最短的路走。

\subsection*{六章知识地图}\label{ux516dux7ae0ux77e5ux8bc6ux5730ux56fe}
\addcontentsline{toc}{subsection}{六章知识地图}

\begin{longtable}[]{@{}
  >{\raggedright\arraybackslash}p{(\linewidth - 6\tabcolsep) * \real{0.1176}}
  >{\raggedright\arraybackslash}p{(\linewidth - 6\tabcolsep) * \real{0.2941}}
  >{\raggedright\arraybackslash}p{(\linewidth - 6\tabcolsep) * \real{0.2941}}
  >{\raggedright\arraybackslash}p{(\linewidth - 6\tabcolsep) * \real{0.2941}}@{}}
\toprule\noalign{}
\begin{minipage}[b]{\linewidth}\raggedright
章
\end{minipage} & \begin{minipage}[b]{\linewidth}\raggedright
回答的问题
\end{minipage} & \begin{minipage}[b]{\linewidth}\raggedright
核心工具
\end{minipage} & \begin{minipage}[b]{\linewidth}\raggedright
关键结论
\end{minipage} \\
\midrule\noalign{}
\endhead
\bottomrule\noalign{}
\endlastfoot
第一章 绪论 & 信号与系统各自怎么描述 & 信号的基本运算、奇异函数、LTI
判定 & 线性＝齐次性＋可加性；时不变＝参数不随时间变 \\
第二章 时域分析 & 不解变换，直接求响应 & 微分方程、算子 \(p\)、卷积 &
\(y(t)=y_{zi}(t)+y_{zs}(t)\)，\(y_{zs}(t)=f(t)*h(t)\) \\
第三章 正交分解 & 信号为什么能拆成三角函数之和 & 正交函数集、傅里叶级数
& 周期信号频谱离散；帕塞瓦尔定理给出功率分配 \\
第四章 频域分析 & 非周期信号怎么办 & 傅里叶变换、\(H(\mathrm{j}\omega)\)
& 时域卷积＝频域相乘；无失真传输与抽样定理的条件 \\
第五章 复频域分析 & 傅里叶变换不收敛的信号怎么办 &
拉普拉斯变换、部分分式、\(s\) 域电路 &
微分方程→代数方程；收敛域决定信号类型 \\
第六章 系统函数 & 不看输入，只看系统本身 & 零极点、罗斯判据 &
极点位置决定响应波形与稳定性 \\
\end{longtable}

一句话概括六章的关系：\textbf{第二、三、四、五章是四种''求响应的算法''，第六章是反过来''由响应看系统本质''。}

\subsection*{三条贯穿全程的恒等式}\label{ux4e09ux6761ux8d2fux7a7fux5168ux7a0bux7684ux6052ux7b49ux5f0f}
\addcontentsline{toc}{subsection}{三条贯穿全程的恒等式}

\begin{enumerate}
\def\labelenumi{\arabic{enumi}.}
\tightlist
\item
  \textbf{卷积定义}是时域一切结论的源头，卷积的四个性质（交换、分配、结合、与冲激/阶跃的关系）必须张口就来：
  \[
  f_1(t)*f_2(t)=\int_{-\infty}^{+\infty}f_1(\tau)f_2(t-\tau)\,\mathrm{d}\tau
  \]
\item
  \textbf{卷积定理}把时域最难算的运算变成乘法，这是第四章和第五章存在的全部理由。
\item
  \textbf{特征函数性质}说明 \(\mathrm{e}^{st}\) 是 LTI
  系统的特征函数，\(H(s)\) 是特征值： \[
  \mathrm{e}^{st}\ \longrightarrow\ H(s)\,\mathrm{e}^{st}
  \] 把 \(s=\mathrm{j}\omega\) 代进去就是频率响应，所以
  \(H(\mathrm{j}\omega)\) 只是 \(H(s)\)
  在虚轴上的截线------这一点把第五章和第四章连成一条线。
\end{enumerate}

\subsection*{建议的三轮复习路线}\label{ux5efaux8baeux7684ux4e09ux8f6eux590dux4e60ux8defux7ebf}
\addcontentsline{toc}{subsection}{建议的三轮复习路线}

\begin{longtable}[]{@{}
  >{\raggedright\arraybackslash}p{(\linewidth - 6\tabcolsep) * \real{0.2500}}
  >{\raggedright\arraybackslash}p{(\linewidth - 6\tabcolsep) * \real{0.2500}}
  >{\raggedright\arraybackslash}p{(\linewidth - 6\tabcolsep) * \real{0.2500}}
  >{\raggedright\arraybackslash}p{(\linewidth - 6\tabcolsep) * \real{0.2500}}@{}}
\toprule\noalign{}
\begin{minipage}[b]{\linewidth}\raggedright
轮次
\end{minipage} & \begin{minipage}[b]{\linewidth}\raggedright
时间
\end{minipage} & \begin{minipage}[b]{\linewidth}\raggedright
做什么
\end{minipage} & \begin{minipage}[b]{\linewidth}\raggedright
用本书哪部分
\end{minipage} \\
\midrule\noalign{}
\endhead
\bottomrule\noalign{}
\endlastfoot
第一轮 & 4 天 & 按章通读''考点清单 →
核心概念与公式''，把公式抄一遍并写下成立条件 & 各章前两节 \\
第二轮 & 3 天 &
只做''典型题型与解法''：先遮住解答自己做例题，再看同类题思路 &
各章第三节 \\
第三轮 & 1 天 & 考前只看附录：公式速查表 + 易错点清单 & 附录 B、附录
C \\
\end{longtable}

第二轮是提分的关键：本书每个题型都有''题型特征''一句话，考场上先靠这句话识别题型，再套步骤，比从头推导快得多。

\subsection*{全书记号约定}\label{ux5168ux4e66ux8bb0ux53f7ux7ea6ux5b9a}
\addcontentsline{toc}{subsection}{全书记号约定}

本书全篇统一使用下列记号，与课堂 PPT
一致。遇到与其他教材写法不同的地方，以此表为准。

\begin{longtable}[]{@{}
  >{\raggedright\arraybackslash}p{(\linewidth - 4\tabcolsep) * \real{0.3333}}
  >{\raggedright\arraybackslash}p{(\linewidth - 4\tabcolsep) * \real{0.3333}}
  >{\raggedright\arraybackslash}p{(\linewidth - 4\tabcolsep) * \real{0.3333}}@{}}
\toprule\noalign{}
\begin{minipage}[b]{\linewidth}\raggedright
记号
\end{minipage} & \begin{minipage}[b]{\linewidth}\raggedright
含义
\end{minipage} & \begin{minipage}[b]{\linewidth}\raggedright
说明
\end{minipage} \\
\midrule\noalign{}
\endhead
\bottomrule\noalign{}
\endlastfoot
\(y(t)\) & 系统的响应 & 课本与作业原卷记作
\(r(t)\)，本书为与各章统一改记 \(y(t)\)，两者等价 \\
\(u(t)\) & 单位阶跃函数 & 全篇用 \(u\)，不写作 \(\varepsilon(t)\) \\
\(\delta(t)\) & 单位冲激函数 & 强度为 1，面积为 1 \\
\(y_{zi}(t)\)、\(y_{zs}(t)\) & 零输入响应、零状态响应 & 下标分别表示
zero-input、zero-state \\
\(y_{n}(t)\)、\(y_{f}(t)\) & 自然响应、受迫响应 &
对应经典法的通解、特解 \\
\(h(t)\)、\(g(t)\) & 冲激响应、阶跃响应 & \(g(t)=h(t)*u(t)\) \\
\(\mathcal{F}[f(t)]=F(\mathrm{j}\omega)\) & 傅里叶变换 & 正变换用
\(\mathcal{F}\)，反变换用 \(\mathcal{F}^{-1}\) \\
\(\mathcal{L}[f(t)]=F(s)\) & 拉普拉斯变换 & 必须同时写出收敛域 \\
\(\mathrm{j}\) & 虚数单位 & 正体，不用 \(i\) \\
\(\mathrm{Sa}(x)=\dfrac{\sin x}{x}\) & 抽样函数 & 与
\(\mathrm{sinc}(x)=\dfrac{\sin \pi x}{\pi x}\) 相差一个 \(\pi\)
尺度，只在明确说明时混用 \\
\(\omega_{0}=\dfrac{2\pi}{T}\) & 周期信号的基频 & 第 3 章课件记作
\(\Omega\)，是同一个量；第 5、6 章极点的虚部另用
\(\omega_{i}\)，不与之混用 \\
\end{longtable}

\subsection*{怎么用这本书}\label{ux600eux4e48ux7528ux8fd9ux672cux4e66}
\addcontentsline{toc}{subsection}{怎么用这本书}

每章固定四节，含义如下：

\begin{itemize}
\tightlist
\item
  \textbf{考点清单}：一张表告诉你考什么、怎么考、多重。带 ★★★
  的必须能默写公式并独立做完例题。
\item
  \textbf{核心概念与公式}：结论 + 成立条件 +
  性质表。公式后面的括号条件不是装饰，漏了会算错。
\item
  \textbf{典型题型与解法}：每个题型给''题型特征 / 解题步骤 / 例题 /
  同类题''。例题取自附加作业并给完整解答，同类题只给关键思路------那些题请自己动手补全过程。
\item
  \textbf{易错点}：本节最短，但考前最值钱，都是作业里真实踩过的坑。
\end{itemize}

附录 A 是专题横向对比（同一问题在三个域里分别怎么做），附录 B
是考前一天的公式速查，附录 C 是易错点总汇。

\section{绪论}\label{ux7eeaux8bba}

\subsection{考点清单}\label{ux8003ux70b9ux6e05ux5355}

\begin{longtable}[]{@{}llc@{}}
\toprule\noalign{}
考点 & 常见考法 & 重要度 \\
\midrule\noalign{}
\endhead
\bottomrule\noalign{}
\endlastfoot
信号的分类 & 给波形或表达式，判连续/离散、确定/随机 & ★★ \\
周期性与周期 & 给两个以上正弦叠加，判周期性并求\(T\) & ★★★ \\
能量信号与功率信号 & 算总能量\(E\)、平均功率\(P\)，判类型 & ★★ \\
信号的简单处理 & 由\(f(t)\)画\(f(at+b)\)的波形，说明变换顺序 & ★★★ \\
奇异函数 & 用\(u(t)\)写波形表达式；用取样特性算积分 & ★★★ \\
线性与时不变性判定 & 给方程判线性、判时不变 & ★★★ \\
LTI 系统的分析路线 & 概念题：时域法与变换域法各解决什么 & ★ \\
\end{longtable}

\subsection{核心概念与公式}\label{ux6838ux5fc3ux6982ux5ff5ux4e0eux516cux5f0f}

\subsubsection{信号的分类（四条规则）}\label{ux4fe1ux53f7ux7684ux5206ux7c7bux56dbux6761ux89c4ux5219}

信号是信息的承载方式，是随时间变化的物理量。四条分类规则互不排斥，一道题可能同时问两条。

\begin{longtable}[]{@{}
  >{\raggedright\arraybackslash}p{(\linewidth - 4\tabcolsep) * \real{0.2500}}
  >{\raggedright\arraybackslash}p{(\linewidth - 4\tabcolsep) * \real{0.5000}}
  >{\raggedright\arraybackslash}p{(\linewidth - 4\tabcolsep) * \real{0.2500}}@{}}
\toprule\noalign{}
\begin{minipage}[b]{\linewidth}\raggedright
规则
\end{minipage} & \begin{minipage}[b]{\linewidth}\raggedright
分出的两类
\end{minipage} & \begin{minipage}[b]{\linewidth}\raggedright
判据
\end{minipage} \\
\midrule\noalign{}
\endhead
\bottomrule\noalign{}
\endlastfoot
时间函数的确定性 & 确定信号 / 随机信号 & 能否用确定的时间函数写出 \\
自变量取值是否连续 & 连续信号 / 离散信号 & 自变量\(t\)是否连续取值 \\
是否存在周期现象 & 周期信号 / 非周期信号 &
是否存在满足定义的最小正周期 \\
时间范围 & 因果信号 / 非因果信号 & \(t<0\)时是否恒为零 \\
\end{longtable}

连续时间信号（CTS）：除若干不连续的时间点外，每个时间点上都有对应的信号值，如
\(f(t)=1+\sin\omega t\)。离散时间信号（DTS）：信号只在某些不连续的时间点上有值，其余时间点
无定义，如\(f(n)\)。

因果信号：\(\forall t\in(-\infty,0)\)，\(f(t)\equiv 0\)。\textbf{什么时候用}：讨论单边拉普拉斯变
换、判断卷积积分下限时都要先看因果性。

\subsubsection{周期性与周期}\label{ux5468ux671fux6027ux4e0eux5468ux671f}

连续周期信号的定义与周期：

\[ \forall t\in(-\infty,+\infty),\quad f(t+mT)=f(t),\quad m=0,\pm 1,\pm 2,\cdots \]

满足上式的最小正数\(T\)称为周期。离散情形把自变量换成整数\(n\)、把周期记作正整数\(N\)：

\[ \forall n\in(-\infty,+\infty),\quad f(n+mN)=f(n),\quad m=0,\pm 1,\pm 2,\cdots \]

满足上式的最小正整数\(N\)称为周期。全篇离散信号的自变量一律用\(n\)（\(f(n)\)、\(u(n)\)、
\(\delta(n)\)、\(f(n-n_{0})\) 等）；求和式的哑指标用\(k\)，例如
\(\displaystyle\sum_{k}f(k)\delta(n-k)=f(n)\)，注意\(k\)只是哑指标、\(n\)才是自变量。

复合周期信号的周期判据（本条是本课程第一章最爱考的结论）：

\begin{itemize}
\tightlist
\item
  两个周期信号的周期之比为\textbf{有理数}时，其和才是周期信号；
\item
  和信号的周期等于两个周期的最小公倍数；
\item
  若周期之比为无理数，和信号不是周期信号。
\end{itemize}

\subsubsection{能量信号与功率信号}\label{ux80fdux91cfux4fe1ux53f7ux4e0eux529fux7387ux4fe1ux53f7}

信号的总能量与平均功率（区间取对称的\([-T,T]\)）：

\[ E=\lim_{T\to+\infty}\int_{-T}^{T}\left|f(t)\right|^{2}\mathrm{d}t,\qquad
P=\lim_{T\to+\infty}\frac{1}{2T}\int_{-T}^{T}\left|f(t)\right|^{2}\mathrm{d}t \]

\begin{itemize}
\tightlist
\item
  能量信号：\(0<E<+\infty\)；此时必有\(P=0\)。
\item
  功率信号：\(0<P<+\infty\)；此时必有\(E\to+\infty\)。
\item
  两者互不相容，但\textbf{并非非此即彼}：既非能量信号又非功率信号的信号是存在的（如
  \(f(t)=\dfrac{1}{t}u(t-1)\)）。
\end{itemize}

两个必须背下来的结论：

\begin{longtable}[]{@{}lccl@{}}
\toprule\noalign{}
信号 & 总能量\(E\) & 平均功率\(P\) & 类型 \\
\midrule\noalign{}
\endhead
\bottomrule\noalign{}
\endlastfoot
\(f(t)=1,\ 0<t<1\)（其余为\(0\)） & \(1\) & \(0\) & 能量信号 \\
\(f(t)=A\sin\omega t\) & \(\to\infty\) & \(A^{2}/2\) & 功率信号 \\
\end{longtable}

\subsubsection{信号的简单处理}\label{ux4fe1ux53f7ux7684ux7b80ux5355ux5904ux7406}

四种基本处理（都以\(t\)为操作对象，不作为的常数不要跟着乱动）：

\begin{longtable}[]{@{}
  >{\raggedright\arraybackslash}p{(\linewidth - 4\tabcolsep) * \real{0.3000}}
  >{\raggedright\arraybackslash}p{(\linewidth - 4\tabcolsep) * \real{0.4000}}
  >{\raggedright\arraybackslash}p{(\linewidth - 4\tabcolsep) * \real{0.3000}}@{}}
\toprule\noalign{}
\begin{minipage}[b]{\linewidth}\raggedright
处理
\end{minipage} & \begin{minipage}[b]{\linewidth}\raggedright
表达式
\end{minipage} & \begin{minipage}[b]{\linewidth}\raggedright
效果
\end{minipage} \\
\midrule\noalign{}
\endhead
\bottomrule\noalign{}
\endlastfoot
相加 & \(f(t)=f_{1}(t)+f_{2}(t)\) & 同一时刻函数值代数相加 \\
相乘 & \(f(t)=f_{1}(t)\times f_{2}(t)\) & 同一时刻函数值相乘 \\
平移 & \(f(t)\to f(t-t_{0})\) & \(t_{0}>0\)右移，\(t_{0}<0\)左移 \\
反褶 & \(f(t)\to f(-t)\) & 以纵轴为轴翻转 \\
尺度变换 & \(f(t)\to f(at)\) &
\(\left|a\right|>1\)压缩，\(\left|a\right|<1\)展宽 \\
标量乘法 & \(f(t)\to af(t)\) & 只改幅度，不改位置 \\
\end{longtable}

\textbf{混合运算的顺序}（本课程强调的操作要点）：分步次序可以不同，但每一步的处理都是针对变量
\(t\)进行的，不是针对括号里的整块式子。两种标准顺序：

\begin{enumerate}
\def\labelenumi{\arabic{enumi}.}
\tightlist
\item
  先尺度、再反褶、再平移、最后标量相乘。
\item
  先平移、再反褶、再尺度、最后标量相乘。
\end{enumerate}

顺序不同则平移量不同：\textbf{先平移再尺度压缩\(a\)倍时，平移量要除以\(a\)}；先尺度后平移则直接平
移给定数值。以课件练习\(f(1-2t)\)为例，两种做法都能得到正确结果：

\[ f(t)\ \longrightarrow\ f(t+1)\ \longrightarrow\ f(-t+1)\ \longrightarrow\ f(-2t+1) \]

依次为左移\(1\)、反褶、压缩\(2\)倍。

\[ f(t)\ \longrightarrow\ f(-t)\ \longrightarrow\ f(-2t)\ \longrightarrow\ f\!\left(-2(t-0.5)\right)=f(-2t+1) \]

依次为反褶、压缩\(2\)倍、右移\(0.5\)。

验证手段只有一个：\textbf{抓关键点}。若已知\(f(t)\)在\(t=t_{k}\)处的值，先解出使括号内等于\(t_{k}\)
的那个\(t\)，再逐点搬移。这一步比''跟着感觉画''可靠得多。

\subsubsection{奇异函数：阶跃与冲激}\label{ux5947ux5f02ux51fdux6570ux9636ux8dc3ux4e0eux51b2ux6fc0}

单位阶跃函数（全篇记\(u(t)\)，课件写作\(\varepsilon(t)\)，是同一个函数）：

\[ u(t)=\begin{cases}1, & t>0\\ 0, & t<0\end{cases} \]

\(t=0\)处的取值在工程应用中不重要，常约定为\(1\)、\(1/2\)或不予定义，\textbf{不影响任何积分与变换的
结果}------凡遇到\(u(t)\)出现在被积函数或变换核里，都按\(t=0\)是零测度点处理即可。工程背景是开关
模型；用途是\textbf{限定函数的时间范围}。

单位冲激函数：持续时间无穷小、瞬间幅度无穷大、涵盖面积恒为\(1\)的理想信号。

\[ \delta(t)=0\ (t\ne 0),\qquad \delta(t)=\infty\ (t=0),\qquad
\int_{-\infty}^{+\infty}\delta(\tau)\,\mathrm{d}\tau=1 \]

强度为\(A\)、发生在\(t_{0}\)的冲激记作\(A\delta(t-t_{0})\)；图形只画''位置
+ 强度''两件事。

冲激函数的六条性质：

\begin{longtable}[]{@{}
  >{\raggedright\arraybackslash}p{(\linewidth - 4\tabcolsep) * \real{0.3000}}
  >{\raggedright\arraybackslash}p{(\linewidth - 4\tabcolsep) * \real{0.4000}}
  >{\raggedright\arraybackslash}p{(\linewidth - 4\tabcolsep) * \real{0.3000}}@{}}
\toprule\noalign{}
\begin{minipage}[b]{\linewidth}\raggedright
性质
\end{minipage} & \begin{minipage}[b]{\linewidth}\raggedright
表达式
\end{minipage} & \begin{minipage}[b]{\linewidth}\raggedright
说明
\end{minipage} \\
\midrule\noalign{}
\endhead
\bottomrule\noalign{}
\endlastfoot
时移 & \(\delta(t-t_{0})\)在\(t=t_{0}\)处面积仍为\(1\) &
位置变了，强度不变 \\
偶函数 & \(\delta(t)=\delta(-t)\) & 反褶后图形不变 \\
与阶跃的关系 &
\(\delta(t)=\dfrac{\mathrm{d}}{\mathrm{d}t}u(t)\)，\(u(t)=\displaystyle\int_{-\infty}^{t}\delta(\tau)\mathrm{d}\tau\)
& 阶跃的导数是冲激 \\
筛选 & \(f(t)\delta(t-t_{0})=f(t_{0})\delta(t-t_{0})\) &
与\(f\)相乘只保留\(t_{0}\)处的值 \\
取样 &
\(\displaystyle\int_{-\infty}^{\infty}f(t)\delta(t-t_{0})\mathrm{d}t=f(t_{0})\)
& 乘积的面积等于冲激处的函数值 \\
定积分限 &
区间含\(t_{0}\)时\(\int_{\alpha}^{\beta}f(t)\delta(t-t_{0})\mathrm{d}t=f(t_{0})\)
& 区间不含\(t_{0}\)时结果为\(0\) \\
\end{longtable}

取样特性是''算积分''而不是''算函数值''------最后的载体是数，不是\(t\)的函数，这是判卷时最常见的
丢分点。

\subsubsection{系统、线性与时不变性}\label{ux7cfbux7edfux7ebfux6027ux4e0eux65f6ux4e0dux53d8ux6027}

系统是由若干相互联系的单元组成、具有某种功能、用以达到某种目的的有机整体。描述方法有输入
输出方程、框图、状态变量三种。

线性（线性 = 齐次性 + 叠加性）：

\begin{itemize}
\tightlist
\item
  齐次性：\(e(t)\to y(t)\)则\(ke(t)\to ky(t)\)；
\item
  叠加性：\(e_{1}\to y_{1}\)、\(e_{2}\to y_{2}\)则\(e_{1}+e_{2}\to y_{1}+y_{2}\)。
\end{itemize}

等价判据：\textbf{先线性运算再经过系统 = 先经过系统再线性运算}。

时不变：由\(e(t)\to y(t)\)必能推出\(e(t-t_{0})\to y(t-t_{0})\)。等价判据：\textbf{先时移再经过系统
= 先经过系统再时移}，即''输出响应与输入施加的起始时刻无关''。

两个必须记住的特例：

\begin{itemize}
\tightlist
\item
  \(y(t)=e(t)+C\ (C\ne 0)\)严格说不是线性系统，但工程上常见，按\textbf{增量线性}处理：只要零状
  态响应满足线性即可。等价地，把响应拆成零输入响应与零状态响应，只要求零状态响应线性。
\item
  乘法器\(y(t)=e_{1}(t)e_{2}(t)\)严格说也不是线性系统，但通信中用得多，本课程仍作研究对象。
\end{itemize}

零输入响应：输入为零时仅由初始储能引起的响应。零状态响应：初始储能为零时仅由输入引起的响应。

\subsubsection{LTI
系统的分析路线}\label{lti-ux7cfbux7edfux7684ux5206ux6790ux8defux7ebf}

LTI（Linear
Time-Invariant）系统是本课程的主线：这类系统大量存在、已有完整分析方法、易于
综合实现。分析分三步：建立数学模型 → 进行分析 → 物理解释。

\begin{itemize}
\tightlist
\item
  时域法：经典法（通解 + 特解）、近代时域法（零输入 + 零状态，卷积）。
\item
  变换域法：频域（傅里叶变换）、复频域（拉普拉斯变换）。
\item
  状态方程法：多输入多输出、便于计算机求解。
\end{itemize}

时域法结果直观但求解困难；变换域法把微分方程变成代数方程，但要两次变换、计算量大。第
2章 走时域法，第 3--6 章走变换域法。

\subsection{典型题型与解法}\label{ux5178ux578bux9898ux578bux4e0eux89e3ux6cd5}

\subsubsection{题型 1-1
判断信号的类型（连续/离散、确定/随机、能量/功率）}\label{ux9898ux578b-1-1-ux5224ux65adux4fe1ux53f7ux7684ux7c7bux578bux8fdeux7eedux79bbux6563ux786eux5b9aux968fux673aux80fdux91cfux529fux7387}

\textbf{题型特征}：给出一组波形或表达式，要求逐条说明属于哪一类信号；或给定信号求\(E\)、\(P\)并判
断类型。

\textbf{解题步骤}： 1. 看自变量：\(t\)连续取值 →
连续信号；只在离散点上取值（竖线、点列）→ 离散信号。波形是
否''平滑''、是否''有跳变''与连续/离散无关。 2.
看能否写出确定的函数式：能 → 确定信号；只能给统计特性 → 随机信号。 3.
算\(E\)与\(P\)：先用分段表达式写出\(\left|f\right|^{2}\)的积分区间，再取极限。
4.
按\(0<E<\infty\)、\(0<P<\infty\)分类；两者都不满足时如实写''既非能量信号也非功率信号''。

\textbf{例题}（【附加作业1 第1题】）：图 1 给出 (a)--(f)
六个子图，要求判断各信号是连续信号还 是离散信号。

解答过程：判据只有一条------\textbf{时间变量\(t\)是否连续取值}。

\begin{enumerate}
\def\labelenumi{\arabic{enumi}.}
\tightlist
\item
  (a)、(b)、(d)、(e)
  四个子图的横轴是连续的时间轴，曲线上每一点都有对应的信号值（其中 (
  b)、(e)
  的波形存在跳变点，只影响''连续性''的措辞，不影响''连续时间信号''的判断），故为连
  续信号。
\item
  (c)、(f) 两个子图只在离散的时刻取值，其余时刻无定义，故为离散信号。
\end{enumerate}

结论：(a)、(b)、(d)、(e) 为连续信号；(c)、(f) 为离散信号。

补充：若把''是否存在跳变''当成判据，会把 (b)、(e)
误判为离散信号。正确提法是''连续时间信号
除若干不连续的时间点外，每个时间点上都有对应的信号值''，跳变点被这句''除外''覆盖了。

\subsubsection{题型 1-2
判断周期性并求周期}\label{ux9898ux578b-1-2-ux5224ux65adux5468ux671fux6027ux5e76ux6c42ux5468ux671f}

\textbf{题型特征}：给出若干正弦（余弦）信号之和，问是否为周期信号；若是，求周期\(T\)。

\textbf{解题步骤}： 1.
把每一项写成\(A\sin(\omega_{i}t)\)或\(A\cos(\omega_{i}t)\)的标准形式，读出\(\omega_{i}\)。
2. 求每一项的周期\(T_{i}=2\pi/\omega_{i}\)。 3.
判据：各\(T_{i}\)两两之比为有理数 ⇒ 是周期信号；出现无理数比 ⇒
非周期信号。 4.
求\(T\)：找最小的一组正整数\(m_{i}\)使\(m_{1}T_{1}=m_{2}T_{2}=\cdots=T\)，这个公共值就是周
期（等价说法：各周期的最小公倍数）。 5.
检验：把\(T\)除以每个\(T_{i}\)都必须是正整数。

\textbf{例题}（【附加作业1
第2题】）：判断下列信号是否为周期信号，若是求周期\(T\)。

1）\(a\sin t-b\sin(3t)\)；2）\(a\sin(4t)+b\cos(7t)\)；3）
\(a\sin\dfrac{5t}{2}+b\cos\dfrac{6t}{5}+c\sin\dfrac{t}{7}\)。

解答：

1）\(\omega_{1}=1\)，\(\omega_{2}=3\)，故\(T_{1}=2\pi\)，\(T_{2}=\dfrac{2\pi}{3}\)。比值
\(T_{1}/T_{2}=3\)为有理数，是周期信号。取\(m_{1}=1\)、\(m_{2}=3\)，得\(T=2\pi\)。检验：
\(T/T_{1}=1\)，\(T/T_{2}=3\)，均为正整数。

2）\(\omega_{1}=4\)，\(\omega_{2}=7\)，故\(T_{1}=\dfrac{\pi}{2}\)，\(T_{2}=\dfrac{2\pi}{7}\)。比
值\(T_{1}/T_{2}=\dfrac{7}{4}\)为有理数，是周期信号。取\(m_{1}=4\)、\(m_{2}=7\)，得
\(T=4T_{1}=7T_{2}=2\pi\)。

3）\(\omega_{1}=\dfrac{5}{2}\)，\(\omega_{2}=\dfrac{6}{5}\)，\(\omega_{3}=\dfrac{1}{7}\)，故

\[ T_{1}=\frac{4\pi}{5},\qquad T_{2}=\frac{5\pi}{3},\qquad T_{3}=14\pi \]

三周期之比均为有理数，是周期信号。求公共倍数：把三个周期同乘以\(\dfrac{15}{\pi}\)，得
\(12,\ 25,\ 210\)；而

\[ 12=2^{2}\cdot 3,\quad 25=5^{2},\quad 210=2\cdot 3\cdot 5\cdot 7
\ \Rightarrow\ \mathrm{lcm}=2^{2}\cdot 3\cdot 5^{2}\cdot 7=2100 \]

故\(T=2100\cdot\dfrac{\pi}{15}=140\pi\)。检验：\(\dfrac{140\pi}{4\pi/5}=175\)，
\(\dfrac{140\pi}{5\pi/3}=84\)，\(\dfrac{140\pi}{14\pi}=10\)，均为正整数。

三个小题的答案依次为\(T=2\pi\)、\(T=2\pi\)、\(T=140\pi\)。

\subsubsection{题型 1-3
波形的平移、反褶与尺度变换}\label{ux9898ux578b-1-3-ux6ce2ux5f62ux7684ux5e73ux79fbux53cdux8936ux4e0eux5c3aux5ea6ux53d8ux6362}

\textbf{题型特征}：已知\(f(t)\)的波形，要求画出\(f(at+b)\)或\(af(at+b)+c\)的波形，并说明变换顺序。

\textbf{解题步骤}： 1.
先把目标式写成\(f\big(a(t-t_{0})\big)\)的\textbf{标准形}，读出压缩比\(a\)与平移量\(t_{0}\)。
2. 按''尺度 → 反褶 → 平移 → 标量相乘''逐步作图画三张中间图。 3.
用关键点搬移核对：若原图有拐点\((t_{k},f_{k})\)，则新图对应点为
\((t_{0}+t_{k}/a,\ af_{k})\)。 4.
写出最终分段表达式，标清每段端点坐标与斜率。

\textbf{例题}（【附加作业1
第3题】）：已知\(x(t)\)是支撑在\((-1,1]\)上的斜线\(x(t)=t+1\)（由
\((-1,0)\) 直线升到\((1,2)\)，\(t>1\)与\(t\le -1\)处为\(0\)），画出
\(y(t)=2x\left(-\dfrac{1}{3}t+\dfrac{2}{3}\right)\)的波形。

解答：先把括号内提出因子\(-\dfrac{1}{3}\)：

\[ -\frac{1}{3}t+\frac{2}{3}=-\frac{1}{3}(t-2) \]

所以\(y(t)=2x\left[-\dfrac{1}{3}(t-2)\right]\)，即\(a=-\dfrac{1}{3}\)、\(t_{0}=2\)。按题目要求的
顺序分四步，每一步都只对变量\(t\)操作：

\begin{enumerate}
\def\labelenumi{\arabic{enumi}.}
\tightlist
\item
  \textbf{尺度变换}：\(x(t)\to x\left(\dfrac{t}{3}\right)\)。因\(\left|a\right|=\dfrac{1}{3}<1\)，
  支撑区\((-1,1]\)被展宽为\((-3,3]\)，端点\((-1,0)\)、\((1,2)\)变成\((-3,0)\)、\((3,2)\)，纵向不变。
\item
  \textbf{反褶}：\(x\left(\dfrac{t}{3}\right)\to x\left(-\dfrac{t}{3}\right)\)，以纵轴为轴翻转，
  支撑区变为\([-3,3)\)，斜线由''上升''变''下降''：左端\((3,0)\)、右端\((-3,2)\)。
\item
  \textbf{平移}：把\(t\)换成\(t-2\)，得\(x\left(-\dfrac{t-2}{3}\right)\)，整体\textbf{右移
  2}，支撑区变为 \([-1,5)\)，端点变为\((5,0)\)与\((-1,2)\)。
\item
  \textbf{标量相乘}：幅度乘\(2\)，端点变为\((5,0)\)与\((-1,4)\)，得\(y(t)\)。
\end{enumerate}

关键点核对：映射关系为\((t_{k},x_{k})\to\left(2-3t_{k},\ 2x_{k}\right)\)。把原图三个关键点
逐个代入：

\[ (-1,0)\to(5,0),\qquad (0,1)\to(2,2),\qquad (1,2)\to(-1,4) \]

三点共线，故\(y(t)\)是支撑在\([-1,5)\)上的下降直线。这个映射本身也说明了反褶的效果：原图右端
\(t=1\)映到新图左端\(t=-1\)。

结果写成闭式便于核对：

\[ y(t)=2x\!\left(-\frac{t}{3}+\frac{2}{3}\right)=\frac{10-2t}{3},\qquad -1\le t<5 \]

\(t=-1\)处\(y=4\)、\(t=2\)处\(y=2\)、\(t=5\)处\(y=0\)，与原斜线\(x(t)=t+1\)逐点对应。

若改按''平移 → 反褶 →
尺度''操作，则第二步的平移量是\(+2\)、第三步压缩\(1/3\)时平移量变成
\(2\times 3=6\)的等效右移，最终仍落在\(t=2\)。两种顺序结果相同，但中间图不同，\textbf{答题时写清自
己用的是哪一种}。

\textbf{同类题}： - 【附加作业1 第10题】图 4
的\(f(t)\)是支撑在\((-1,1]\)上的斜线\(f(t)=\dfrac{t+1}{2}\)（由
\((-1,0)\)升到\((1,1)\)），要求画六个波形：\(f(2t)\)、\(f(t)u(t)\)、\(f(t-2)u(t)\)、
\(f(t-2)u(t-2)\)、\(f(2-t)\)、\(f(-2-t)u(-t)\)。统一用关键点映射：\(f(at+b)\)把\((t_{k},f_{k})\)
搬到\(\left((t_{k}-b)/a,\ f_{k}\right)\)，再按乘上的\(u(\cdot)\)的截断条件剪掉另一半。六个
答案的关键点依次为\((-0.5,0)\to(0.5,1)\)；\((0,0.5)\to(1,1)\)；\((1,0)\to(3,1)\)；
\((2,0.5)\to(3,1)\)；\((1,1)\to(3,0)\)；\((-3,1)\to(-1,0)\)。

\subsubsection{题型 1-4
用阶跃函数表示波形；用冲激取样特性求积分}\label{ux9898ux578b-1-4-ux7528ux9636ux8dc3ux51fdux6570ux8868ux793aux6ce2ux5f62ux7528ux51b2ux6fc0ux53d6ux6837ux7279ux6027ux6c42ux79efux5206}

\textbf{题型特征}：两种小计算常一起考------(i)
给分段波形，要求''不能用分段函数、必须用阶跃函数'' 写表达式；(ii)
给含\(\delta\)的定积分，要求算数值。

\textbf{解题步骤（阶跃表示）}： 1.
从\(t=-\infty\)向右扫描，记下每一处跳变的位置\(t_{k}\)与跳变量\(\Delta_{k}\)（上跳为正）。
2. 累加得\(f(t)=\sum_{k}\Delta_{k}u(t-t_{k})\)，这就是''跳变叠加法''。
3. 若某段是斜坡而不是跳变，把该段拆成''跳变 +
斜率\(\times\)斜坡''两项，或对\(f'(t)\)用跳变叠 加法再积分。 4.
门函数型：\(f(t)=A\left[u(t-t_{1})-u(t-t_{2})\right]\)，只在\(t_{1}<t<t_{2}\)内取\(A\)。

\textbf{解题步骤（取样积分）}： 1.
找出冲激所在时刻\(t_{0}\)与强度\(A\)。 2. 把\(\delta\)换成''取样''：
\(\displaystyle\int_{\alpha}^{\beta}f(t)A\delta(t-t_{0})\mathrm{d}t=A f(t_{0})\)，\textbf{前提是
\(\alpha<t_{0}<\beta\)}。 3. 若\(t_{0}\)落在积分区间之外，直接得\(0\)。
4.
若\(f(t)\)在\(t_{0}\)处是不定式（如\(\dfrac{\sin 2t}{t}\)在\(t=0\)），先用筛选特性
\(f(t)\delta(t-t_{0})=f(t_{0})\delta(t-t_{0})\)把\(f\)换成它在\(t_{0}\)处的极限值，再取样。
5. 结果是数，不带\(t\)。

\textbf{例题}（【附加作业1 第8题】）：利用冲激函数的取样特性求下列积分。

\[ \int_{-\infty}^{+\infty}\delta(t-2)\sin t\,\mathrm{d}t,\qquad
\int_{-\infty}^{+\infty}\frac{\sin 2t}{t}\delta(t)\,\mathrm{d}t,\qquad
\int_{0}^{+\infty}\mathrm{e}^{-2t}\delta(t+2)\,\mathrm{d}t \]

解答：

\begin{enumerate}
\def\labelenumi{(\arabic{enumi})}
\tightlist
\item
  冲激位于\(t_{0}=2\)、强度为\(1\)；\(f(t)=\sin t\)在\(t=2\)处连续，且\(2\in(-\infty,+\infty)\)，
  可以直接取样：
\end{enumerate}

\[ \int_{-\infty}^{+\infty}\delta(t-2)\sin t\,\mathrm{d}t=\sin 2 \]

\begin{enumerate}
\def\labelenumi{(\arabic{enumi})}
\setcounter{enumi}{1}
\tightlist
\item
  冲激位于\(t_{0}=0\)，\(0\in(-\infty,+\infty)\)；但\(f(t)=\dfrac{\sin 2t}{t}\)在\(t=0\)处无定
  义，必须先取极限。把它改写成抽样函数：
\end{enumerate}

\[ \frac{\sin 2t}{t}=2\cdot\frac{\sin 2t}{2t}=2\,\mathrm{Sa}(2t),
\qquad \mathrm{Sa}(x)=\frac{\sin x}{x} \]

这里\(\mathrm{Sa}(x)\)就是抽样函数（与 \(\mathrm{sinc}\)
的定义等价，本章全用此记号）。
因\(\displaystyle\lim_{t\to 0}\mathrm{Sa}(2t)=1\)，故\(f(0)=2\)，于是

\[ \int_{-\infty}^{+\infty}\frac{\sin 2t}{t}\delta(t)\,\mathrm{d}t
=\int_{-\infty}^{+\infty}2\delta(t)\,\mathrm{d}t=2 \]

\begin{enumerate}
\def\labelenumi{(\arabic{enumi})}
\setcounter{enumi}{2}
\tightlist
\item
  冲激位于\(t_{0}=-2\)，而积分区间是\([0,+\infty)\)，\(-2\notin[0,+\infty)\)，冲激点不在区间
  内，故
\end{enumerate}

\[ \int_{0}^{+\infty}\mathrm{e}^{-2t}\delta(t+2)\,\mathrm{d}t=0 \]

三个答案依次为\(\sin 2\)、\(2\)、\(0\)。第 (3)
小题正是''先看积分限、再看被积函数''的典型：若只
顾取样而写成\(\mathrm{e}^{-2\cdot(-2)}=\mathrm{e}^{4}\)，就完全错了。

\textbf{同类题}： - 【附加作业1 第9题】图 3
三个波形用阶跃函数写表达式：(a) 在\([0,T]\)内四级阶梯上升、
\(t=T\)后归零；(b)
两个宽为\(\tau\)的矩形脉冲，分别起于\(t_{1}\)与\(t_{2}\)；(c)
两个半波整流的 正弦。一律用跳变叠加法定系数，答案为

\[ \begin{aligned}
\text{(a)}\quad i(t)&=\frac{I}{4}\left[u(t)+u\left(t-\frac{T}{4}\right)
+u\left(t-\frac{T}{2}\right)+u\left(t-\frac{3T}{4}\right)\right]\\
&\quad -I\,u(t-T)
\end{aligned} \]

\[ \text{(b)}\quad i(t)=I\left[u(t-t_{1})-u(t-t_{1}-\tau)\right]
+I\left[u(t-t_{2})-u(t-t_{2}-\tau)\right] \]

\[ \text{(c)}\quad i(t)=I\sin t\left[u(t)-u(t-\pi)\right]
+I\sin(t-\pi)\left[u(t-\pi)-u(t-2\pi)\right] \]

三问都用''每级跳变量 \(\times\)
该级起始时刻的阶跃''叠加，最后再减去截止时刻的阶跃； (c)
的第二项必须写成 \(\sin(t-\pi)\) 而不是 \(\sin t\)，否则平移后相位不对。

\subsubsection{题型 1-5
判断系统的线性与时不变性}\label{ux9898ux578b-1-5-ux5224ux65adux7cfbux7edfux7684ux7ebfux6027ux4e0eux65f6ux4e0dux53d8ux6027}

\textbf{题型特征}：给出一个输入输出方程（或框图），问''是否具备线性、是否具备时不变性''，两问必
须分别作答。

\textbf{解题步骤}： 1.
\textbf{先扫描非线性项}：把方程整理成''\(y(t)\)及其各阶导数 =
\(e(t)\)及其各阶导数''后，逐项检查
有没有\(y\)与\(e\)的乘积、\(y\)的平方、\(\left|y\right|\)、\(y\cdot y'\)这类项。有一项就足以判非
线性。 2.
\textbf{齐次性检验}：设\(e(t)\to y(t)\)，把\(ky(t)\)代回左端，看是否恰好等于\(k e(t)\)；不等则齐
次性不成立。 3.
\textbf{叠加性检验}：设\(e_{1}\to y_{1}\)、\(e_{2}\to y_{2}\)，看\(e_{1}+e_{2}\)的响应是否等于
\(y_{1}+y_{2}\)。 4.
\textbf{时不变检验}：看所有系数是否显含\(t\)；再用''先时移再经过系统 =
先经过系统再时移''验证。 5.
结论分两句写：``具备/不具备线性；具备/不具备时不变性''。

\textbf{例题}（【附加作业1 第4题】）：判断方程

\[ \frac{\mathrm{d}^{2}y(t)}{\mathrm{d}t^{2}}+2y(t)\frac{\mathrm{d}y(t)}{\mathrm{d}t}-3y(t)=e(t) \]

描述的系统是否具备线性、是否具备时不变性。

解答（原卷把响应记作 r，本资料按统一记号写 y，下文同此，不再逐一注明）：

\begin{enumerate}
\def\labelenumi{\arabic{enumi}.}
\tightlist
\item
  \textbf{找非线性项}：左端第二项是\(2y(t)\cdot\dfrac{\mathrm{d}y(t)}{\mathrm{d}t}\)，即\(y\)与其
  一阶导数的\textbf{乘积}。有这一项，系统就不可能线性。
\item
  \textbf{齐次性检验（给具体反例）}：设\(e(t)\to y(t)\)。若线性成立，则\(2e(t)\)必须给出\(2y(t)\)。
  把\(2y\)代入方程左端：
\end{enumerate}

\[ \frac{\mathrm{d}^{2}(2y)}{\mathrm{d}t^{2}}+2(2y)\frac{\mathrm{d}(2y)}{\mathrm{d}t}-3(2y)
=2y''+8yy'-6y \]

而右端\(2e(t)=2\left(y''+2yy'-3y\right)=2y''+4yy'-6y\)。两者相差\(4yy'\)，只有
\(yy'\equiv 0\)
时才相等。故''\(2e\to 2y\)``不成立，\textbf{齐次性不满足}。 3.
\textbf{叠加性同样不满足}：设\(e_{1}\to y_{1}\)、\(e_{2}\to y_{2}\)、
\(e_{1}+e_{2}\to y_{\Sigma}\)，
把\(y_{1}+y_{2}\)代入左端会多出交叉项\(2\left(y_{1}y_{2}'+y_{2}y_{1}'\right)\)，故
\(y_{\Sigma}\ne y_{1}+y_{2}\)，系统\textbf{不具备线性}。 4.
\textbf{时不变检验}：方程各项系数都是常数（\(1\)、\(2\)、\(-3\)），不显含\(t\)。设\(e(t)\to y(t)\)，
令\(e_{1}(t)=e(t-t_{0})\)，则\(y_{1}(t)=y(t-t_{0})\)代入后与原方程逐项相同------求导与常数乘法
都不引入\(t\)的显式依赖。故\textbf{具备时不变性}。

结论：\textbf{非线性、时不变}。

两句话记住这两件事为什么互相独立：出现乘积项，说明两路激励的响应会彼此干扰，这是''非线性''；
系数不随\(t\)变，说明今天做与明天做结果一样，这是''时不变''。一个说''叠加不成立''，一个说''平移
之后规律不变''。

判定范式（课件同款四例，务必记住结论）：

\begin{longtable}[]{@{}lccl@{}}
\toprule\noalign{}
方程 & 线性？ & 时不变？ & 判别依据 \\
\midrule\noalign{}
\endhead
\bottomrule\noalign{}
\endlastfoot
\(3y'(t)+5y(t)=e'(t)+e(t)\) & 是 & 是 & 一次组合、常系数 \\
\(y(t)=e^{2}(t)\) & 否 & 是 & 出现\(e\)的平方 \\
\(y(t)=t\,e(t)\) & 是 & 否 & 系数显含\(t\) \\
\(y(t)=e(t)+C,\ C\ne 0\) & 否（增量线性） & 是 & 有非零常数项 \\
\end{longtable}

\subsection{易错点}\label{ux6613ux9519ux70b9}

\begin{enumerate}
\def\labelenumi{\arabic{enumi}.}
\tightlist
\item
  \textbf{把''两个周期信号之和必为周期信号''当定理} →
  先算周期比，比值是无理数时和信号不是周 期信号 →
  反例\(f(t)=\sin t+\sin(\sqrt{2}\,t)\)。
\item
  \textbf{求周期时对频率取最小公倍数} → 必须对\textbf{周期}
  \(T_{i}=2\pi/\omega_{i}\)取最小公倍数 →
  频率与周期互为倒数，对频率取公倍数会得到完全不同的数。
\item
  \textbf{求完\(T\)不验算} → 用\(T/T_{i}\)必须全是正整数回代一次 →
  三个以上分量叠加时最容易漏 掉某个分量。
\item
  \textbf{波形变换时只对括号整体操作} →
  每一步都必须针对变量\(t\)操作；先尺度后平移时平移量是
  给定值，先平移后尺度时平移量要除以\(a\)→ 顺序搞错会让波形整体错位。
\item
  \textbf{把\(\delta(t)\)当普通函数求值，写\(\delta(0)=1\)} →
  \(\delta(t)\)只在积分号下才有意义， 用取样特性化为\(f(t_{0})\) →
  \(\delta\)是广义函数，强度是面积不是幅值。
\item
  \textbf{取样积分不看积分限} →
  先判\(\alpha<t_{0}<\beta\)是否成立，不成立直接答\(0\)→ 这是本
  章最容易被扣分的一步。
\item
  \textbf{判线性时漏掉常数项、乘积项与大信号情形} →
  有非零常数项时写''增量线性''，并说明''零
  状态响应满足线性''；出现\(y\cdot y'\)这类乘积项时一律直接判非线性 →
  课程把\(y=e+C\)与乘法
  器作为特例单列，而乘积项是完全的非线性，不能按特例放过。
\item
  \textbf{把''因果''与''连续''混为一谈} →
  因果看时间范围（\(t<0\)是否恒零），连续看自变量是否连续 取值 →
  两条分类规则彼此独立。
\end{enumerate}

\section{连续时间系统的时域分析}\label{ux8fdeux7eedux65f6ux95f4ux7cfbux7edfux7684ux65f6ux57dfux5206ux6790}

\subsection{考点清单}\label{ux8003ux70b9ux6e05ux5355-1}

\begin{longtable}[]{@{}
  >{\raggedright\arraybackslash}p{(\linewidth - 4\tabcolsep) * \real{0.2500}}
  >{\raggedright\arraybackslash}p{(\linewidth - 4\tabcolsep) * \real{0.4167}}
  >{\centering\arraybackslash}p{(\linewidth - 4\tabcolsep) * \real{0.3333}}@{}}
\toprule\noalign{}
\begin{minipage}[b]{\linewidth}\raggedright
考点
\end{minipage} & \begin{minipage}[b]{\linewidth}\raggedright
常见考法
\end{minipage} & \begin{minipage}[b]{\linewidth}\centering
重要度
\end{minipage} \\
\midrule\noalign{}
\endhead
\bottomrule\noalign{}
\endlastfoot
微分方程的算子表示 & 写出\(H(p)\)、\(D(p)\)、\(N(p)\)；算子运算律 &
★★ \\
零输入响应与自然频率 & 给方程与初始条件求\(y_{zi}(t)\)、列自然频率 &
★★★ \\
冲激响应\(h(t)\) & 部分分式法求\(h(t)\)，含\(m\ge n\)与重根 & ★★★ \\
阶跃响应\(g(t)\) & 由\(h(t)\)积分求\(g(t)\)，或反向求\(h\) & ★★ \\
卷积定义与图解法 & 由两波形用图解法求卷积并画波形 & ★★★ \\
卷积的性质 & 用\(\delta\)、\(\delta'\)、\(u\)、微积分、时移性质化简 &
★★★ \\
卷积积分限的确定 & 分段写结果，判对每段上下限 & ★★★ \\
零状态响应与全响应 & 近代时域法三步走完 & ★★★ \\
响应的分解 & 零输入/零状态、自然/受迫、瞬态/稳态互判 & ★★★ \\
子系统组合的冲激响应 & 串联求卷积、并联求和 & ★★ \\
\end{longtable}

\subsection{核心概念与公式}\label{ux6838ux5fc3ux6982ux5ff5ux4e0eux516cux5f0f-1}

\subsubsection{微分方程与算子表示}\label{ux5faeux5206ux65b9ux7a0bux4e0eux7b97ux5b50ux8868ux793a}

线性非时变系统（单输入单输出）的一般方程（\(a_{n}=1\)的归一化形式）：

\[ y^{(n)}(t)+a_{n-1}y^{(n-1)}(t)+\cdots+a_{1}y'(t)+a_{0}y(t)
= b_{m}e^{(m)}(t)+b_{m-1}e^{(m-1)}(t)+\cdots+b_{1}e'(t)+b_{0}e(t) \]

标准初始条件为\(y(0),y'(0),\cdots,y^{(n-1)}(0)\)。本质上它描述的是输入变化与输出变化之间的
动态关系。

微分算子：

\[ p=\frac{\mathrm{d}}{\mathrm{d}t},\qquad p^{n}=\frac{\mathrm{d}^{n}}{\mathrm{d}t^{n}},\qquad
\frac{1}{p}=\int_{-\infty}^{t}(\ \cdot\ )\,\mathrm{d}\tau \]

于是方程可写为\(D(p)y(t)=N(p)e(t)\)，进而

\[ H(p)=\frac{N(p)}{D(p)}
=\frac{b_{m}p^{m}+b_{m-1}p^{m-1}+\cdots+b_{1}p+b_{0}}
{p^{n}+a_{n-1}p^{n-1}+\cdots+a_{1}p+a_{0}},\qquad y(t)=H(p)e(t) \]

\(H(p)\)称为转移算子。\(D(p)=0\)的根\(\lambda_{1},\cdots,\lambda_{n}\)就是系统的\textbf{自然频率}（
特征根）。

算子运算律（只在下列三条内可放心用）：

\begin{longtable}[]{@{}ll@{}}
\toprule\noalign{}
法则 & 成立条件 \\
\midrule\noalign{}
\endhead
\bottomrule\noalign{}
\endlastfoot
\(mp+np=(m+n)p\) & \(m,n\)任意整数 \\
\(p^{m}p^{n}=p^{m+n}\) & \(m,n\)同为任意正整数（或同为负整数） \\
\(p\cdot\dfrac{1}{p}=1\) & 恒成立 \\
\end{longtable}

而\(\dfrac{1}{p}\cdot p\ne 1\)：因为

\[ \frac{1}{p}\,p\,x(t)=\int_{-\infty}^{t}\frac{\mathrm{d}x(\lambda)}{\mathrm{d}\lambda}\mathrm{d}\lambda=x(t)-x(-\infty) \]

\textbf{推论}：等式两边的公共算子因子\textbf{不能消去}，即由\(px(t)=py(t)\)不能推出\(x(t)=y(t)\)，只能
推出二者相差一个常数。这是算子法区别于代数运算的关键，也是''为什么不能用约分把\(H(p)\)化简
掉一个极点''的原因。

\subsubsection{零输入响应与自然频率}\label{ux96f6ux8f93ux5165ux54cdux5e94ux4e0eux81eaux7136ux9891ux7387}

零输入响应是齐次方程的解：

\[ D(p)y_{zi}(t)=\left(p^{n}+a_{n-1}p^{n-1}+\cdots+a_{1}p+a_{0}\right)y_{zi}(t)=0 \]

三步法：定特征根 → 写形式解 → 代初始条件定系数。

\begin{enumerate}
\def\labelenumi{\arabic{enumi}.}
\tightlist
\item
  令\(D(\lambda)=0\)，把\(p\)当作代数量，解出\(n\)个根\(\lambda_{1},\cdots,\lambda_{n}\)。
\item
  形式解：

  \begin{itemize}
  \tightlist
  \item
    无重根：\(\displaystyle y_{zi}(t)=\sum_{i=1}^{n}C_{i}\mathrm{e}^{\lambda_{i}t}\)；
  \item
    \(\lambda_{1}\)为\(k\)重根：
    \(\displaystyle y_{zi}(t)=\sum_{i=1}^{k}C_{i}t^{\,i-1}\mathrm{e}^{\lambda_{1}t}+\sum_{i=k+1}^{n}C_{i}\mathrm{e}^{\lambda_{i}t}\)。
  \end{itemize}
\item
  代入初始条件得线性方程组：
\end{enumerate}

\[ y_{zi}^{(j)}(0)=\sum_{i=1}^{n}C_{i}\lambda_{i}^{\,j},\qquad j=0,1,\cdots,n-1 \]

若初始条件给的是\(y(0^{-})\)与\(y(0^{+})\)不一致（激励在\(t=0\)有跳变、出现冲激项），要用\textbf{冲
激匹配法}先从\(0^{-}\)跳到\(0^{+}\)，再用\(y_{zi}\)的初始条件。本课程默认不区分二者，按
\(0^{-}\)处理。

\subsubsection{奇异函数与信号的时域分解}\label{ux5947ux5f02ux51fdux6570ux4e0eux4fe1ux53f7ux7684ux65f6ux57dfux5206ux89e3}

阶跃、冲激的定义与六条性质见第 1
章的``奇异函数：阶跃与冲激''一节，此处只补两条本章要用的结论。

任意（有始）信号表示为冲激函数的积分：

\[ f(t)=\int_{0}^{\infty}f(\tau)\delta(t-\tau)\,\mathrm{d}\tau
=\int_{0}^{t}f(\tau)\delta(t-\tau)\,\mathrm{d}\tau \]

表示为阶跃函数的积分（杜阿美型分解）：

\[ f(t)=f(0)u(t)+\int_{0^{+}}^{t}f'(\tau)u(t-\tau)\,\mathrm{d}\tau \]

选冲激而不用阶跃的原因：冲激分解只需对激励\textbf{取样}、不必求导，而阶跃分解需要\(f'(\tau)\)，
要求激励连续可导，条件更强、计算更繁。

\subsubsection{冲激响应与阶跃响应}\label{ux51b2ux6fc0ux54cdux5e94ux4e0eux9636ux8dc3ux54cdux5e94}

\begin{itemize}
\tightlist
\item
  阶跃响应\(g(t)\)：系统对\(u(t)\)的零状态响应；
\item
  冲激响应\(h(t)\)：系统对\(\delta(t)\)的零状态响应。
\end{itemize}

两者的关系（线性系统，知一即可）：

\[ g(t)=\int_{0^{-}}^{t}h(\tau)\,\mathrm{d}\tau,\qquad h(t)=\frac{\mathrm{d}}{\mathrm{d}t}g(t) \]

求\(h(t)\)采用 Heaviside
部分分式分解法，核心是把复杂系统拆成若干简单系统之和。\textbf{一阶结论
必须先背}：

\[ h'(t)-\lambda h(t)=k\delta(t)\ \Longrightarrow\ h(t)=k\mathrm{e}^{\lambda t}u(t) \]

简记为：\(\dfrac{k}{p-\lambda}\delta(t)=k\mathrm{e}^{\lambda t}u(t)\)。

一般系统（\(m<n\)，特征根无重根）：

\[ H(p)=\sum_{i=1}^{n}\frac{k_{i}}{p-\lambda_{i}}
\ \Longrightarrow\ h(t)=\sum_{i=1}^{n}k_{i}\mathrm{e}^{\lambda_{i}t}u(t) \]

三种情况必须分清：

\begin{longtable}[]{@{}
  >{\raggedright\arraybackslash}p{(\linewidth - 4\tabcolsep) * \real{0.2917}}
  >{\raggedright\arraybackslash}p{(\linewidth - 4\tabcolsep) * \real{0.3958}}
  >{\raggedright\arraybackslash}p{(\linewidth - 4\tabcolsep) * \real{0.3125}}@{}}
\toprule\noalign{}
\begin{minipage}[b]{\linewidth}\raggedright
分子分母阶数
\end{minipage} & \begin{minipage}[b]{\linewidth}\raggedright
\(H(p)\)的分部分式
\end{minipage} & \begin{minipage}[b]{\linewidth}\raggedright
\(h(t)\)的构成
\end{minipage} \\
\midrule\noalign{}
\endhead
\bottomrule\noalign{}
\endlastfoot
\(m<n\) & \(\sum_{i}\dfrac{k_{i}}{p-\lambda_{i}}\) &
只有\(\mathrm{e}^{\lambda_{i}t}u(t)\)项 \\
\(m=n\) & \(C_{0}+\sum_{i}\dfrac{k_{i}}{p-\lambda_{i}}\) &
多一个\(C_{0}\delta(t)\) \\
\(m>n\) & \(\sum C_{j}p^{\,j}+\sum_{i}\dfrac{k_{i}}{p-\lambda_{i}}\) &
多出\(\delta^{(j)}(t)\)项 \\
\end{longtable}

有重根时用这条：

\[ \frac{k}{(p-\lambda)^{n}}\delta(t)=\frac{k\,t^{\,n-1}}{(n-1)!}\mathrm{e}^{\lambda t}u(t),\qquad n\ge 2 \]

\subsubsection{卷积的定义、图解与性质}\label{ux5377ux79efux7684ux5b9aux4e49ux56feux89e3ux4e0eux6027ux8d28}

叠加积分（卷积）：

\[ y_{zs}(t)=e(t)*h(t)=\int_{-\infty}^{\infty}e(\tau)h(t-\tau)\,\mathrm{d}\tau \]

推导链（三步都要能复述）：\(e(t)=\int e(\tau)\delta(t-\tau)\mathrm{d}\tau\)；系统时不变
⇒ \(\delta(t-\tau)\to h(t-\tau)\)；系统齐次 + 叠加
⇒\(\int e(\tau)h(t-\tau)\mathrm{d}\tau\)。所以卷积不是人为定义的技巧，而是''线性
+ 时不变'' 共同推出的结果。

因果情形（\(e\)与\(h\)都在\(t<0\)为零）：

\[ y_{zs}(t)=\int_{0}^{t}e(\tau)h(t-\tau)\,\mathrm{d}\tau \]

图解法四步：\textbf{换元 → 反褶 → 平移 → 相乘叠加（积分）}。

卷积性质表（前三条与乘法相同，后四条不同）：

\begin{longtable}[]{@{}
  >{\raggedright\arraybackslash}p{(\linewidth - 4\tabcolsep) * \real{0.2500}}
  >{\raggedright\arraybackslash}p{(\linewidth - 4\tabcolsep) * \real{0.3333}}
  >{\raggedright\arraybackslash}p{(\linewidth - 4\tabcolsep) * \real{0.4167}}@{}}
\toprule\noalign{}
\begin{minipage}[b]{\linewidth}\raggedright
性质
\end{minipage} & \begin{minipage}[b]{\linewidth}\raggedright
表达式
\end{minipage} & \begin{minipage}[b]{\linewidth}\raggedright
成立条件
\end{minipage} \\
\midrule\noalign{}
\endhead
\bottomrule\noalign{}
\endlastfoot
交换律 & \(u*v=v*u\) & 恒成立 \\
分配律 & \(u*(v+w)=u*v+u*w\) & 恒成立 \\
结合律 & \(u*(v*w)=(u*v)*w\) & 恒成立 \\
微分 & \(\dfrac{\mathrm{d}}{\mathrm{d}t}(u*v)=u'*v=u*v'\) &
\(u,v\)可微 \\
积分 &
\(\displaystyle\int_{-\infty}^{t}(u*v)\mathrm{d}\tau=\left(\int_{-\infty}^{t}u\right)*v=u*\left(\int_{-\infty}^{t}v\right)\)
& 恒成立 \\
多重微积分 & \(u^{(m)}*v^{(n)}=(u*v)^{(m+n)}\) & \(m,n\)为整数 \\
时移 &
\(u(t)*v(t)=y(t)\Rightarrow u(t-t_{1})*v(t-t_{2})=y(t-t_{1}-t_{2})\) &
恒成立 \\
\end{longtable}

三个特殊函数（几乎每张卷子都用）：

\begin{longtable}[]{@{}
  >{\raggedright\arraybackslash}p{(\linewidth - 4\tabcolsep) * \real{0.2727}}
  >{\raggedright\arraybackslash}p{(\linewidth - 4\tabcolsep) * \real{0.2727}}
  >{\raggedright\arraybackslash}p{(\linewidth - 4\tabcolsep) * \real{0.4545}}@{}}
\toprule\noalign{}
\begin{minipage}[b]{\linewidth}\raggedright
卷积
\end{minipage} & \begin{minipage}[b]{\linewidth}\raggedright
结果
\end{minipage} & \begin{minipage}[b]{\linewidth}\raggedright
记忆方式
\end{minipage} \\
\midrule\noalign{}
\endhead
\bottomrule\noalign{}
\endlastfoot
\(f(t)*\delta(t-t_{0})\) & \(f(t-t_{0})\) & \(\delta\)是卷积的单位元 \\
\(f(t)*\delta^{(n)}(t-t_{0})\) & \(f^{(n)}(t-t_{0})\) &
与\(\delta\)的导数卷积 = 求导再平移 \\
\(f(t)*u(t)\) & \(\displaystyle\int_{-\infty}^{t}f(\tau)\mathrm{d}\tau\)
& 与阶跃卷积 = 积分 \\
\end{longtable}

\subsubsection{卷积积分限的确定（可操作判据）}\label{ux5377ux79efux79efux5206ux9650ux7684ux786eux5b9aux53efux64cdux4f5cux5224ux636e}

这是图解法唯一真正的难点。设\(f_{1}\)只在\([a,b]\)上非零、\(f_{2}\)只在\([c,d]\)上非零，求
\(f_{1}(t)*f_{2}(t)\)。

\textbf{第一步：把两个''活动区间''写成\(\tau\)的函数。}

\begin{itemize}
\tightlist
\item
  \(f_{1}(\tau)\)的非零区间：\(\tau\in[a,b]\)，固定不动；
\item
  \(f_{2}(t-\tau)\)的非零条件：\(c\le t-\tau\le d\)，即\(\tau\in[t-d,\ t-c]\)，随\(t\)整体右移。
\end{itemize}

\textbf{第二步：写出重叠区间与存在的条件。}

\[ \tau\in\left[\max(a,\ t-d),\ \min(b,\ t-c)\right] \]

上下限必须满足\(\max(a,t-d)<\min(b,t-c)\)，即该区间非空；否则重叠为零，卷积值为\(0\)。

\textbf{第三步：只用两个区间的四个端点就够划分。}

四个分界点（按大小排序后即为分段点）：

\[ t=a+c,\qquad t=a+d,\qquad t=b+c,\qquad t=b+d \]

含义非常直观：\(t=a+c\)右端刚接触，\(t=a+d\)时\(f_{2}\)的左端扫到\(a\)，\(t=b+c\)时\(f_{2}\)的右端扫
到\(b\)，\(t=b+d\)完全离开。

\begin{longtable}[]{@{}
  >{\raggedright\arraybackslash}p{(\linewidth - 8\tabcolsep) * \real{0.1923}}
  >{\raggedright\arraybackslash}p{(\linewidth - 8\tabcolsep) * \real{0.2308}}
  >{\raggedright\arraybackslash}p{(\linewidth - 8\tabcolsep) * \real{0.1923}}
  >{\raggedright\arraybackslash}p{(\linewidth - 8\tabcolsep) * \real{0.1923}}
  >{\raggedright\arraybackslash}p{(\linewidth - 8\tabcolsep) * \real{0.1923}}@{}}
\toprule\noalign{}
\begin{minipage}[b]{\linewidth}\raggedright
位置关系
\end{minipage} & \begin{minipage}[b]{\linewidth}\raggedright
\(t\)的范围
\end{minipage} & \begin{minipage}[b]{\linewidth}\raggedright
积分下限
\end{minipage} & \begin{minipage}[b]{\linewidth}\raggedright
积分上限
\end{minipage} & \begin{minipage}[b]{\linewidth}\raggedright
几何图像
\end{minipage} \\
\midrule\noalign{}
\endhead
\bottomrule\noalign{}
\endlastfoot
未进入 & \(t<a+c\) & --- & --- & 无重叠，结果\(0\) \\
头部进入 & \(a+c\le t\)，且\(t<a+d\)、\(t<b+c\) & \(a\) & \(t-c\) &
只吃进\(f_{2}\)的右端 \\
完全进入 & \(a+d\le t<b+c\) & \(a\) & \(b\) &
全重叠，结果是斜率的线性函数 \\
尾部离开 & \(t\ge b+c\)，且\(t<b+d\) & \(t-d\) & \(b\) &
只剩\(f_{2}\)的左端 \\
已离开 & \(t\ge b+d\) & --- & --- & 无重叠，结果\(0\) \\
\end{longtable}

（``完全进入''一段只有在\(a+d<b+c\)时才存在；若\(a+d\ge b+c\)，两段部分重叠直接衔接。）

\textbf{三条校验规则}（算完必查，能挡住绝大多数错）：

\begin{enumerate}
\def\labelenumi{\arabic{enumi}.}
\tightlist
\item
  结果的分段点数 = 四个分界点中互不相同的个数；
\item
  有限时宽信号的卷积，时宽\textbf{相加}：结果支撑区长度为\((b-a)+(d-c)\)，起点为\(a+c\)；
\item
  两个因果信号卷积，结果必因果，分段点全在\(t\ge 0\)一侧。
\end{enumerate}

\subsubsection{全响应的四种分解}\label{ux5168ux54cdux5e94ux7684ux56dbux79cdux5206ux89e3}

\[ y(t)=y_{zi}(t)+y_{zs}(t) \]

\begin{longtable}[]{@{}
  >{\raggedright\arraybackslash}p{(\linewidth - 4\tabcolsep) * \real{0.2857}}
  >{\raggedright\arraybackslash}p{(\linewidth - 4\tabcolsep) * \real{0.3571}}
  >{\raggedright\arraybackslash}p{(\linewidth - 4\tabcolsep) * \real{0.3571}}@{}}
\toprule\noalign{}
\begin{minipage}[b]{\linewidth}\raggedright
分解名
\end{minipage} & \begin{minipage}[b]{\linewidth}\raggedright
两个分量
\end{minipage} & \begin{minipage}[b]{\linewidth}\raggedright
界定依据
\end{minipage} \\
\midrule\noalign{}
\endhead
\bottomrule\noalign{}
\endlastfoot
零输入 / 零状态 & \(y_{zi}\)只由初始状态决定；\(y_{zs}\)只由激励决定 &
谁引起 \\
自然 / 受迫 &
\(y_{n}\)是与自然频率\(\lambda_{i}\)对应的项；\(y_{f}\)是与激励极点对应的项
& 极点归属 \\
瞬态 / 稳态 & \(y_{tr}\)随\(t\)增长趋于零；\(y_{ss}\)不趋于零 &
长时间行为 \\
通解 / 特解 & 经典法：齐次解 + 特解 & 求解方式 \\
\end{longtable}

关键从属关系（考试常以判断形式出现）：\(y_{zi}\)是自然响应的一部分，但自然响应还包括
\(y_{zs}\)中与\(\lambda_{i}\)对应的那部分；\(y_{zs}\)除含受迫响应外，还含一部分自然响应。换句
话说，\(y_{zi}\)含于\(y_{n}\)，\(y_{f}\)含于\(y_{zs}\)。所以''\(y_{zi}=\)自然响应''\,``\(y_{zs}=\)受迫响
应''都是错的。

对稳定系统：自然响应必属于瞬态响应；受迫响应可能是瞬态也可能是稳态，取决于激励。

指数激励\(e(t)=\mathrm{e}^{st}u(t)\)（\(s\ne\lambda_{i}\)）的响应分解：

\[ y(t)=\sum_{i=1}^{n}\left(c_{i}-\frac{k_{i}}{s-\lambda_{i}}\right)\mathrm{e}^{\lambda_{i}t}u(t)
+\sum_{i=1}^{n}\frac{k_{i}}{s-\lambda_{i}}\mathrm{e}^{st}u(t) \]

前一项是自然响应（系数被激励修正过），后一项是受迫响应。若\(s\)恰等于某个\(\lambda_{i}\)，出
现\(\mathrm{e}^{\lambda t}u(t)*\mathrm{e}^{\lambda t}u(t)=t\,\mathrm{e}^{\lambda t}u(t)\)，
这就是\textbf{谐振}。

\subsection{典型题型与解法}\label{ux5178ux578bux9898ux578bux4e0eux89e3ux6cd5-1}

\subsubsection{题型 2-1
求自然频率与零输入响应}\label{ux9898ux578b-2-1-ux6c42ux81eaux7136ux9891ux7387ux4e0eux96f6ux8f93ux5165ux54cdux5e94}

\textbf{题型特征}：给微分方程（或转移算子）与未加激励时的初始条件，要求自然频率\(\lambda_{i}\)
与零输入响应\(y_{zi}(t)\)。

\textbf{解题步骤}： 1. 写出\(D(p)\)，取方程左端并入算子形式。 2.
令\(D(\lambda)=0\)解特征根，\textbf{重根要标明重数}。 3.
按根的重数写形式解：无重根为\(\sum C_{i}\mathrm{e}^{\lambda_{i}t}\)；\(k\)重根必须补齐
\(t\,\mathrm{e}^{\lambda t},\,t^{2}\mathrm{e}^{\lambda t},\cdots\)共\(k\)项；共轭复根
\(-\alpha\pm\mathrm{j}\beta\)用\(\mathrm{e}^{-\alpha t}\left(C_{1}\cos\beta t
   +C_{2}\sin\beta t\right)\)。 4.
用初始条件列方程组，解\(C_{i}\)；\textbf{待定系数的个数必须等于初始条件的个数}。
5.
答句写全：``自然频率为\(\lambda=\cdots\)，零输入响应\(y_{zi}(t)=\cdots\)（\(t\ge 0\)）''。

\textbf{例题}（【附加作业1
第7题】）：已知系统的微分方程与未加激励下的初始条件为

\[ \frac{\mathrm{d}^{3}}{\mathrm{d}t^{3}}y(t)+2\frac{\mathrm{d}^{2}}{\mathrm{d}t^{2}}y(t)
+\frac{\mathrm{d}}{\mathrm{d}t}y(t)=3\frac{\mathrm{d}}{\mathrm{d}t}e(t)+e(t),
\qquad y(0)=y'(0)=0,\quad y''(0)=1 \]

求系统的零输入响应，并指出自然频率。

解答：

\begin{enumerate}
\def\labelenumi{\arabic{enumi}.}
\tightlist
\item
  写转移算子（分子来自激励端、分母来自响应端）：
\end{enumerate}

\[ H(p)=\frac{3p+1}{p^{3}+2p^{2}+p}=\frac{3p+1}{p\,(p+1)^{2}} \]

\begin{enumerate}
\def\labelenumi{\arabic{enumi}.}
\setcounter{enumi}{1}
\tightlist
\item
  令\(D(\lambda)=\lambda(\lambda+1)^{2}=0\)，得
\end{enumerate}

\[ \lambda_{1}=0,\qquad \lambda_{2}=\lambda_{3}=-1 \]

其中\(\lambda=-1\)是\textbf{二重根}，故自然频率为\(0\)与\(-1\)（\(-1\)计二重）。

\begin{enumerate}
\def\labelenumi{\arabic{enumi}.}
\setcounter{enumi}{2}
\tightlist
\item
  二重根必须补 \(t\mathrm{e}^{-t}\) 这一项，形式解写成
\end{enumerate}

\[ y_{zi}(t)=C_{1}+C_{2}\mathrm{e}^{-t}+C_{3}t\,\mathrm{e}^{-t} \]

这里有三个待定系数，正好对应三个初始条件。若偷懒只写\(C_{1}+C_{2}\mathrm{e}^{-t}\)，两个未
知数无法同时满足三个初始条件------这是本题型最容易崩的地方。

\begin{enumerate}
\def\labelenumi{\arabic{enumi}.}
\setcounter{enumi}{3}
\tightlist
\item
  逐阶求导：
\end{enumerate}

\[ \begin{aligned}
y_{zi}'(t)&=-C_{2}\mathrm{e}^{-t}+C_{3}\mathrm{e}^{-t}-C_{3}t\,\mathrm{e}^{-t}
=\left(C_{3}-C_{2}\right)\mathrm{e}^{-t}-C_{3}t\,\mathrm{e}^{-t}\\
y_{zi}''(t)&=-\left(C_{3}-C_{2}\right)\mathrm{e}^{-t}-C_{3}\mathrm{e}^{-t}
+C_{3}t\,\mathrm{e}^{-t}=\left(C_{2}-2C_{3}\right)\mathrm{e}^{-t}
+C_{3}t\,\mathrm{e}^{-t}
\end{aligned} \]

\begin{enumerate}
\def\labelenumi{\arabic{enumi}.}
\setcounter{enumi}{4}
\tightlist
\item
  令\(t=0\)代初始条件，得线性方程组
\end{enumerate}

\[ \begin{aligned}
y_{zi}(0)&=C_{1}+C_{2}=0\\
y_{zi}'(0)&=C_{3}-C_{2}=0\\
y_{zi}''(0)&=C_{2}-2C_{3}=1
\end{aligned} \]

由第二式\(C_{3}=C_{2}\)，代入第三式得\(C_{2}-2C_{2}=-C_{2}=1\)，故\(C_{2}=-1\)、\(C_{3}=-1\)，
再由第一式\(C_{1}=1\)。

\begin{enumerate}
\def\labelenumi{\arabic{enumi}.}
\setcounter{enumi}{5}
\tightlist
\item
  结论：
\end{enumerate}

\[ y_{zi}(t)=1-\left(1+t\right)\mathrm{e}^{-t},\qquad t\ge 0;\qquad
\lambda_{1}=0,\ \lambda_{2}=\lambda_{3}=-1 \]

验算：\(y_{zi}(0)=1-1=0\)；\(y_{zi}'(0)=\left(1+t\right)\mathrm{e}^{-t}-\mathrm{e}^{-t}
   =t\,\mathrm{e}^{-t}\big|_{0}=0\)；\(y_{zi}''(0)=\mathrm{e}^{-t}-t\mathrm{e}^{-t}\big|_{
   0}=1\)， 三个初始条件全部满足。

\textbf{同类题}： - 【附加作业1
第6题】三组转移算子与同一组初始条件\(y(0)=1\)、\(y'(0)=2\)，正好覆盖三种根型：
(1) \(H(p)=\dfrac{p+3}{p^{2}+3p+2}\)，\(\lambda=-1,-2\)，
\(y_{zi}=\left(4\mathrm{e}^{-t}-3\mathrm{e}^{-2t}\right)u(t)\)； (2)
\(H(p)=\dfrac{p+3}{p^{2}+2p+2}\)，共轭复根\(\lambda=-1\pm\mathrm{j}\)，
\(y_{zi}=\mathrm{e}^{-t}\left(\cos t+3\sin t\right)u(t)\)； (3)
\(H(p)=\dfrac{p+3}{p^{2}+2p+1}\)，二重根\(\lambda=-1\)，
\(y_{zi}=\mathrm{e}^{-t}\left(1+3t\right)u(t)\)。
注意各小问只需求分母的根，分子\(p+3\)不参与自然频率的计算。 -
【附加作业2
第5题】的零输入部分同一套流程：\(D(\lambda)=\lambda^{2}+3\lambda+2
  =(\lambda+1)(\lambda+2)\)，由\(y_{zi}(0^{-})=1\)、\(y_{zi}'(0^{-})=2\)解得\(C_{1}=4\)、
\(C_{2}=-3\)，即\(y_{zi}(t)=\left(4\mathrm{e}^{-t}-3\mathrm{e}^{-2t}\right)u(t)\)。

\subsubsection{题型 2-2
用部分分式法求冲激响应}\label{ux9898ux578b-2-2-ux7528ux90e8ux5206ux5206ux5f0fux6cd5ux6c42ux51b2ux6fc0ux54cdux5e94}

\textbf{题型特征}：给\(H(p)=N(p)/D(p)\)或微分方程，要求\(h(t)\)；常与''再由\(h\)反推某个待定系数''组
合成一题。

\textbf{解题步骤}： 1.
把\(H(p)\)的分母写成\((p-\lambda_{1})(p-\lambda_{2})\cdots\)。 2.
比较\(m\)与\(n\)：\(m<n\)直接分部分式；\(m=n\)先分出常数\(C_{0}\)；\(m>n\)先做多项式除法。
3. 用待定系数或留数法求各\(k_{i}\)。 4.
逐项用\(\dfrac{k}{p-\lambda}\delta(t)=k\mathrm{e}^{\lambda t}u(t)\)写出\(h(t)\)，重根用
\(\dfrac{t^{n-1}}{(n-1)!}\)公式。 5. \textbf{每次都要补上\(u(t)\)}。

\textbf{例题}（【附加作业2 第4题】）：某 LTI 因果连续系统的微分方程为
\(y''(t)+5y'(t)+6y(t)=e(t)\)，激励\(e(t)=\mathrm{e}^{-t}u(t)\)，全响应为
\(c\,\mathrm{e}^{-t}u(t)\)。求 (1) 初始状态\(y(0),y'(0)\)；(2)
系数\(c\)。要求用近代时域法求解。

解答：

第一步，写转移算子：

\[ H(p)=\frac{1}{p^{2}+5p+6}=\frac{1}{(p+2)(p+3)} \]

部分分式：

\[ H(p)=\frac{1}{p+2}-\frac{1}{p+3} \]

第二步，冲激响应（\(m=0<n=2\)，无\(\delta\)项）：

\[ h(t)=\left(\mathrm{e}^{-2t}-\mathrm{e}^{-3t}\right)u(t) \]

第三步，零状态响应。用基本卷积对（\(a\ne b\)，\(t\ge 0\)）：

\[ \mathrm{e}^{at}u(t)*\mathrm{e}^{bt}u(t)=\frac{\mathrm{e}^{at}-\mathrm{e}^{bt}}{a-b}\,u(t) \]

于是

\[ \begin{aligned}
\left(\mathrm{e}^{-2t}u\right)*\left(\mathrm{e}^{-t}u\right)
&=\frac{\mathrm{e}^{-2t}-\mathrm{e}^{-t}}{-2-(-1)}=\mathrm{e}^{-t}-\mathrm{e}^{-2t}\\
\left(\mathrm{e}^{-3t}u\right)*\left(\mathrm{e}^{-t}u\right)
&=\frac{\mathrm{e}^{-3t}-\mathrm{e}^{-t}}{-3-(-1)}=\frac{1}{2}\left(\mathrm{e}^{-t}-\mathrm{e}^{-3t}\right)
\end{aligned} \]

\[ y_{zs}(t)=h(t)*e(t)=\frac{1}{2}\mathrm{e}^{-t}-\mathrm{e}^{-2t}+\frac{1}{2}\mathrm{e}^{-3t},\qquad t\ge 0 \]

第四步，全响应并比较系数。\(D(p)\)的根为\(\lambda=-2,-3\)，故

\[ y_{zi}(t)=C_{1}\mathrm{e}^{-2t}+C_{2}\mathrm{e}^{-3t} \]

\[ y(t)=y_{zi}+y_{zs}=\frac{1}{2}\mathrm{e}^{-t}+\left(C_{1}-1\right)\mathrm{e}^{-2t}
+\left(C_{2}+\frac{1}{2}\right)\mathrm{e}^{-3t} \]

题目给出全响应\(y(t)=c\,\mathrm{e}^{-t}u(t)\)，只含\(\mathrm{e}^{-t}\)一项，故两端对应项系数
相等给出

\[ C_{1}-1=0\ \Rightarrow\ C_{1}=1,\qquad C_{2}+\frac{1}{2}=0\ \Rightarrow\ C_{2}=-\frac{1}{2},\qquad c=\frac{1}{2} \]

第五步，回求初始状态。零输入响应满足\(y_{zi}(0)=y(0)\)、\(y_{zi}'(0)=y'(0)\)：

\[ y(0)=C_{1}+C_{2}=1-\frac{1}{2}=\frac{1}{2},\qquad
y'(0)=-2C_{1}-3C_{2}=-2+\frac{3}{2}=-\frac{1}{2} \]

结论：\(y(0)=\dfrac{1}{2}\)，\(y'(0)=-\dfrac{1}{2}\)，\(c=\dfrac{1}{2}\)。

验算：把\(y(t)=\frac{1}{2}\mathrm{e}^{-t}\)代回原方程，左端
\(=\frac{1}{2}\mathrm{e}^{-t}-\frac{5}{2}\mathrm{e}^{-t}+3\mathrm{e}^{-t}=\mathrm{e}^{-t}\)，
等于右端，结果自洽。

\textbf{同类题}： - 【附加作业2 第5题】的第一步求\(h(t)\)是同一套路：
\(H(p)=\dfrac{3}{p^{2}+3p+2}=\dfrac{3}{p+1}-\dfrac{3}{p+2}\)，故
\(h(t)=\left(3\mathrm{e}^{-t}-3\mathrm{e}^{-2t}\right)u(t)\)。

\subsubsection{题型 2-3
卷积法求零状态响应与全响应}\label{ux9898ux578b-2-3-ux5377ux79efux6cd5ux6c42ux96f6ux72b6ux6001ux54cdux5e94ux4e0eux5168ux54cdux5e94}

\textbf{题型特征}：给转移算子/微分方程、激励表达式与初始状态，要求零状态响应、零输入响应与全
响应，且要求''采用近代时域法''。

\textbf{解题步骤}： 1. 求\(H(p)\)，读出\(D(p)\)的根（自然频率）。 2.
求\(y_{zi}\)：形式解代初始条件。 3. 求\(h(t)=H(p)\delta(t)\)。 4.
求\(y_{zs}=e*h\)：用基本卷积对，注意\(u(t)\)与收敛性。 5.
全响应\(y=y_{zi}+y_{zs}\)，并按需指出自然/受迫、瞬态/稳态分量。

\textbf{例题}（【附加作业2
第5题】）：系统算子表示为\(p^{2}y(t)+3py(t)+2y(t)=3e(t)\)，
\(e(t)=\mathrm{e}^{-4t}u(t)\)，初始状态\(y_{zi}(0^{-})=1\)、
\(\dfrac{\mathrm{d}}{\mathrm{d}t}y_{zi}(0^{-})=2\)。求\(y_{zs}(t)\)、\(y_{zi}(t)\)与全响应
\(y(t)\)。

解答：

1）零输入响应。特征方程\(\lambda^{2}+3\lambda+2=0\)，得\(\lambda_{1}=-1\)、\(\lambda_{2}=-2\)
（自然频率）。形式解

\[ y_{zi}(t)=C_{1}\mathrm{e}^{-t}+C_{2}\mathrm{e}^{-2t} \]

代初始条件：

\[ \begin{cases}C_{1}+C_{2}=1\\ -C_{1}-2C_{2}=2\end{cases} \]

两式相加得\(-C_{2}=3\)，即\(C_{2}=-3\)；回代得\(C_{1}=4\)。故

\[ y_{zi}(t)=4\mathrm{e}^{-t}-3\mathrm{e}^{-2t},\qquad t\ge 0 \]

2）冲激响应。

\[ H(p)=\frac{3}{p^{2}+3p+2}=\frac{3}{(p+1)(p+2)}=\frac{3}{p+1}-\frac{3}{p+2} \]

\[ h(t)=\left(3\mathrm{e}^{-t}-3\mathrm{e}^{-2t}\right)u(t) \]

3）零状态响应。\(e(t)=\mathrm{e}^{-4t}u(t)\)，用基本卷积对：

\[ \begin{aligned}
\left(\mathrm{e}^{-t}u\right)*\left(\mathrm{e}^{-4t}u\right)
&=\frac{\mathrm{e}^{-t}-\mathrm{e}^{-4t}}{-1-(-4)}=\frac{1}{3}\left(\mathrm{e}^{-t}-\mathrm{e}^{-4t}\right)\\
\left(\mathrm{e}^{-2t}u\right)*\left(\mathrm{e}^{-4t}u\right)
&=\frac{\mathrm{e}^{-2t}-\mathrm{e}^{-4t}}{-2-(-4)}=\frac{1}{2}\left(\mathrm{e}^{-2t}-\mathrm{e}^{-4t}\right)
\end{aligned} \]

\[ \begin{aligned}
y_{zs}(t)&=3\cdot\frac{1}{3}\left(\mathrm{e}^{-t}-\mathrm{e}^{-4t}\right)
-3\cdot\frac{1}{2}\left(\mathrm{e}^{-2t}-\mathrm{e}^{-4t}\right)\\
&=\mathrm{e}^{-t}-1.5\,\mathrm{e}^{-2t}+0.5\,\mathrm{e}^{-4t},\qquad t\ge 0
\end{aligned} \]

4）全响应：

\[ \begin{aligned}
y(t)&=y_{zi}+y_{zs}\\
&=\left(4+1\right)\mathrm{e}^{-t}+\left(-3-1.5\right)\mathrm{e}^{-2t}+0.5\,\mathrm{e}^{-4t}\\
&=5\mathrm{e}^{-t}-4.5\,\mathrm{e}^{-2t}+0.5\,\mathrm{e}^{-4t},\qquad t\ge 0
\end{aligned} \]

5）分量归属：\(5\mathrm{e}^{-t}-4.5\mathrm{e}^{-2t}\)属于自然响应（与\(\lambda=-1,-2\)对应，
其中\(-1.5\mathrm{e}^{-2t}\)与\(1\mathrm{e}^{-t}\)在\(y_{zs}\)内部就已出现）；
\(0.5\mathrm{e}^{-4t}\)属于受迫响应（与激励极点\(s=-4\)对应）。三项都随\(t\to\infty\)趋于零，
故全部是瞬态响应，稳态响应为零。

\textbf{同类题}： - 课件 2.9
节练习（方程\(y''+5y'+6y=e\)，\(e=\mathrm{e}^{-t}u(t)\)，\(y(0^{-})=3.5\)、
\(y'(0^{-})=-8.5\)）走的是完全相同五步，只是数值不同：\(h\)与\(y_{zs}\)的表达式与题型
2-2 的
例题一致，再叠加\(y_{zi}=C_{1}\mathrm{e}^{-2t}+C_{2}\mathrm{e}^{-3t}\)即可。

\subsubsection{题型 2-4
图解法计算卷积并确定积分限}\label{ux9898ux578b-2-4-ux56feux89e3ux6cd5ux8ba1ux7b97ux5377ux79efux5e76ux786eux5b9aux79efux5206ux9650}

\textbf{题型特征}：给两个波形（多为矩形、斜坡等分段直线），要求用图解法求卷积并画出结果波形。

\textbf{解题步骤}： 1.
写出两个信号的支撑区（非零区间）\([a,b]\)与\([c,d]\)，以及每段上的表达式。
2.
按四步做图：换元\(f_{1}(t)\to f_{1}(\tau)\)、反褶\(f_{2}(t)\to f_{2}(-\tau)\)、平移
\(f_{2}(-\tau)\to f_{2}(t-\tau)\)（\(t>0\)右移、\(t<0\)左移）、相乘后对\(\tau\)积分。
3.
用``卷积积分限的确定''一节的四个分界点\(a+c,\ a+d,\ b+c,\ b+d\)排序，划分\(t\)的区间。
4.
逐段定上下限并计算。若其中一个是常数（如\(f_{1}\equiv A\)），直接把积分化成''常数
\(\times\) 两块图形的重合面积''，即\(y(t)=A\cdot S\)，比硬算积分快得多。
5. 画结果波形，标出各段端点坐标，并校验连续性、峰值与时宽。

\textbf{例题}（【附加作业2 第3题】）：用图解法求图 3
所示两个信号的卷积。其中

\[ f_{1}(t)=A\quad(0<t<1),\qquad
f_{2}(t)=\frac{B}{2}t\quad(0<t<2) \]

两者在区间外均为\(0\)。

解答：支撑区\([a,b]=[0,1]\)、\([c,d]=[0,2]\)。反褶平移后\(f_{2}(t-\tau)\)的非零区间为
\(\tau\in[t-2,\ t]\)。四个分界点

\[ a+c=0,\qquad b+c=1,\qquad a+d=2,\qquad b+d=3 \]

排序为\(0,1,2,3\)，正好把\(t\)轴分成五段，与``卷积积分限的确定''一节的判据一致。

\(f_{1}\equiv A\)是常数，故每段都用面积法：把两图形的重合面积记为\(S\)，则
\(y(t)=f_{1}(t)*f_{2}(t)=A\cdot S\)，只需按几何图形算\(S\)，不必写积分。

\begin{enumerate}
\def\labelenumi{\arabic{enumi}.}
\tightlist
\item
  \(t\le 0\)：两支撑区不重叠，\(S=0\)，\(y(t)=0\)。
\item
  \(0<t\le 1\)：重叠段为\(0<\tau<t\)，图形是以\(t\)为底、以\(\dfrac{B}{2}t\)为高的直角三角形，
\end{enumerate}

\[ S=\frac{1}{2}\cdot t\cdot\frac{B}{2}t=\frac{B}{4}t^{2},
\qquad y(t)=\frac{AB}{4}t^{2} \]

\begin{enumerate}
\def\labelenumi{\arabic{enumi}.}
\setcounter{enumi}{2}
\tightlist
\item
  \(1<t\le 2\)：重叠段为\(0<\tau<1\)，图形是梯形，两底高分别为\(\dfrac{B}{2}t\)与
  \(\dfrac{B}{2}(t-1)\)，
\end{enumerate}

\[ y(t)=A\cdot\frac{1}{2}\left[\frac{B}{2}t+\frac{B}{2}(t-1)\right]
=\frac{AB}{4}\left(2t-1\right) \]

\begin{enumerate}
\def\labelenumi{\arabic{enumi}.}
\setcounter{enumi}{3}
\tightlist
\item
  \(2<t\le 3\)：重叠段为\(t-2<\tau<1\)，宽度\(3-t\)，两端高度分别为\(B\)与\(\dfrac{B}{2}(t-1)\)，
\end{enumerate}

\[ y(t)=A\cdot\frac{1}{2}\left(3-t\right)\left[B+\frac{B}{2}(t-1)\right]
=\frac{AB}{4}\left(3+2t-t^{2}\right) \]

\begin{enumerate}
\def\labelenumi{\arabic{enumi}.}
\setcounter{enumi}{4}
\tightlist
\item
  \(t>3\)：两支撑区不再重叠，\(y(t)=0\)。
\end{enumerate}

汇总：

\[ f_{1}(t)*f_{2}(t)=\begin{cases}
0, & t\le 0\\
\dfrac{AB}{4}t^{2}, & 0<t\le 1\\
\dfrac{AB}{4}\left(2t-1\right), & 1<t\le 2\\
\dfrac{AB}{4}\left(3+2t-t^{2}\right), & 2<t\le 3\\
0, & t>3
\end{cases} \]

三重校验：

\begin{enumerate}
\def\labelenumi{\arabic{enumi}.}
\tightlist
\item
  \textbf{端点连续}：\(t=1\)处两式都给\(\dfrac{AB}{4}\)；\(t=2\)处两式都给\(\dfrac{3AB}{4}\)。
\item
  \textbf{时宽相加}：\(f_{1}\)时宽\(1\)、\(f_{2}\)时宽\(2\)，结果支撑区应为\((0,3)\)、长\(3\)，与上面的分段
  一致。
\item
  \textbf{面积守恒}：\(\displaystyle\int y\,\mathrm{d}t=\left(\int f_{1}\right)
  \left(\int f_{2}\right)=A\cdot\frac{B}{2}\cdot 2=A B\)；把四段积分相加得
  \(\dfrac{AB}{12}+\dfrac{AB}{2}+\dfrac{5AB}{12}=AB\)，通过。峰值\(\dfrac{3AB}{4}\)在\(t=2\)处。
\end{enumerate}

画图：波形从\((0,0)\)按抛物线升到\(\left(1,\dfrac{AB}{4}\right)\)，再线性升到
\(\left(2,\dfrac{3AB}{4}\right)\)，之后按抛物线降到\((3,0)\)。抛物段与直线段的连接处切线不连
续，这不是错误------被积函数在\(t\)跨过分段点时换了一支。

\textbf{同类题}： - 变式（课件 2.8 节同法例）：\(f(\tau)=2\) 在
\((4,5)\) 上、\(h(\tau)=1\) 在 \((1,3)\) 上，四界点 为
\(a+c=5\)、\(b+c=6\)、\(a+d=7\)、\(b+d=8\)，五段结果为
\(0\)、\(2t-10\)、\(2\)、\(16-2t\)、\(0\)。
与本题对照可见：面积法的前提是\textbf{其中一个信号在重合段上是常数}（本题
\(f_{1}\equiv A\)、 该变式中 \(h\equiv 1\) 都满足）；若两个信号都随
\(\tau\) 变化，就只能逐段写定积分，而上下限 仍然由四界点给出。

\subsubsection{题型 2-5
用卷积性质化简，求子系统组合的冲激响应}\label{ux9898ux578b-2-5-ux7528ux5377ux79efux6027ux8d28ux5316ux7b80ux6c42ux5b50ux7cfbux7edfux7ec4ux5408ux7684ux51b2ux6fc0ux54cdux5e94}

\textbf{题型特征}：两种常见小计算------(i)
激励或冲激响应含\(\delta(t-t_{0})\)，要求用性质直接写
出卷积而不作积分；(ii)
给若干子系统的\(h_{i}(t)\)与联接方式，求组合系统的\(h(t)\)。

\textbf{解题步骤}： 1.
把含冲激的因子识别出来，用\(f(t)*\delta(t-t_{0})=f(t-t_{0})\)和
\(f(t)*\delta^{(n)}(t-t_{0})=f^{(n)}(t-t_{0})\)逐项拆。 2.
展开后用线性性合并同类项。 3.
看框图的联接方式再套公式：\textbf{串联（级联）用卷积}
\(h=h_{1}*h_{2}\)，\textbf{并联用相加}
\(h=h_{1}+h_{2}\)；\textbf{直通支路相当于\(h=\delta(t)\)}，它与其他支路仍是并联关系。
4.
冲激之间的卷积用\(\delta(t-t_{1})*\delta(t-t_{2})=\delta(t-t_{1}-t_{2})\)。
5.
与激励级联时，可先把整个系统的\(h(t)\)求出来，再与\(e(t)\)卷积；也可写成
\(y(t)=e(t)*h(t)\)并利用\(e(t)*\delta(t)=e(t)\)把直通支路直接落成\(e(t)\)。

\textbf{例题}（【附加作业2
第2题】）：已知信号\(x(t)\)与\(h(t)\)的波形，其中
\(h(t)=2\delta(t+2)+\delta(t-2)\)，要求画出卷积\(y(t)=x(t)*h(t)\)的波形。

解答：先分离冲激，再用时移性质，最后用线性性叠加：

\[ \begin{aligned}
y(t)&=x(t)*\left[2\delta(t+2)+\delta(t-2)\right]\\
&=2\,x(t)*\delta(t+2)+x(t)*\delta(t-2)\\
&=2x(t+2)+x(t-2)
\end{aligned} \]

三步理由：分配律拆开；\(f(t)*\delta(t-t_{0})=f(t-t_{0})\)逐项搬移；标量系数按齐次性提到外面。

画图要点：把\(x(t)\)的波形\textbf{左移 2} 得\(x(t+2)\)，幅度乘
2；再把\(x(t)\) \textbf{右移 2} 得
\(x(t-2)\)，幅度不变；两者逐点相加即为\(y(t)\)。关键点坐标用
\((t_{k},x_{k})\to(t_{k}-2,\ 2x_{k})\)与\((t_{k},x_{k})\to(t_{k}+2,\ x_{k})\)分别搬两次。注
意\(t+2\)是左移、\(t-2\)是右移，符号与方向相反是最常见的错。

\textbf{同类题}： - 【附加作业2
第1题】\(h_{1}(t)=\delta(t-1)\)、\(h_{2}(t)=\delta(t)-\delta(t-3)\)，图
1 的结构
是\textbf{直通支路、\(h_{1}\)支路、``\(h_{1}\)与\(h_{2}\)级联''支路，三条支路并联}（输出相加）。
并联用相加、级联用卷积，先写总冲激响应：

\[ h(t)=\delta(t)+h_{1}(t)+h_{1}(t)*h_{2}(t) \]

再逐项算。级联那一条：

\[ h_{1}(t)*h_{2}(t)=\delta(t-1)*\left[\delta(t)-\delta(t-3)\right]
=\delta(t-1)-\delta(t-4) \]

代回并合并同延迟项（\(\delta(t-1)\)出现两次）：

\[ h(t)=\delta(t)+\delta(t-1)+\delta(t-1)-\delta(t-4)
=\delta(t)+2\delta(t-1)-\delta(t-4) \]

有了\(h(t)\)，系统对任意激励的响应就是

\[ y(t)=e(t)*h(t)=e(t)*\left[\delta(t)+h_{1}(t)+h_{1}(t)*h_{2}(t)\right] \]

三个要点：直通支路贡献\(e(t)\)本身；两条含\(h_{1}\)的支路把\(e(t)\)延迟\(1\)后叠加，故
\(\delta(t-1)\)的系数是\(2\)；只有\(h_{1}\)与\(h_{2}\)级联的那条才会出现\(4\)的延迟（\(1+3\)）。

\subsubsection{题型 2-6
由两组全响应拆出零输入与零状态分量}\label{ux9898ux578b-2-6-ux7531ux4e24ux7ec4ux5168ux54cdux5e94ux62c6ux51faux96f6ux8f93ux5165ux4e0eux96f6ux72b6ux6001ux5206ux91cf}

\textbf{题型特征}：同一系统、同一初始状态、两种不同激励，给出两个全响应，要求第三种激励下的全
响应。题面常强调''初始状态不变''。

\textbf{解题步骤}： 1.
写分解式：\(y_{1}=y_{zi}+y_{zs}^{(1)}\)，\(y_{2}=y_{zi}+y_{zs}^{(2)}\)。
2.
用齐次性把\(y_{zs}^{(2)}\)表示为\(y_{zs}^{(1)}\)的倍数（激励成倍数，零状态响应同倍数）。
3. 两式相减消去\(y_{zi}\)，解出\(y_{zs}^{(1)}\)。 4. 回代求\(y_{zi}\)。
5. 对第三种激励，\(y_{3}=y_{zi}+k\,y_{zs}^{(1)}\)。

\textbf{例题}（【附加作业1 第5题】）：一 LTI
系统具有非零初始状态。激励为\(e(t)\)时全响应

\[ y_{1}(t)=\mathrm{e}^{-t}+2\cos(\pi t),\qquad t>0 \]

初始状态不变、激励为\(2e(t)\)时全响应

\[ y_{2}(t)=3\cos(\pi t),\qquad t>0 \]

求同样初始状态下激励为\(3e(t)\)时的全响应\(y_{3}(t)\)。

解答：设激励\(e(t)\)单独作用（零状态）产生的响应为\(y_{zs}(t)\)，初始状态单独作用产生的响应
为\(y_{zi}(t)\)。由线性性，激励\(2e(t)\)时的零状态响应为\(2y_{zs}(t)\)，而\(y_{zi}(t)\)与激励无
关、保持不变。于是

\[ \begin{cases}y_{zi}(t)+y_{zs}(t)=\mathrm{e}^{-t}+2\cos(\pi t)\\[2pt]
y_{zi}(t)+2y_{zs}(t)=3\cos(\pi t)\end{cases} \]

两式相减，得零状态响应：

\[ y_{zs}(t)=3\cos(\pi t)-\mathrm{e}^{-t}-2\cos(\pi t)=\cos(\pi t)-\mathrm{e}^{-t} \]

回代第一式，得零输入响应：

\[ y_{zi}(t)=\mathrm{e}^{-t}+2\cos(\pi t)-\left[\cos(\pi t)-\mathrm{e}^{-t}\right]
=2\mathrm{e}^{-t}+\cos(\pi t) \]

激励为\(3e(t)\)时：

\[ \begin{aligned}
y_{3}(t)&=y_{zi}(t)+3y_{zs}(t)
=\left[2\mathrm{e}^{-t}+\cos(\pi t)\right]+3\left[\cos(\pi t)-\mathrm{e}^{-t}\right]\\
&=4\cos(\pi t)-\mathrm{e}^{-t},\qquad t>0
\end{aligned} \]

答案：\(y_{3}(t)=4\cos(\pi t)-\mathrm{e}^{-t}\)，\(t>0\)。

自检：三条响应在\(t\to 0^{+}\)处的值分别为\(y_{1}(0)=3\)、\(y_{2}(0)=3\)、\(y_{3}(0)=3\)。激励成
倍变化时零状态分量在\(t=0^{+}\)处为零（\(\cos 0-\mathrm{e}^{0}=0\)），故全响应的初值只由
\(y_{zi}(0)=3\)决定，与算得的\(y_{zi}(0)=2+1=3\)吻合，结果自洽。

\subsection{易错点}\label{ux6613ux9519ux70b9-1}

\begin{enumerate}
\def\labelenumi{\arabic{enumi}.}
\tightlist
\item
  \textbf{把\(y_{zi}\)当作自然响应、把\(y_{zs}\)当作受迫响应} → 自然响应
  = \(y_{zi}\) + \(y_{zs}\)中 与\(\lambda_{i}\)对应的项；受迫响应 =
  \(y_{zs}\)中与激励极点对应的项 → 两个分法的划分依
  据不同，不能一一对应。
\item
  \textbf{求\(h(t)\)漏掉\(u(t)\)} → 每项指数后面都补\(u(t)\) →
  零状态响应默认\(t<0\)无输出，漏 \(u(t)\)等于承认\(t<0\)也有响应。
\item
  \textbf{\(m\ge n\)时漏掉\(\delta(t)\)或\(\delta^{(m-n)}(t)\)项} →
  动笔前先比较\(m\)与\(n\)，\(m=n\)
  补\(C_{0}\delta(t)\)，\(m>n\)先做多项式除法 →
  长除法漏项是本题型失分的主因。
\item
  \textbf{卷积一律写\(\int_{-\infty}^{\infty}\)不分段} →
  先定支撑区，再用四个分界点 \(a+c,a+d,b+c,b+d\)划区间，逐段写上下限 →
  不重叠的区间上被积函数恒为零，硬套全域积分得 不到分段结果。
\item
  \textbf{上下限顺序颠倒} → 保证\(\max(a,t-d)<\min(b,t-c)\)再积 →
  上限小于下限时物理上就是'' 无重叠''，应写\(0\)而不是负值。
\item
  \textbf{算子运算中约掉公因子} →
  保留\(D(p)\)的全部因子，用\(p\cdot\frac{1}{p}=1\)但
  \(\frac{1}{p}\cdot p\ne 1\)的差别提醒自己 →
  微分与积分次序不可交换，约分等价于假设 \(x(-\infty)=0\)。
\item
  \textbf{用\(y(0^{+})\)当零输入响应的初始条件} →
  先判\(t=0\)是否有冲激项或跳变，用冲激匹配法从
  \(0^{-}\)跳到\(0^{+}\)，再解\(C_{i}\)→
  方程右端含\(\delta\)时\(y(0^{+})\ne y(0^{-})\)。
\item
  \textbf{对两个因子各求一次导再相乘} →
  一次求导只能加在\textbf{一个}因子上：
  \(\dfrac{\mathrm{d}}{\mathrm{d}t}\left(f_{1}*f_{2}\right)=f_{1}'*f_{2}=f_{1}*f_{2}'\)；
  写成\(f_{1}'*f_{2}'\)就多求了一次导，它等于\(\dfrac{\mathrm{d}^{2}}{\mathrm{d}t^{2}}
  \left(f_{1}*f_{2}\right)\)而不是一阶导数 → 依据多重微积分性质
  \(f_{1}^{(m)}*f_{2}^{(n)}=\left(f_{1}*f_{2}\right)^{(m+n)}\)，两个因子上的阶数会\textbf{累加}，
  不是取其一。
\item
  \textbf{谐振情形当成普通情形} →
  检查激励的\(s\)是否等于某个\(\lambda_{i}\)，相等时卷积出现
  \(t\mathrm{e}^{\lambda t}\) →
  \(\mathrm{e}^{\lambda t}u*\mathrm{e}^{\lambda t}u=t\mathrm{e}^{\lambda t}u\)，不能约成零。
\item
  \textbf{图解法反褶后忘记''平移方向与\(t\)的符号相反''} →
  \(h(-\tau)\)右移\(t\)（\(t>0\)）得到
  \(h(t-\tau)=h\left[-(\tau-t)\right]\)，是把\(h(-\tau)\)的图形整体\textbf{向右}移动\(t\)
  → 写成 \(h(-(\tau+t))\)就变成左移了。
\end{enumerate}

\section{连续信号的正交分解}\label{ux8fdeux7eedux4fe1ux53f7ux7684ux6b63ux4ea4ux5206ux89e3}

\subsection{考点清单}\label{ux8003ux70b9ux6e05ux5355-2}

\begin{longtable}[]{@{}
  >{\raggedright\arraybackslash}p{(\linewidth - 4\tabcolsep) * \real{0.2500}}
  >{\raggedright\arraybackslash}p{(\linewidth - 4\tabcolsep) * \real{0.4167}}
  >{\centering\arraybackslash}p{(\linewidth - 4\tabcolsep) * \real{0.3333}}@{}}
\toprule\noalign{}
\begin{minipage}[b]{\linewidth}\raggedright
考点
\end{minipage} & \begin{minipage}[b]{\linewidth}\raggedright
常见考法
\end{minipage} & \begin{minipage}[b]{\linewidth}\centering
重要度
\end{minipage} \\
\midrule\noalign{}
\endhead
\bottomrule\noalign{}
\endlastfoot
正交函数与正交函数集 & 判断两函数在给定区间上是否正交 & ★★ \\
正交分解的最佳系数 & 要求方均误差最小，求近似式的系数 & ★★★ \\
完备正交函数集 & 归一化、正交化、完备性三条含义判断 & ★★ \\
三角形式傅里叶级数 & 给波形或表达式求 \(a_0\)、\(a_n\)、\(b_n\) & ★★★ \\
指数形式傅里叶级数 & 与三角形式互化，写出 \(c_n\) 与负频率项 & ★★★ \\
周期信号的频谱 & 画单边/双边频谱，读谱得幅度与相位 & ★★★ \\
频谱的离散性、谐波性、收敛性 & 判断频谱图形是否合法、估计包络衰减 &
★★ \\
周期信号的功率与帕塞瓦尔定理 & 求平均功率、有效值、方均误差 & ★★★ \\
周期信号的频带宽度 & 按三种定义求 \(B\)，判断脉宽与带宽关系 & ★★ \\
函数的奇偶性与级数成分 & 判断级数只含正弦项、余弦项还是二者皆有 & ★★★ \\
\end{longtable}

\subsection{核心概念与公式}\label{ux6838ux5fc3ux6982ux5ff5ux4e0eux516cux5f0f-2}

\subsubsection{正交函数与正交函数集}\label{ux6b63ux4ea4ux51fdux6570ux4e0eux6b63ux4ea4ux51fdux6570ux96c6}

\textbf{正交}：在区间 \((t_1,t_2)\) 上若

\[
\int_{t_1}^{t_2} f_1(t)f_2(t)\,\mathrm{d}t = 0
\]

则称 \(f_1(t)\) 与 \(f_2(t)\) 在该区间上正交。复函数的内积取共轭：

\[
\int_{t_1}^{t_2} f_1(t)f_2^{*}(t)\,\mathrm{d}t = 0
\]

什么时候用：题目给两个函数和一个区间，问''是否正交''，
只需算定积分看是否为零；零则正交，非零则不正交。

\textbf{正交函数集}：若 \(n\) 个函数 \(\{f_1,f_2,\cdots,f_n\}\) 在区间
\((t_1,t_2)\) 上两两正交，则称为正交函数集；若还满足每个函数自身内积为
\(1\)， 则称为归一化正交函数集。

\textbf{易混点}：``正交''只在\textbf{指定的区间}上成立。同一对函数换了区间，
结论可以完全相反，所以答题时区间必须与题目一致。

\subsubsection{信号的正交分解与最佳系数}\label{ux4fe1ux53f7ux7684ux6b63ux4ea4ux5206ux89e3ux4e0eux6700ux4f73ux7cfbux6570}

用单个函数近似：设 \(f(t)\approx c_1f_1(t)\)，方均误差定义为

\[
\overline{\varepsilon^{2}(t)}
=\frac{1}{t_2-t_1}\int_{t_1}^{t_2}\bigl[f(t)-c_1f_1(t)\bigr]^{2}\,\mathrm{d}t
\]

对 \(c_1\) 求极小（令导数
\(\frac{\mathrm{d}}{\mathrm{d}c_1}\overline{\varepsilon^{2}}=0\)），
得方均误差最小的\textbf{最佳系数}

\[
c_1=\frac{\int_{t_1}^{t_2} f(t)f_1(t)\,\mathrm{d}t}
{\int_{t_1}^{t_2} f_1^{2}(t)\,\mathrm{d}t}
\]

成立条件：\(f_1(t)\) 在区间上平方可积且不恒为零。

用 \(n\) 个正交函数近似：\(f(t)\approx\sum_{i=1}^{n}c_if_i(t)\)，各系数
\textbf{互不影响}，每个都按上式单独求：

\[
c_i=\frac{\int_{t_1}^{t_2} f(t)f_i(t)\,\mathrm{d}t}
{\int_{t_1}^{t_2} f_i^{2}(t)\,\mathrm{d}t},\qquad i=1,2,\cdots,n
\]

此时最小方均误差为

\[
\overline{\varepsilon^{2}(t)}
=\frac{1}{t_2-t_1}\left[\int_{t_1}^{t_2} f^{2}(t)\,\mathrm{d}t
-\sum_{i=1}^{n}c_i^{2}\int_{t_1}^{t_2} f_i^{2}(t)\,\mathrm{d}t\right]\ge 0
\]

最后这个''\(\ge 0\)``就是\textbf{贝塞尔不等式}，它给了一个上界：
任何有限项正交展开都不可能比原信号能量还大。

什么时候用：凡是出现''方均误差最小''\,``最佳近似''``正弦信号近似表示''
字样，一律用''内积除以自内积''，不要去解微分方程。

\subsubsection{完备正交函数集}\label{ux5b8cux5907ux6b63ux4ea4ux51fdux6570ux96c6}

\begin{longtable}[]{@{}
  >{\raggedright\arraybackslash}p{(\linewidth - 4\tabcolsep) * \real{0.2727}}
  >{\raggedright\arraybackslash}p{(\linewidth - 4\tabcolsep) * \real{0.4545}}
  >{\raggedright\arraybackslash}p{(\linewidth - 4\tabcolsep) * \real{0.2727}}@{}}
\toprule\noalign{}
\begin{minipage}[b]{\linewidth}\raggedright
条件
\end{minipage} & \begin{minipage}[b]{\linewidth}\raggedright
数学表述
\end{minipage} & \begin{minipage}[b]{\linewidth}\raggedright
作用
\end{minipage} \\
\midrule\noalign{}
\endhead
\bottomrule\noalign{}
\endlastfoot
归一化 & \(\int_{t_1}^{t_2} f_i f_i^{*}\,\mathrm{d}t=1\) &
系数公式分母化为 \(1\) \\
正交化 & \(\int_{t_1}^{t_2} f_i f_j^{*}\,\mathrm{d}t=0,\ i\ne j\) &
各系数可独立求解 \\
完备性 & 可用其线性组合表示任意信号 & 展开无残差 \\
\end{longtable}

\textbf{完备性}的准确含义：不存在与集合中每个函数都正交的非零函数。
若集合不完备，则再多项也无法把方均误差压到零。

什么时候用：问''为什么傅里叶级数能表示任意周期信号''------因为三角函数集
与复指数函数集都是完备正交函数集。

\subsubsection{三角形式傅里叶级数}\label{ux4e09ux89d2ux5f62ux5f0fux5085ux91ccux53f6ux7ea7ux6570}

设 \(f(t)\) 周期为 \(T\)，基频 \(\omega_{0}=\frac{2\pi}{T}\)（课件记作
\(\Omega\)， 同一个量；本资料全篇统一用 \(\omega_{0}\)），则

\[
f(t)=\frac{a_0}{2}+\sum_{n=1}^{+\infty}\bigl(a_n\cos n\omega_{0} t
+b_n\sin n\omega_{0} t\bigr)
\]

系数由''内积除以自内积''直接得出（这就是 3.2 的最佳系数）：

\[
a_n=\frac{2}{t_2-t_1}\int_{t_1}^{t_2} f(t)\cos n\omega_{0} t\,\mathrm{d}t,
\qquad n=0,1,2,\cdots
\]

\[
b_n=\frac{2}{t_2-t_1}\int_{t_1}^{t_2} f(t)\sin n\omega_{0} t\,\mathrm{d}t,
\qquad n=1,2,\cdots
\]

其中 \(t_2-t_1=T\)；若取主值区间 \((-\frac{T}{2},\frac{T}{2})\)，
系数就是 \(\frac{2}{T}\int_{-T/2}^{T/2}\cdots\mathrm{d}t\)。 注意
\(a_0\) 是把 \(n=0\) 代入 \(a_n\) 的公式得到的，所以 \textbf{直流分量是
\(\frac{a_0}{2}\) 而不是 \(a_0\)}。

\textbf{成立条件（狄利克雷条件）}，三条缺一不可：

\begin{enumerate}
\def\labelenumi{\arabic{enumi}.}
\tightlist
\item
  在一周期内连续，或只有有限个第一类间断点；
\item
  在一周期内极大值与极小值的数目有限；
\item
  在一周期内绝对可积，即
  \(\int_{t_1}^{t_1+T}|f(t)|\,\mathrm{d}t<\infty\)。
\end{enumerate}

在间断点 \(t_0\) 处，级数收敛到左右极限的平均值
\(\frac{f(t_0^{+})+f(t_0^{-})}{2}\)，不等于 \(f(t_0)\) 本身。

\textbf{另一种等价写法（幅度-相角形式）}：

\[
f(t)=\frac{A_0}{2}+\sum_{n=1}^{+\infty}A_n\cos(n\omega_{0} t+\varphi_n),
\qquad A_n=\sqrt{a_n^{2}+b_n^{2}}
\]

其中 \(a_n=A_n\cos\varphi_n\)，\(b_n=-A_n\sin\varphi_n\)，
\(\varphi_n=-\arctan\frac{b_n}{a_n}\)，且 \(\varphi_n\) 的象限由
\((a_n,-b_n)\) 所在象限决定，不能只看反正切主值。

\subsubsection{指数形式傅里叶级数及三种形式换算}\label{ux6307ux6570ux5f62ux5f0fux5085ux91ccux53f6ux7ea7ux6570ux53caux4e09ux79cdux5f62ux5f0fux6362ux7b97}

\[
f(t)=\sum_{n=-\infty}^{+\infty}c_ne^{\mathrm{j}n\omega_{0} t},
\qquad
c_n=\frac{1}{T}\int_{t_1}^{t_1+T} f(t)e^{-\mathrm{j}n\omega_{0} t}\,\mathrm{d}t
\]

由于 \(e^{\mathrm{j}n\omega_{0} t}\) 与 \(e^{-\mathrm{j}n\omega_{0} t}\)
成对出现，
复指数形式一定含\textbf{负频率}。负频率没有物理意义，只是数学上的配对项，
但画双边频谱时必须画出来。

对实信号 \(f(t)\)，三种形式之间的换算是固定的：

\[
c_n=\frac{1}{2}\left(a_n-\mathrm{j}b_n\right),
\qquad
c_{-n}=\frac{1}{2}\left(a_n+\mathrm{j}b_n\right)=c_n^{*},
\qquad
c_0=\frac{a_0}{2}
\]

\begin{longtable}[]{@{}
  >{\raggedright\arraybackslash}p{(\linewidth - 6\tabcolsep) * \real{0.1463}}
  >{\raggedright\arraybackslash}p{(\linewidth - 6\tabcolsep) * \real{0.1951}}
  >{\raggedright\arraybackslash}p{(\linewidth - 6\tabcolsep) * \real{0.4146}}
  >{\raggedright\arraybackslash}p{(\linewidth - 6\tabcolsep) * \real{0.2439}}@{}}
\toprule\noalign{}
\begin{minipage}[b]{\linewidth}\raggedright
形式
\end{minipage} & \begin{minipage}[b]{\linewidth}\raggedright
直流项
\end{minipage} & \begin{minipage}[b]{\linewidth}\raggedright
第 \(n\) 次谐波项
\end{minipage} & \begin{minipage}[b]{\linewidth}\raggedright
系数范围
\end{minipage} \\
\midrule\noalign{}
\endhead
\bottomrule\noalign{}
\endlastfoot
三角（正弦余弦） & \(\frac{a_0}{2}\) &
\(a_n\cos n\omega_{0} t+b_n\sin n\omega_{0} t\) & \(n\ge 1\) \\
三角（幅度相角） & \(\frac{A_0}{2}\) &
\(A_n\cos(n\omega_{0} t+\varphi_n)\) & \(n\ge 1\) \\
复指数 & \(c_0=\frac{a_0}{2}\) & 成对出现 \(c_n\) 与 \(c_{-n}\) 两项 &
全部整数 \\
\end{longtable}

三条换算口诀：\(c_0=\frac{a_0}{2}=\frac{A_0}{2}\)；\(|c_n|=\frac{A_n}{2}\)（\(n\ge1\)），
且 \(|c_{-n}|=|c_n|\)；\(c_n\) 的辐角就是 \(+\varphi_n\)，而 \(c_{-n}\)
的辐角是 \(-\varphi_n\)。

\subsubsection{周期信号的频谱}\label{ux5468ux671fux4fe1ux53f7ux7684ux9891ux8c31}

\textbf{频谱}是把各频率分量的系数画在''频率轴''上得到的图形， 横坐标是
\(\omega\)（或 \(n\omega_{0}\)），纵坐标是幅度或相位。

\textbf{单边频谱}：只画 \(\omega\ge 0\)，纵坐标为 \(A_n\) 与
\(\varphi_n\)， \(n=0\) 处的值是 \(\frac{A_0}{2}\)（即直流分量）。

\textbf{双边频谱}：画 \(\omega\) 全部取值，纵坐标为 \(|c_n|\) 与
\(\arg c_n\)。
实信号的双边幅度谱关于纵轴\textbf{偶对称}，相位谱\textbf{奇对称}，
所以只画 \(\omega\ge 0\) 一半也能还原全部信息。

\textbf{幅度换算}：\(n\ge 1\)
时单边幅度是双边幅度的两倍，\(A_n=2|c_n|\)； \(n=0\) 时单边画
\(\frac{A_0}{2}=c_0\)，双边画 \(|c_0|\)，二者相等。

\textbf{三条特点}（判断频谱图是否合法的三条标尺）：

\begin{itemize}
\tightlist
\item
  离散性：谱线是一条条不连续的线，不是连续曲线；
\item
  谐波性：谱线只出现在基频 \(\omega_{0}\) 的整数倍处，间隔恒为
  \(\omega_{0}\)；
\item
  收敛性：\(n\to\infty\) 时 \(|c_n|\to 0\)，包络随频率衰减。
\end{itemize}

\textbf{周期矩形脉冲}（幅度 \(E\)、脉宽 \(\tau\)、周期
\(T\)）是必背模型：

\[
c_n=\frac{E\tau}{T}\,\mathrm{Sa}\!\left(\frac{n\omega_{0}\tau}{2}\right)
=\frac{E\tau}{T}\cdot\frac{\sin(n\omega_{0}\tau/2)}{n\omega_{0}\tau/2}
\]

其中抽样函数 \(\mathrm{Sa}(x)=\frac{\sin x}{x}\)（与 \(\mathrm{sinc}\)
定义相差 一个 \(\pi\) 因子，本资料统一用 \(\mathrm{Sa}\)），且
\(\mathrm{Sa}(0)=1\)。 包络的第一个零点在
\(\frac{n\omega_{0}\tau}{2}=\pi\)，即 \(f=\frac{1}{\tau}\)。

\subsubsection{周期信号的功率与帕塞瓦尔定理}\label{ux5468ux671fux4fe1ux53f7ux7684ux529fux7387ux4e0eux5e15ux585eux74e6ux5c14ux5b9aux7406}

周期信号功率有限、能量无限，所以只谈\textbf{平均功率}。设周期为 \(T\)：

\[
P=\frac{1}{T}\int_{t_1}^{t_1+T} f^{2}(t)\,\mathrm{d}t
=\frac{a_0^{2}}{4}+\sum_{n=1}^{+\infty}\frac{a_n^{2}+b_n^{2}}{2}
=\frac{A_0^{2}}{4}+\sum_{n=1}^{+\infty}\frac{A_n^{2}}{2}
=\sum_{n=-\infty}^{+\infty}|c_n|^{2}
\]

成立条件：\(f(t)\) 周期且在一个周期内平方可积（功率有限）。

\textbf{帕塞瓦尔定理}说的就是上式最后一步的等式：时域平均功率等于
各次谐波平均功率之和，功率可以''按频率分配''。由此得到有效值

\[
f_{\mathrm{rms}}=\sqrt{P}
=\sqrt{\frac{A_0^{2}}{4}+\sum_{n=1}^{+\infty}\frac{A_n^{2}}{2}}
\]

注意有效值\textbf{不能}逐项相加，只有功率能相加。

对非周期信号改谈能量，对应式是 \[
E=\int_{-\infty}^{+\infty}f^{2}(t)\,\mathrm{d}t
=\frac{1}{2\pi}\int_{-\infty}^{+\infty}|F(\mathrm{j}\omega)|^{2}\,\mathrm{d}\omega
\] （成立条件：\(f(t)\) 能量有限，也就是平方可积。这里要区分两个条件：
``平方可积''与''绝对可积''是两回事，不能互推。平方可积不一定绝对可积，
反例 \(f(t)=t^{-3/4}u(t-1)\) 满足 \(\int|f|^{2}\mathrm{d}t<\infty\) 但
\(\int|f|\mathrm{d}t=\infty\)；绝对可积也不一定平方可积，反例
\(f(t)=t^{-1/2}\) 在 \((0,1)\)
上。傅里叶\textbf{变换的存在条件}用的是绝对可积
（见``傅里叶变换的定义与存在条件''一节），而帕塞瓦尔等式只要求信号能量有限。）

\subsubsection{周期信号的频带宽度}\label{ux5468ux671fux4fe1ux53f7ux7684ux9891ux5e26ux5bbdux5ea6}

由于收敛性，可以只保留主要分量、丢掉高频尾巴，被保留的频率区间
就是\textbf{频带}。三种定义必须在题目中认准是哪一种：

\begin{enumerate}
\def\labelenumi{\arabic{enumi}.}
\tightlist
\item
  以最大幅度的 \(\frac{1}{10}\) 为限；
\item
  以幅度谱的\textbf{第一个过零点}为限（最常用）；
\item
  以包含总能量（功率）的 \(90\%\) 为限。
\end{enumerate}

按第 2 种定义，周期矩形脉冲的频带宽度为

\[
B_{\omega}=\frac{2\pi}{\tau}\ \mathrm{rad/s},
\qquad B_{f}=\frac{1}{\tau}\ \mathrm{Hz}
\]

即
\(B_{\omega}\tau=2\pi\)、\(B_{f}\tau=1\)：\textbf{脉宽越窄，带宽越宽}。
更一般的结论是''边沿变化越快（跳变越陡），频带越宽''，
所以三角波的频带比方波窄、正弦波最窄。

\subsection{典型题型与解法}\label{ux5178ux578bux9898ux578bux4e0eux89e3ux6cd5-2}

\subsubsection{题型 3-1
用方均误差最小求正交分解系数并验证正交性}\label{ux9898ux578b-3-1-ux7528ux65b9ux5747ux8befux5deeux6700ux5c0fux6c42ux6b63ux4ea4ux5206ux89e3ux7cfbux6570ux5e76ux9a8cux8bc1ux6b63ux4ea4ux6027}

\textbf{题型特征}：给一个区间和一个具体波形（方波、三角波等），
要求''用某正弦/余弦信号近似，方均误差最小''，或''证明与某函数正交''。

\textbf{解题步骤}：

\begin{enumerate}
\def\labelenumi{\arabic{enumi}.}
\tightlist
\item
  写清区间 \((t_1,t_2)\)，把近似式写成 \(f(t)\approx c_1f_1(t)\)；
\item
  套最佳系数公式：分子算 \(\int_{t_1}^{t_2} f(t)f_1(t)\,\mathrm{d}t\)，
  分母算
  \(\int_{t_1}^{t_2} f_1^{2}(t)\,\mathrm{d}t\)，分段函数按段积分；
\item
  分母常用半角公式 \(\sin^{2}\theta=\frac{1-\cos 2\theta}{2}\) 化简；
\item
  验证正交只算 \(\int_{t_1}^{t_2} f(t)g(t)\,\mathrm{d}t\)，
  分段代入后看是否为零，\textbf{不要}去算系数。
\end{enumerate}

\textbf{例题}（【附加作业3 第1题】）：已知

\[
f(t)=\begin{cases}
1, & 0<t<\pi\\
-1, & \pi<t<2\pi
\end{cases}
\]

\begin{enumerate}
\def\labelenumi{(\arabic{enumi})}
\tightlist
\item
  在同一区间上用正弦信号近似 \(f(t)\)，要求方均误差最小，
  写出该正弦信号的表达式；(2) 证明 \(f(t)\) 与同一区间上的
  \(\cos nt\)（\(n\) 为整数）正交。
\end{enumerate}

解 (1)：区间为 \((0,2\pi)\)，取 \(f_1(t)=\sin t\)，设近似式为
\(f(t)\approx c_{12}\sin t\)。由最佳系数公式

\[
c_{12}=\frac{\int_{0}^{2\pi} f(t)\sin t\,\mathrm{d}t}
{\int_{0}^{2\pi}\sin^{2}t\,\mathrm{d}t}
\]

分子分段积分：

\[
\int_{0}^{2\pi} f(t)\sin t\,\mathrm{d}t
=\int_{0}^{\pi}\sin t\,\mathrm{d}t-\int_{\pi}^{2\pi}\sin t\,\mathrm{d}t
=2-(-2)=4
\]

其中 \(-\cos t\big|_{0}^{\pi}=2\)，\(-\cos t\big|_{\pi}^{2\pi}=-2\)。
分母用半角公式：

\[
\int_{0}^{2\pi}\sin^{2}t\,\mathrm{d}t
=\int_{0}^{2\pi}\frac{1-\cos 2t}{2}\,\mathrm{d}t
=\frac{1}{2}\cdot 2\pi-\frac{1}{4}\sin 2t\Big|_{0}^{2\pi}=\pi
\]

所以 \(c_{12}=\frac{4}{\pi}\approx 1.273\)，即近似正弦信号为

\[
f(t)\approx\frac{4}{\pi}\sin t
\]

（这里 \(f_1=\sin t\) 未归一化，分母 \(\pi\ne 1\)，所以系数不是
\(\frac{4}{\pi}\) 之外 的别的值；若改写成归一化基
\(\frac{\sin t}{\sqrt{\pi}}\)，系数才变成 \(\frac{4}{\sqrt{\pi}}\)，
两者给出的近似信号完全一样。）

解 (2)：\(n\) 为整数，直接算内积：

\[
\int_{0}^{2\pi} f(t)\cos nt\,\mathrm{d}t
=\int_{0}^{\pi}\cos nt\,\mathrm{d}t-\int_{\pi}^{2\pi}\cos nt\,\mathrm{d}t
\]

当 \(n\ne 0\) 时，

\[
=\frac{1}{n}\sin nt\Big|_{0}^{\pi}-\frac{1}{n}\sin nt\Big|_{\pi}^{2\pi}
=\frac{\sin n\pi}{n}-\frac{\sin 2n\pi-\sin n\pi}{n}=0
\]

当 \(n=0\) 时 \(\cos 0=1\)，积分为 \(\pi-\pi=0\)。故对所有整数 \(n\)
内积为零， \(f(t)\) 与 \(\cos nt\) 在 \((0,2\pi)\) 上正交。

\textbf{同类题}：

\begin{itemize}
\tightlist
\item
  【附加作业3 第3题】判断级数只含正弦项还是余弦项，关键是看
  ``在一个完整周期上延拓后的奇偶性''，而不是看区间内的表达式。
\end{itemize}

\subsubsection{题型 3-2
由波形展开三角形式傅里叶级数并控制方均误差}\label{ux9898ux578b-3-2-ux7531ux6ce2ux5f62ux5c55ux5f00ux4e09ux89d2ux5f62ux5f0fux5085ux91ccux53f6ux7ea7ux6570ux5e76ux63a7ux5236ux65b9ux5747ux8befux5dee}

\textbf{题型特征}：给波形图（三角波、方波、锯齿波），要求在指定区间
展开为有限项三角级数，并使方均误差小于总能量的某个百分比。

\textbf{解题步骤}：

\begin{enumerate}
\def\labelenumi{\arabic{enumi}.}
\tightlist
\item
  写出分段表达式，读出周期 \(T\) 与基频 \(\omega_{0}=\frac{2\pi}{T}\)；
\item
  判断奇偶性：偶函数 \(b_n=0\)，奇函数 \(a_n=0\)、\(a_0=0\)，
  可省一半积分；
\item
  算 \(a_0\)（或 \(a_n\)）：分段积分，含 \(t\cos nt\) 的项用分部积分，
  注意 \(\cos n\pi=(-1)^{n}\)；
\item
  算总能量 \(E=\int_{t_1}^{t_2} f^{2}(t)\,\mathrm{d}t\)；
\item
  用
  \(\overline{\varepsilon^{2}}=\frac{1}{t_2-t_1}\bigl[E-\sum c_i^{2}k_i\bigr]\)
  逐项累加，比较 \(\frac{\overline{\varepsilon^{2}}}{E}\)
  与题目给的百分比， 取到刚好满足为止；
\item
  写出有限项级数并注明 \(n\) 的取值范围。
\end{enumerate}

\textbf{例题}（【附加作业3 第2题】）：三角信号在 \((-\pi,\pi)\) 上为
关键点 \((-\pi,0)\)、\((0,A)\)、\((\pi,0)\) 连成的折线，即

\[
f(t)=\begin{cases}
A\left(1+\dfrac{t}{\pi}\right), & -\pi\le t\le 0\\[4pt]
A\left(1-\dfrac{t}{\pi}\right), & 0\le t\le\pi
\end{cases}
\]

要求有限项三角级数与实际信号的方均误差小于 \(f(t)\) 总能量的 \(1\%\)。

解：\(f(t)\) 在 \((-\pi,\pi)\) 上是偶函数，故 \(b_n=0\)；又 \(T=2\pi\)，
\(\omega_{0}=\frac{2\pi}{T}=1\)。

\[
a_0=\frac{2}{T}\int_{0}^{T} f(t)\,\mathrm{d}t
=\frac{1}{\pi}\left[\int_{-\pi}^{0}A\left(1+\frac{t}{\pi}\right)\mathrm{d}t
+\int_{0}^{\pi}A\left(1-\frac{t}{\pi}\right)\mathrm{d}t\right]
\]

两个积分都等于 \(\frac{A\pi}{2}\)（三角形面积除以底），故

\[
a_0=\frac{1}{\pi}\left(\frac{A\pi}{2}+\frac{A\pi}{2}\right)=A
\]

于是直流分量为 \(\frac{a_0}{2}=\frac{A}{2}\)。再算余弦系数：

\[
a_n=\frac{2}{T}\int_{-\pi}^{\pi} f(t)\cos nt\,\mathrm{d}t
=\frac{1}{\pi}\cdot 2\int_{0}^{\pi}A\left(1-\frac{t}{\pi}\right)\cos nt\,\mathrm{d}t
\]

括号里第一项
\(\int_{0}^{\pi}\cos nt\,\mathrm{d}t=0\)（\(n\ge 1\)），第二项

\[
\int_{0}^{\pi}t\cos nt\,\mathrm{d}t
=\frac{t\sin nt}{n}\Big|_{0}^{\pi}+\frac{\cos nt}{n^{2}}\Big|_{0}^{\pi}
=\frac{\cos n\pi-1}{n^{2}}
\]

所以

\[
a_n=\frac{2A}{\pi}\cdot\left(-\frac{1}{\pi}\right)
\cdot\frac{\cos n\pi-1}{n^{2}}=\frac{2A}{n^{2}\pi^{2}}\bigl(1-\cos n\pi\bigr)
=\begin{cases}
\dfrac{4A}{n^{2}\pi^{2}}, & n=1,3,5,\cdots\\[6pt]
0, & n=2,4,6,\cdots
\end{cases}
\]

于是

\[
f(t)=A\left[\frac{1}{2}+\frac{4}{\pi^{2}}\left(\cos t+\frac{1}{3^{2}}\cos 3t
+\frac{1}{5^{2}}\cos 5t+\cdots\right)\right]
\]

总能量

\[
E=\int_{-\pi}^{\pi}f^{2}(t)\,\mathrm{d}t
=2\int_{0}^{\pi}A^{2}\left(1-\frac{t}{\pi}\right)^{2}\mathrm{d}t
=\frac{2A^{2}\pi}{3}
\]

只取直流一项（\(c_1=\frac{A}{2}\)、\(k_1=2\pi\)）时

\[
\overline{\varepsilon_0^{2}}
=\frac{1}{2\pi}\left(\frac{2\pi A^{2}}{3}-\frac{A^{2}}{4}\cdot 2\pi\right)
=\frac{A^{2}}{12}
\]

\[
\frac{\overline{\varepsilon_0^{2}}}{E}=\frac{A^{2}/12}{2\pi A^{2}/3}
=\frac{1}{8\pi}\approx 3.98\%>1\%
\]

不满足。再取直流加基波（\(c_2=\frac{4A}{\pi^{2}}\)、\(k_2=\pi\)）：

\[
\overline{\varepsilon_1^{2}}
=\frac{1}{2\pi}\left[\frac{2\pi A^{2}}{3}-\frac{A^{2}}{4}\cdot 2\pi
-\frac{16A^{2}}{\pi^{4}}\cdot\pi\right]
=\frac{A^{2}}{12}-\frac{8A^{2}}{\pi^{4}}
\]

\[
\frac{\overline{\varepsilon_1^{2}}}{E}
=\frac{0.08333A^{2}-0.08213A^{2}}{2.0944A^{2}}\approx 0.0575\%<1\%
\]

满足要求。故至少取两项（直流加基波）：

\[
f(t)\approx A\left(\frac{1}{2}+\frac{4}{\pi^{2}}\cos t\right)
\]

\textbf{同类题}：

\begin{itemize}
\tightlist
\item
  【附加作业3 第5题】求 \((\cos\omega_c t)^{2}\) 的直流分量：先用
  \(\cos^{2}\theta=\frac{1+\cos 2\theta}{2}\) 展开，直流分量即为常数项
  \(0.5\)。
\end{itemize}

\subsubsection{题型 3-3
判断傅里叶级数只含正弦项还是余弦项}\label{ux9898ux578b-3-3-ux5224ux65adux5085ux91ccux53f6ux7ea7ux6570ux53eaux542bux6b63ux5f26ux9879ux8fd8ux662fux4f59ux5f26ux9879}

\textbf{题型特征}：给一个函数和\textbf{指定区间}（常为
\(\left(0,\frac{T}{2}\right)\)），
问''仅有余弦项、仅有正弦项、还是二者均有''。

\textbf{解题步骤}：

\begin{enumerate}
\def\labelenumi{\arabic{enumi}.}
\tightlist
\item
  在指定区间之外做\textbf{周期延拓}，得到整周期的波形；
\item
  看延拓后的波形关于纵轴是否对称（偶 → 只有余弦和直流； 奇 →
  只有正弦）；
\item
  再看是否满足 \(f(t+\frac{T}{2})=-f(t)\)（奇谐 → 只有奇次谐波） 或
  \(f(t+\frac{T}{2})=f(t)\)（偶谐 → 只有偶次谐波与直流）；
\item
  结论一句话写清''延拓后为奇函数，故仅有正弦项''。
\end{enumerate}

\textbf{例题}：已知周期为 \(T\) 的信号在 \(\left(0,\frac{T}{2}\right)\)
内 \(f(t)=E\)，在 \(\left(\frac{T}{2},T\right)\) 内 \(f(t)=-E\)，
判断其三角级数含哪些成分。

解：取主值区间 \(\left(-\frac{T}{2},\frac{T}{2}\right)\)，波形为
\(\left(-\frac{T}{2},0\right)\) 上
\(-E\)、\(\left(0,\frac{T}{2}\right)\) 上 \(+E\)。 由于
\(f(-t)=-f(t)\)，是奇函数，故 \(a_0=0\)、\(a_n=0\)，只含正弦项。 又
\(f\left(t+\frac{T}{2}\right)=-f(t)\)，是奇谐函数，故只含奇次谐波。
其系数为 \(b_n=\frac{4E}{n\pi}\)（\(n\) 为奇数）\(0\)（\(n\)
为偶数），即

\[
f(t)=\frac{4E}{\pi}\left(\sin\omega_{0} t+\frac{1}{3}\sin 3\omega_{0} t
+\frac{1}{5}\sin 5\omega_{0} t+\cdots\right)
\]

\textbf{同类题}：

\begin{itemize}
\tightlist
\item
  【附加作业3 第3题】\(f(t)=A\cos\frac{2\pi}{T}t\) 在
  \(\left(0,\frac{T}{2}\right)\)
  上延拓后为奇函数，故\textbf{仅含正弦项}； 陷阱在于 \(f(t)\)
  本身是余弦表达式，但看的是延拓波形的奇偶性。
\end{itemize}

\subsubsection{题型 3-4
非正弦周期信号的三种级数形式换算}\label{ux9898ux578b-3-4-ux975eux6b63ux5f26ux5468ux671fux4fe1ux53f7ux7684ux4e09ux79cdux7ea7ux6570ux5f62ux5f0fux6362ux7b97}

\textbf{题型特征}：给几个不同频率的正弦/余弦项相加（常含相移
\(\frac{\pi}{2}\)）， 要求周期与基频、余弦形式三角级数、指数形式级数。

\textbf{解题步骤}：

\begin{enumerate}
\def\labelenumi{\arabic{enumi}.}
\tightlist
\item
  把每一项写成 \(\cos(\omega_i t+\theta_i)\) 的标准形式， 利用
  \(\sin\alpha=\cos\left(\alpha-\frac{\pi}{2}\right)\)；
\item
  各 \(\omega_i\) 除以最小角频率得到谐波次数，最小角频率即
  \(\omega_{0}\)；
\item
  周期取各分量周期的最小公倍数：\(T=\mathrm{lcm}(T_1,T_2)\)， 再由
  \(\omega_{0}=\frac{2\pi}{T}\) 核对；
\item
  套 \(c_n=\frac{A_n}{2}e^{\mathrm{j}\varphi_n}\)、\(c_{-n}=c_n^{*}\)
  逐项写， 漏掉负频率项是最常见失分点。
\end{enumerate}

\textbf{例题}（【附加作业3 第4题】）：已知

\[
f(t)=\sin 4\pi t+3\cos\left(0.5\pi t+\frac{\pi}{2}\right)
\]

求周期与基频、余弦形式三角级数、指数形式级数。

解 1)：两项角频率为 \(4\pi\ \mathrm{rad/s}\) 与
\(0.5\pi\ \mathrm{rad/s}\)， 对应周期
\(T_1=\frac{2\pi}{4\pi}=0.5\ \mathrm{s}\)、
\(T_2=\frac{2\pi}{0.5\pi}=4\ \mathrm{s}\)。取最小公倍数

\[
T=4\ \mathrm{s},\qquad \omega_{0}=\frac{2\pi}{T}=0.5\pi\ \mathrm{rad/s}
\]

且 \(\frac{4\pi}{0.5\pi}=8\)，故第二项是 8 次谐波，第一项是基波。

解 2)：先把两项统一成余弦形式。由
\(\sin\alpha=\cos\left(\alpha-\frac{\pi}{2}\right)\) 得

\[
\sin 4\pi t=\cos\left(4\pi t-\frac{\pi}{2}\right)
=\cos\left(8\omega_{0} t-\frac{\pi}{2}\right)
\]

第二项已是余弦形式，即
\(3\cos\left(\omega_{0} t+\frac{\pi}{2}\right)\)。故

\[
f(t)=\cos\left(8\omega_{0} t-\frac{\pi}{2}\right)
+3\cos\left(\omega_{0} t+\frac{\pi}{2}\right)
\]

对照 \(f(t)=\frac{A_0}{2}+\sum A_n\cos(n\omega_{0} t+\varphi_n)\)：
\(A_0=0\)；\(A_1=3\)、\(\varphi_1=\frac{\pi}{2}\)；\(A_8=1\)、\(\varphi_8=-\frac{\pi}{2}\)；
其余 \(A_n=0\)。即 1 次与 8 次谐波存在，2 至 7 次谐波为零。

解 3)：由 \(c_n=\frac{A_n}{2}e^{\mathrm{j}\varphi_n}\) 得

\[
c_1=\frac{3}{2}e^{\mathrm{j}\pi/2}=\mathrm{j}1.5,\qquad
c_{-1}=c_1^{*}=-\mathrm{j}1.5
\]

\[
c_8=\frac{1}{2}e^{-\mathrm{j}\pi/2}=-\mathrm{j}0.5,\qquad
c_{-8}=c_8^{*}=\mathrm{j}0.5
\]

于是指数形式级数为

\[
f(t)=\mathrm{j}1.5\,e^{\mathrm{j}\omega_{0} t}-\mathrm{j}1.5\,e^{-\mathrm{j}\omega_{0} t}
-\mathrm{j}0.5\,e^{\mathrm{j}8\omega_{0} t}+\mathrm{j}0.5\,e^{-\mathrm{j}8\omega_{0} t}
\]

校验：平均功率 \(P=\frac{A_1^{2}}{2}+\frac{A_8^{2}}{2}=4.5+0.5=5\)， 与
\(P=\sum|c_n|^{2}=1.5^{2}+1.5^{2}+0.5^{2}+0.5^{2}=5\) 一致。

\textbf{同类题}：

\begin{itemize}
\tightlist
\item
  【附加作业3 第2题】的偶对称三角波只需算 \(a_0\) 与 \(a_n\)，
  \(b_n\equiv 0\)，这就是''先判奇偶、再动笔积分''的收益。
\end{itemize}

\subsubsection{题型 3-5
求周期信号的平均功率、有效值与直流分量}\label{ux9898ux578b-3-5-ux6c42ux5468ux671fux4fe1ux53f7ux7684ux5e73ux5747ux529fux7387ux6709ux6548ux503cux4e0eux76f4ux6d41ux5206ux91cf}

\textbf{题型特征}：给周期波形或级数，要求平均功率、有效值，
或只问''波形中的直流分量''。

\textbf{解题步骤}：

\begin{enumerate}
\def\labelenumi{\arabic{enumi}.}
\tightlist
\item
  直流分量 = 一个周期内的平均值 = \(\frac{a_0}{2}=c_0\)；
\item
  平均功率用帕塞瓦尔定理：频域式
  \(P=\frac{A_0^{2}}{4}+\sum_{n\ge1}\frac{A_n^{2}}{2}\)， 时域式
  \(P=\frac{1}{T}\int_{T}f^{2}\,\mathrm{d}t\)，两法互相校验；
\item
  有效值 \(f_{\mathrm{rms}}=\sqrt{P}\)，只能对功率开方，不能逐项开方；
\item
  若只给 \(c_n\)，直接用 \(P=\sum_{n}|c_n|^{2}\)，含负频率项。
\end{enumerate}

\textbf{例题}：幅度为 \(\pm E\)、周期为 \(T\)
的对称方波，求平均功率与有效值。

解：时域法

\[
\begin{aligned}
P&=\frac{1}{T}\int_{0}^{T}f^{2}(t)\,\mathrm{d}t
=\frac{1}{T}\int_{0}^{T}E^{2}\,\mathrm{d}t=E^{2}
\end{aligned}
\]

故有效值为 \(E\)。频域法校验：由题型 3-3 的例题，
\(b_n=\frac{4E}{n\pi}\)（\(n\) 为奇数），

\[
P=\sum_{n=1,3,5,\cdots}\frac{b_n^{2}}{2}
=\frac{8E^{2}}{\pi^{2}}\sum_{n=1,3,5,\cdots}\frac{1}{n^{2}}
=\frac{8E^{2}}{\pi^{2}}\cdot\frac{\pi^{2}}{8}=E^{2}
\]

两法一致（用到了
\(\sum_{n=1,3,5,\cdots}\frac{1}{n^{2}}=\frac{\pi^{2}}{8}\)）。

\textbf{同类题}：

\begin{itemize}
\tightlist
\item
  【附加作业3 第5题】\((\cos\omega_c t)^{2}=0.5+0.5\cos 2\omega_c t\)，
  常数项就是直流分量 \(0.5\)；同时可知其平均功率为
  \(0.5^{2}+\frac{0.5^{2}}{2}=0.375\)，与直接积分 \(\frac{3}{8}\) 相同。
\end{itemize}

\subsection{易错点}\label{ux6613ux9519ux70b9-2}

\begin{enumerate}
\def\labelenumi{\arabic{enumi}.}
\tightlist
\item
  \textbf{最佳系数忘记除以自内积} → 必须用
  \(c_i=\frac{\int f f_i\,\mathrm{d}t}{\int f_i^{2}\,\mathrm{d}t}\) →
  正交分解的系数是''投影''，基函数未归一化时分母不为 \(1\)。
\item
  \textbf{把 \(a_0\) 当直流分量} → 直流分量是 \(\frac{a_0}{2}=c_0\) →
  展开式写的是 \(\frac{a_0}{2}+\sum\)，\(a_0\) 前面自带
  \(\frac{1}{2}\)。
\item
  \textbf{双边频谱把 \(|c_n|\) 当单边幅度用} → \(n\ge 1\) 时
  \(A_n=2|c_n|\) → 单边频谱把负频率那一半的能量''折叠''到了正频率上。
\item
  \textbf{复指数级数漏写负频率项} → 每一项都要成对写
  \(c_ne^{\mathrm{j}n\omega_{0} t}+c_{-n}e^{-\mathrm{j}n\omega_{0} t}\)
  → 实信号必然成对出现，漏项会丢掉一半幅度。
\item
  \textbf{把方均误差与总能量混为一谈} → 方均误差带 \(\frac{1}{t_2-t_1}\)
  的平均因子，比的是''平均功率''级别的量 → 题目说''小于总能量的
  \(1\%\)``时，两边量纲都要先摆清楚再比。
\item
  \textbf{区间不是 \(\left(-\frac{T}{2},\frac{T}{2}\right)\) 时仍用
  \(\frac{2}{T}\)} → 一般区间用 \(\frac{2}{t_2-t_1}\) →
  系数公式因子来自基函数的自内积。
\item
  \textbf{判断级数成分时看函数表达式而不看延拓波形} →
  必须在整个周期上判断奇偶性 → 半周期的表达式奇偶与整周期延拓的
  奇偶可以完全不同。
\item
  \textbf{奇谐函数与偶谐函数混淆} → 奇谐 \(f(t+\frac{T}{2})=-f(t)\)
  只含奇次谐波；偶谐只含偶次谐波与直流 → 记''负号对应奇次''。
\item
  \textbf{周期矩形脉冲 \(c_n\) 漏 \(\frac{\tau}{T}\) 或漏
  \(\frac{1}{2}\)} →
  \(c_n=\frac{E\tau}{T}\mathrm{Sa}\left(\frac{n\omega_{0}\tau}{2}\right)\)
  → 直接由
  \(\frac{1}{T}\int_{-\tau/2}^{\tau/2}E e^{-\mathrm{j}n\omega_{0} t}\mathrm{d}t\)
  现推一遍最稳。
\item
  \textbf{认为取有限项后级数处处等于原信号} → 只在连续点相等，
  间断点收敛到左右极限平均值，且跳变附近有吉布斯振荡 →
  方均误差可趋于零，但逐点误差不为零。
\end{enumerate}

\section{连续时间系统的频域分析}\label{ux8fdeux7eedux65f6ux95f4ux7cfbux7edfux7684ux9891ux57dfux5206ux6790}

\subsection{考点清单}\label{ux8003ux70b9ux6e05ux5355-3}

\begin{longtable}[]{@{}
  >{\raggedright\arraybackslash}p{(\linewidth - 4\tabcolsep) * \real{0.2500}}
  >{\raggedright\arraybackslash}p{(\linewidth - 4\tabcolsep) * \real{0.4167}}
  >{\centering\arraybackslash}p{(\linewidth - 4\tabcolsep) * \real{0.3333}}@{}}
\toprule\noalign{}
\begin{minipage}[b]{\linewidth}\raggedright
考点
\end{minipage} & \begin{minipage}[b]{\linewidth}\raggedright
常见考法
\end{minipage} & \begin{minipage}[b]{\linewidth}\centering
重要度
\end{minipage} \\
\midrule\noalign{}
\endhead
\bottomrule\noalign{}
\endlastfoot
傅里叶变换的定义与存在条件 & 写出正反变换式，说明充分条件 & ★★★ \\
常用信号的傅里叶变换对 & 门函数、指数、阶跃、冲激、抽样函数 & ★★★ \\
傅里叶变换的性质 & 时移、频移、尺度、卷积、微积分、对称 & ★★★ \\
周期信号的傅里叶变换 & 把级数与冲激串对应起来，写冲激序列的谱 & ★★★ \\
频率响应 \(H(\mathrm{j}\omega)\) & 由微分方程写
\(H(\mathrm{j}\omega)\)，画幅频相频 & ★★★ \\
零状态响应的频域求解 & \(R=HE\) 再逆变换，含部分分式展开 & ★★★ \\
无失真传输条件 & 判断给定 \(H(\mathrm{j}\omega)\) 是否无失真 & ★★ \\
理想低通滤波器 & 求 \(h(t)\)、阶跃响应、上升时间与带宽 & ★★ \\
物理可实现性（佩利-维纳准则） & 判断给定幅频特性能否用实际网络实现 &
★★ \\
抽样定理与信号恢复 & 求最低抽样率、恢复滤波器带宽、判断混叠 & ★★★ \\
\end{longtable}

\subsection{核心概念与公式}\label{ux6838ux5fc3ux6982ux5ff5ux4e0eux516cux5f0f-3}

\subsubsection{傅里叶变换的定义与存在条件}\label{ux5085ux91ccux53f6ux53d8ux6362ux7684ux5b9aux4e49ux4e0eux5b58ux5728ux6761ux4ef6}

正变换与反变换（把非周期信号分解为连续频率的复正弦之和）：

\[
F(\mathrm{j}\omega)=\int_{-\infty}^{+\infty}f(t)e^{-\mathrm{j}\omega t}\,\mathrm{d}t
\]

\[
f(t)=\frac{1}{2\pi}\int_{-\infty}^{+\infty}F(\mathrm{j}\omega)
e^{\mathrm{j}\omega t}\,\mathrm{d}\omega
\]

记作
\(F(\mathrm{j}\omega)=\mathcal{F}[f(t)]\)，\(f(t)=\mathcal{F}^{-1}[F(\mathrm{j}\omega)]\)。

\textbf{存在条件（充分条件）}：\(f(t)\) 在 \((-\infty,+\infty)\)
上满足狄利克雷条件， 即绝对可积
\(\int_{-\infty}^{+\infty}|f(t)|\,\mathrm{d}t<\infty\)，
且只有有限个第一类间断点与有限个极值点。
注意这里的时间区间是整条实轴，不再是''一个周期''。

\textbf{放宽}：引入冲激函数后可定义\textbf{广义傅里叶变换}，
阶跃、直流、正弦等不满足绝对可积的信号也有变换式， 代价是频谱中出现
\(\delta(\omega)\)。

什么时候用：题目问''为什么可以用频域法''，答''满足狄利克雷条件''；
题目给了不绝对可积的信号，答''在广义函数意义下存在傅里叶变换''。

\subsubsection{常用信号的傅里叶变换对}\label{ux5e38ux7528ux4fe1ux53f7ux7684ux5085ux91ccux53f6ux53d8ux6362ux5bf9}

\begin{longtable}[]{@{}
  >{\raggedright\arraybackslash}p{(\linewidth - 6\tabcolsep) * \real{0.1395}}
  >{\raggedright\arraybackslash}p{(\linewidth - 6\tabcolsep) * \real{0.1860}}
  >{\raggedright\arraybackslash}p{(\linewidth - 6\tabcolsep) * \real{0.5349}}
  >{\raggedright\arraybackslash}p{(\linewidth - 6\tabcolsep) * \real{0.1395}}@{}}
\toprule\noalign{}
\begin{minipage}[b]{\linewidth}\raggedright
序号
\end{minipage} & \begin{minipage}[b]{\linewidth}\raggedright
\(f(t)\)
\end{minipage} & \begin{minipage}[b]{\linewidth}\raggedright
\(F(\mathrm{j}\omega)\)
\end{minipage} & \begin{minipage}[b]{\linewidth}\raggedright
条件
\end{minipage} \\
\midrule\noalign{}
\endhead
\bottomrule\noalign{}
\endlastfoot
1 & \(\delta(t)\) & \(1\) & 广义 \\
2 & \(\delta(t-t_0)\) & \(e^{-\mathrm{j}\omega t_0}\) & 广义 \\
3 & \(e^{-\alpha t}u(t)\) & \(\dfrac{1}{\alpha+\mathrm{j}\omega}\) &
\(\alpha>0\) \\
4 & \(e^{-\alpha|t|}\) & \(\dfrac{2\alpha}{\alpha^{2}+\omega^{2}}\) &
\(\alpha>0\) \\
5 & \(G_{\tau}(t)\) &
\(\tau\,\mathrm{Sa}\left(\dfrac{\omega\tau}{2}\right)\) & 门函数 \\
6 & \(\mathrm{Sa}(t)=\dfrac{\sin t}{t}\) & \(\pi G_{2}(\omega)\) &
广义 \\
7 & \(u(t)\) & \(\pi\delta(\omega)+\dfrac{1}{\mathrm{j}\omega}\) &
广义 \\
8 & \(1\) & \(2\pi\delta(\omega)\) & 广义 \\
9 & \(e^{\mathrm{j}\omega_c t}\) & \(2\pi\delta(\omega-\omega_c)\) &
广义 \\
10 & \(\cos\omega_c t\) &
\(\pi\bigl[\delta(\omega+\omega_c)+\delta(\omega-\omega_c)\bigr]\) &
广义 \\
\end{longtable}

说明两点：门函数 \(G_{\tau}(t)\) 在 \(|t|<\frac{\tau}{2}\) 内为
\(1\)、其余为 \(0\)； 序号 3、7、8
这类不绝对可积的信号，其变换式只在广义函数意义下成立，
答题时要写''借助冲激函数''。

\subsubsection{傅里叶变换的性质}\label{ux5085ux91ccux53f6ux53d8ux6362ux7684ux6027ux8d28}

\begin{longtable}[]{@{}
  >{\raggedright\arraybackslash}p{(\linewidth - 4\tabcolsep) * \real{0.3333}}
  >{\raggedright\arraybackslash}p{(\linewidth - 4\tabcolsep) * \real{0.3333}}
  >{\raggedright\arraybackslash}p{(\linewidth - 4\tabcolsep) * \real{0.3333}}@{}}
\toprule\noalign{}
\begin{minipage}[b]{\linewidth}\raggedright
性质
\end{minipage} & \begin{minipage}[b]{\linewidth}\raggedright
时域
\end{minipage} & \begin{minipage}[b]{\linewidth}\raggedright
频域
\end{minipage} \\
\midrule\noalign{}
\endhead
\bottomrule\noalign{}
\endlastfoot
线性 & \(af_1+bf_2\) & \(aF_1+bF_2\) \\
时移 & \(f(t-t_0)\) &
\(F(\mathrm{j}\omega)e^{-\mathrm{j}\omega t_0}\) \\
频移 & \(f(t)e^{\mathrm{j}\omega_c t}\) &
\(F\bigl(\mathrm{j}(\omega-\omega_c)\bigr)\) \\
尺度 & \(f(at)\) &
\(\dfrac{1}{|a|}F\left(\mathrm{j}\dfrac{\omega}{a}\right)\) \\
对称 & \(F(\mathrm{j}t)\) & \(2\pi f(-\omega)\) \\
微分 & \(\dfrac{\mathrm{d}^{n}f}{\mathrm{d}t^{n}}\) &
\((\mathrm{j}\omega)^{n}F(\mathrm{j}\omega)\) \\
积分 & \(\int_{-\infty}^{t}f(\tau)\,\mathrm{d}\tau\) &
\(\frac{F(\mathrm{j}\omega)}{\mathrm{j}\omega}+\pi F(0)\delta(\omega)\) \\
卷积 & \(f_1(t)*f_2(t)\) &
\(F_1(\mathrm{j}\omega)F_2(\mathrm{j}\omega)\) \\
相乘 & \(f_1(t)f_2(t)\) &
\(\dfrac{1}{2\pi}F_1(\mathrm{j}\omega)*F_2(\mathrm{j}\omega)\) \\
\end{longtable}

成立条件与常用推论：

\begin{itemize}
\tightlist
\item
  微分性质要求 \(f(t)\) 连续且
  \(\frac{\mathrm{d}^{n}f}{\mathrm{d}t^{n}}\) 存在傅里叶变换；若
  \(f(t)\) 有跳变，跳变点产生冲激，须逐个写出。
\item
  积分性质中的 \(\pi F(0)\delta(\omega)\) 项只在 \(F(0)\ne 0\)
  时需要保留， 即 \(f(t)\) 有非零直流分量时才有这一项。
\item
  频移性质的常用推论（调制定理）：时域乘一个余弦，频谱被搬移并减半：
\end{itemize}

\[
f(t)\cos\omega_c t\longleftrightarrow
\frac{1}{2}\bigl[F(\mathrm{j}(\omega+\omega_c))
+F(\mathrm{j}(\omega-\omega_c))\bigr]
\] - 尺度性质推论：时域压缩 \(a\) 倍，频域展宽 \(a\) 倍，即
``脉宽与带宽成反比''，与``周期信号的频带宽度''一节的结论一致。 -
相乘性质（频域卷积定理）必须带 \(\frac{1}{2\pi}\) 因子，
这是与卷积定理最容易搞混的地方。

\subsubsection{周期信号的傅里叶变换}\label{ux5468ux671fux4fe1ux53f7ux7684ux5085ux91ccux53f6ux53d8ux6362}

周期信号不满足绝对可积，但把它展开成指数级数后逐项变换， 再用
\(e^{\mathrm{j}n\omega_{0} t}\leftrightarrow 2\pi\delta(\omega-n\omega_{0})\)，得

\[
F(\mathrm{j}\omega)=2\pi\sum_{n=-\infty}^{+\infty}c_n\,\delta(\omega-n\omega_{0}),
\qquad \omega_{0}=\frac{2\pi}{T}
\]

即\textbf{周期信号的频谱是冲激串}，冲激的位置在 \(n\omega_{0}\)， 强度是
\(2\pi c_n\)。这正好对应''周期信号的频谱是离散谱''。

\textbf{均匀冲激序列}（最重要的一个）：

\[
\delta_T(t)=\sum_{k=-\infty}^{+\infty}\delta(t-kT)
\;\longleftrightarrow\;
\omega_{0}\sum_{n=-\infty}^{+\infty}\delta(\omega-n\omega_{0}),
\qquad \omega_{0}=\frac{2\pi}{T}
\]

推导思路三步：① \(\delta_T(t)\) 是周期信号，\(c_n=\frac{1}{T}\)； ②
代入冲激串公式得
\(F(\mathrm{j}\omega)=\frac{2\pi}{T}\sum_n\delta(\omega-n\omega_{0})\)；
③ 注意 \(\frac{2\pi}{T}=\omega_{0}\)，所以时域冲激间隔 \(T\)
与频域冲激间隔 \(\omega_{0}\) 成反比。

\textbf{单脉冲与周期化的关系}：设 \(f_0(t)\) 是 \(f(t)\)
的一个周期（其余补零）， 其变换为 \(F_0(\mathrm{j}\omega)\)，则周期化后

\[
c_n=\frac{1}{T}F_0(\mathrm{j}n\omega_{0})
\]

这条式子把''查表得的单脉冲变换''与''级数系数''直接连起来，很好用。

\subsubsection{系统的频率响应与零状态响应的频域求解}\label{ux7cfbux7edfux7684ux9891ux7387ux54cdux5e94ux4e0eux96f6ux72b6ux6001ux54cdux5e94ux7684ux9891ux57dfux6c42ux89e3}

对由微分方程描述的系统

\[
\sum_{k=0}^{n}a_ky^{(k)}(t)=\sum_{k=0}^{m}b_kf^{(k)}(t)
\]

定义\textbf{频率响应}（传输函数）

\[
H(\mathrm{j}\omega)=
\frac{b_m(\mathrm{j}\omega)^{m}+b_{m-1}(\mathrm{j}\omega)^{m-1}
+\cdots+b_1(\mathrm{j}\omega)+b_0}
{(\mathrm{j}\omega)^{n}+a_{n-1}(\mathrm{j}\omega)^{n-1}
+\cdots+a_1(\mathrm{j}\omega)+a_0}
\]

它与转移算子 \(H(p)\) 形式相同，只需把 \(p\) 换成 \(\mathrm{j}\omega\)；
但 \(H(\mathrm{j}\omega)\)
是代数式，可以自由做代数运算（约分、通分、部分分式）。

写成极坐标形式：

\[
H(\mathrm{j}\omega)=|H(\mathrm{j}\omega)|e^{\mathrm{j}\varphi(\omega)}
\]

\(|H(\mathrm{j}\omega)|\) 称\textbf{幅频特性}，\(\varphi(\omega)\)
称\textbf{相频特性}， 合称系统的频率特性。物理含义：系统对频率为
\(\omega\) 的复正弦 只改变幅度与相位，不改变频率。

\textbf{零状态响应}：设激励频谱
\(E(\mathrm{j}\omega)=\mathcal{F}[e(t)]\)，则

\[
R_{zs}(\mathrm{j}\omega)=H(\mathrm{j}\omega)E(\mathrm{j}\omega),
\qquad
y_{zs}(t)=\mathcal{F}^{-1}\bigl[H(\mathrm{j}\omega)E(\mathrm{j}\omega)\bigr]
=\frac{1}{2\pi}\int_{-\infty}^{+\infty}H(\mathrm{j}\omega)E(\mathrm{j}\omega)
e^{\mathrm{j}\omega t}\,\mathrm{d}\omega
\]

成立条件：系统线性时不变、零状态（初始储能为零）， 且
\(H(\mathrm{j}\omega)\) 存在，即 \(H(s)\) 的收敛域包含虚轴，
等价于极点全在左半平面（系统稳定）。若系统不稳定（存在右半平面极点），
收敛域不含虚轴，\(H(\mathrm{j}\omega)\) 根本不存在，此时只能用
拉普拉斯变换求解（见第 5 章）。

\textbf{周期信号}激励时不必做积分，逐项加权即可：

\[
e(t)=\sum_n c_ne^{\mathrm{j}n\omega_{0} t}
\;\Longrightarrow\;
y_{zs}(t)=\sum_n c_nH(\mathrm{j}n\omega_{0})e^{\mathrm{j}n\omega_{0} t}
\]

对实正弦激励 \(A\cos(\omega t+\theta)\)，响应为

\[
A|H(\mathrm{j}\omega)|\cos\bigl(\omega t+\theta+\varphi(\omega)\bigr)
\]

即''幅度乘幅频、相位加相频''。

三种求 \(H(\mathrm{j}\omega)\) 的途径（结果相同，按题目条件选）： ①
由微分方程系数直接写出；② 由 \(h(t)\) 做傅里叶变换； ③
由电路做相量稳态分析（阻抗法）求输出与输入之比。

\subsubsection{无失真传输条件}\label{ux65e0ux5931ux771fux4f20ux8f93ux6761ux4ef6}

\textbf{无失真}指输出波形与输入波形形状完全相同，只允许幅度放大和整体延时：

\[
y(t)=Kf(t-t_d),\qquad K>0,\ t_d>0
\]

两边取傅里叶变换得
\(Y(\mathrm{j}\omega)=K e^{-\mathrm{j}\omega t_d}F(\mathrm{j}\omega)\)，
故无失真传输的\textbf{充要条件}是

\[
H(\mathrm{j}\omega)=Ke^{-\mathrm{j}\omega t_d}
\]

拆成幅频与相频两条：

\begin{enumerate}
\def\labelenumi{\arabic{enumi}.}
\tightlist
\item
  幅频特性为常数：\(|H(\mathrm{j}\omega)|=K\)，对所有频率一视同仁；
\item
  相频特性为过原点的直线：\(\varphi(\omega)=-\omega t_d\)， 群延时
  \(\tau_g=-\frac{\mathrm{d}\varphi}{\mathrm{d}\omega}=t_d\) 为常数。
\end{enumerate}

两条必须同时满足。幅度不失真但相位非线性，会产生\textbf{相位失真}
（波形畸变但频谱幅度不变）；相位线性但幅度不平坦，会产生
\textbf{幅度失真}。实际系统只要在信号频带内近似满足这两条即可。

\subsubsection{理想低通滤波器}\label{ux7406ux60f3ux4f4eux901aux6ee4ux6ce2ux5668}

\textbf{幅频}在 \(|\omega|\le\omega_c\) 内为 \(1\)、其余为
\(0\)，\textbf{相频}线性， 即

\[
H(\mathrm{j}\omega)=G_{2\omega_c}(\omega)e^{-\mathrm{j}\omega t_d}
\]

\textbf{冲激响应}：由对称性，\(G_{2\omega_c}(\omega)\) 的逆变换为
\(\frac{\omega_c}{\pi}\mathrm{Sa}(\omega_c t)\)，加时移得

\[
h(t)=\frac{\omega_c}{\pi}\,\mathrm{Sa}\bigl[\omega_c(t-t_d)\bigr]
=\frac{\sin\omega_c(t-t_d)}{\pi(t-t_d)}
\]

它在 \(t<0\) 就有非零值，说明理想低通\textbf{非因果}。

\textbf{阶跃响应}：把 \(u(t)\) 与 \(h(t)\) 卷积，得

\[
g(t)=\frac{1}{2}+\frac{1}{\pi}\mathrm{Si}\bigl[\omega_c(t-t_d)\bigr],
\qquad
\mathrm{Si}(y)=\int_{0}^{y}\frac{\sin x}{x}\,\mathrm{d}x
\]

\(\mathrm{Si}(\pm\infty)=\pm\frac{\pi}{2}\)，故
\(g(-\infty)=0\)、\(g(+\infty)=1\)。 波形有\textbf{过冲}（约
\(9\%\)，吉布斯现象）与振铃，且 \(g(t)\) 在 \(t<t_d\)
时已有非零值，同样非因果。

\textbf{上升时间与带宽的矛盾}：由主瓣估算，\(g(t)\) 从近似 \(0\)
升到近似 \(1\) 的时间约

\[
t_r\approx\frac{\pi}{\omega_c}=\frac{1}{2B_f},\qquad B_f=\frac{\omega_c}{2\pi}
\]

即 \(t_rB_f\approx 0.5\)：通带越宽，上升越快、越能跟上快速变化；
通带越窄，响应越迟缓。这就是''带宽与响应速度成反比''。

\subsubsection{物理可实现性：佩利-维纳准则}\label{ux7269ux7406ux53efux5b9eux73b0ux6027ux4f69ux5229-ux7ef4ux7eb3ux51c6ux5219}

时域判据：因果系统的冲激响应必须满足 \(h(t)=0\)（\(t<0\)）。
频域判据（佩利-维纳准则，必要且几乎充分）：

\[
\int_{-\infty}^{+\infty}
\frac{\bigl|\ln|H(\mathrm{j}\omega)|\bigr|}{1+\omega^{2}}\,\mathrm{d}\omega<\infty
\]

三条推论，考试常直接考这三条：

\begin{enumerate}
\def\labelenumi{\arabic{enumi}.}
\tightlist
\item
  \(|H(\mathrm{j}\omega)|\) 可以在某些孤立频率点上为零，但不能在
  一段频带内恒为零（否则积分发散）；
\item
  \(|H(\mathrm{j}\omega)|\) 的衰减速度不能快于指数型，
  所以''矩形幅频''\,``高斯型幅频''都不可实现；
\item
  理想低通、理想高通、理想带通均\textbf{不可实现}；实际滤波器只能
  逼近它，允许过渡带与一定的幅频起伏。
\end{enumerate}

允许的幅频形状：只要 \(|H(\mathrm{j}\omega)|\) 不为零且衰减不过快，
总能用有理函数（由 \(R\)、\(L\)、\(C\) 与受控源构成的网络）逼近实现。

\subsubsection{抽样定理与信号恢复}\label{ux62bdux6837ux5b9aux7406ux4e0eux4fe1ux53f7ux6062ux590d}

设 \(f(t)\) 是\textbf{带限信号}，最高频率为 \(f_m\)（对应
\(\omega_m=2\pi f_m\)）， 用周期 \(T_s\) 的冲激串抽样：

\[
f_s(t)=f(t)\delta_T(t)=\sum_{k}f(kT_s)\delta(t-kT_s)
\]

两边取傅里叶变换（用相乘性质或直接展开）得

\[
F_s(\mathrm{j}\omega)=\frac{1}{T_s}\sum_{k=-\infty}^{+\infty}
F\bigl(\mathrm{j}(\omega-k\omega_s)\bigr),\qquad \omega_s=\frac{2\pi}{T_s}
\]

即抽样后频谱是原频谱以 \(\omega_s\) 为周期\textbf{平移复制}。

\textbf{抽样定理（奈奎斯特）}：若 \(f(t)\) 带限于 \(f_m\)，则当

\[
f_s=\frac{1}{T_s}\ge 2f_m
\quad\Longleftrightarrow\quad
\omega_s\ge 2\omega_m
\]

时，\(F_s(\mathrm{j}\omega)\) 的各次复制不重叠，可用一个理想低通
完整恢复 \(f(t)\)。\(2f_m\) 称奈奎斯特率，\(\frac{1}{2f_m}\)
称奈奎斯特间隔。

\textbf{恢复}：取理想低通
\(H(\mathrm{j}\omega)=T_sG_{2\omega_c}(\omega)\)， 截止频率满足

\[
\omega_m<\omega_c<\omega_s-\omega_m
\]

则
\(F_s(\mathrm{j}\omega)H(\mathrm{j}\omega)=F(\mathrm{j}\omega)\)，时域上即

\[
f(t)=\sum_{k=-\infty}^{+\infty}f(kT_s)\,
\mathrm{Sa}\bigl[\omega_c(t-kT_s)\bigr]
\]

当 \(\omega_c=\frac{\omega_s}{2}\) 时就是标准的
\(\mathrm{Sa}\left[\frac{\omega_s}{2}(t-kT_s)\right]\) 内插公式。

\textbf{混叠}：若
\(f_s<2f_m\)，复制谱相互重叠，高频成分被''折叠''到低频，
无法用线性滤波去除，称为混叠失真。工程上要么提高 \(f_s\)，
要么在抽样前加\textbf{抗混叠前置低通}把 \(f_m\) 压下来。

\subsection{典型题型与解法}\label{ux5178ux578bux9898ux578bux4e0eux89e3ux6cd5-3}

\subsubsection{题型 4-1
由时域表达式或波形直接积分求傅里叶变换}\label{ux9898ux578b-4-1-ux7531ux65f6ux57dfux8868ux8fbeux5f0fux6216ux6ce2ux5f62ux76f4ux63a5ux79efux5206ux6c42ux5085ux91ccux53f6ux53d8ux6362}

\textbf{题型特征}：给分段常数波形（单个或成对的矩形脉冲）或含 \(u(t)\)
的 表达式，要求 \(F(\mathrm{j}\omega)\)。

\textbf{解题步骤}：

\begin{enumerate}
\def\labelenumi{\arabic{enumi}.}
\tightlist
\item
  把波形写成门函数之和（用平移），或写成阶跃函数之差；
\item
  套
  \(G_{\tau}(t)\leftrightarrow\tau\mathrm{Sa}\left(\frac{\omega\tau}{2}\right)\)
  与时移性质，逐项变换；
\item
  用和差化积把结果整理成实函数，如
  \(2\sin\alpha\cos\beta=\sin(\alpha+\beta)+\sin(\alpha-\beta)\)；
\item
  校验：奇对称波形其 \(F(\mathrm{j}\omega)\) 应为纯虚或纯实且奇，
  偶对称波形应为实偶函数。
\end{enumerate}

\textbf{例题}（【附加作业4 第4题】）：\(f(t)\) 在 \([-3,-1]\) 与
\([1,3]\) 上为 \(1\)、 其余为 \(0\)（两个宽 \(2\)、高 \(1\)
的矩形脉冲），求 \(F(\mathrm{j}\omega)\)。

解（方法一，门函数叠加）：两个脉冲可写成一个门函数向左、向右各平移
\(2\)：

\[
f(t)=G_2(t+2)+G_2(t-2)
\]

由
\(G_{\tau}(t)\leftrightarrow\tau\mathrm{Sa}\left(\frac{\omega\tau}{2}\right)\)，
取 \(\tau=2\) 得
\(G_2(t)\leftrightarrow 2\mathrm{Sa}(\omega)\)。再用时移性质

\[
G_2(t+2)\leftrightarrow 2\mathrm{Sa}(\omega)e^{\mathrm{j}2\omega},\qquad
G_2(t-2)\leftrightarrow 2\mathrm{Sa}(\omega)e^{-\mathrm{j}2\omega}
\]

相加并用欧拉公式
\(e^{\mathrm{j}\theta}+e^{-\mathrm{j}\theta}=2\cos\theta\)：

\[
F(\mathrm{j}\omega)=2\mathrm{Sa}(\omega)\left(e^{\mathrm{j}2\omega}
+e^{-\mathrm{j}2\omega}\right)=4\mathrm{Sa}(\omega)\cos 2\omega
\]

再把 \(\mathrm{Sa}(\omega)=\frac{\sin\omega}{\omega}\) 代入并和差化积：

\[
F(\mathrm{j}\omega)=\frac{4\sin\omega\cos 2\omega}{\omega}
=\frac{2}{\omega}\bigl[\sin 3\omega-\sin\omega\bigr]
\]

结果在 \(\omega=0\) 处取极限 \(F(0)=4\)（等于波形面积 \(2+2\)），与
\(\frac{2}{\omega}(\sin3\omega-\sin\omega)\to 2\cdot(3-1)=4\) 一致，
可作校验。

解（方法二，阶跃叠加）：\(f(t)=u(t+3)-u(t+1)+u(t-1)-u(t-3)\)，
代入定义式逐项积分亦可，结果相同。

\textbf{同类题}：

\begin{itemize}
\tightlist
\item
  【附加作业4 第1题】概念题：充分条件是狄利克雷条件，变换对为
  \(F(\mathrm{j}\omega)=\int f(t)e^{-\mathrm{j}\omega t}\mathrm{d}t\) 与
  \(f(t)=\frac{1}{2\pi}\int F(\mathrm{j}\omega)e^{\mathrm{j}\omega t}\mathrm{d}\omega\)，
  注意反变换的 \(\frac{1}{2\pi}\) 不能丢。
\item
  【附加作业4 第6题】\(e(t)=\frac{\sin 2\pi t}{2\pi t}\)：先由
  \(G_{\tau}(t)\leftrightarrow\tau\mathrm{Sa}\left(\frac{\omega\tau}{2}\right)\)
  与对称性得
  \(2\pi G_{4\pi}(\omega)\leftrightarrow 4\pi\mathrm{Sa}(2\pi t)\)，
  两边除 \(2\pi\) 即得
  \(E(\mathrm{j}\omega)=\frac{1}{2}G_{4\pi}(\omega)\)。
\end{itemize}

\subsubsection{题型 4-2
利用微积分特性求折线信号的傅里叶变换}\label{ux9898ux578b-4-2-ux5229ux7528ux5faeux79efux5206ux7279ux6027ux6c42ux6298ux7ebfux4fe1ux53f7ux7684ux5085ux91ccux53f6ux53d8ux6362}

\textbf{题型特征}：波形是梯形、三角等\textbf{分段线性}函数（求导后变成矩形、
二次求导后变成冲激），要求 \(F(\mathrm{j}\omega)\)。

\textbf{解题步骤}：

\begin{enumerate}
\def\labelenumi{\arabic{enumi}.}
\tightlist
\item
  先把波形写成 \(t\) 与 \(u(t-t_i)\) 的乘积之和，方便逐项求导；
\item
  求一阶导：\(t\,u(t-t_0)\) 的导数是 \(u(t-t_0)+t\delta(t-t_0)\)， 而
  \(t\delta(t-t_0)=t_0\delta(t-t_0)\)，别忘了冲激项；
\item
  若一阶导仍是矩形，再求一次得到冲激串；
\item
  用
  \(f^{(n)}(t)\leftrightarrow(\mathrm{j}\omega)^{n}F(\mathrm{j}\omega)\)
  反解，即
  \(F(\mathrm{j}\omega)=(\mathrm{j}\omega)^{-n}\mathcal{F}[f^{(n)}(t)]\)；
\item
  用 \(F(0)=\int f\,\mathrm{d}t\)（波形面积）检验结果。
\end{enumerate}

\textbf{例题}：三角脉冲 \(f(t)=1-|t|\)（\(|t|<1\)）、其余为 \(0\)，
关键点 \((0,1)\)、\((\pm1,0)\)，用微积分特性求其变换。

解：分段表达式为
\(f(t)=(1+t)\)（\(-1\le t\le 0\)）、\(f(t)=(1-t)\)（\(0\le t\le 1\)）。
一阶导是两段矩形：\(f'(t)=1\) 于 \((-1,0)\)、\(f'(t)=-1\) 于
\((0,1)\)，即

\[
f'(t)=G_1\left(t+\frac{1}{2}\right)-G_1\left(t-\frac{1}{2}\right)
\]

再求导得冲激串：

\[
f''(t)=\delta(t+1)-2\delta(t)+\delta(t-1)
\]

由微分性质及
\(\delta(t-t_0)\leftrightarrow e^{-\mathrm{j}\omega t_0}\)：

\[
(\mathrm{j}\omega)^{2}F(\mathrm{j}\omega)
=e^{\mathrm{j}\omega}-2+e^{-\mathrm{j}\omega}
=2\cos\omega-2
\]

所以

\[
F(\mathrm{j}\omega)=-\frac{2\cos\omega-2}{\omega^{2}}
=\frac{2(1-\cos\omega)}{\omega^{2}}
=\left(\frac{\sin(\omega/2)}{\omega/2}\right)^{2}
=\mathrm{Sa}^{2}\!\left(\frac{\omega}{2}\right)
\]

校验：\(F(0)=1\)，而波形面积为 \(1\)（底 \(2\)、高 \(1\)
的三角形面积恰为 \(1\)）， 一致。该结果也是必背变换对：三角脉冲对应
\(\mathrm{Sa}^{2}\)。

\textbf{同类题}：

\begin{itemize}
\tightlist
\item
  【附加作业4 第5题】梯形信号关键点
  \((-2,0)\)、\((-1,1)\)、\((1,1)\)、\((2,0)\)： 先写
  \(f(t)=(t+2)u(t+2)-(t+1)u(t+1)-(t-1)u(t-1)+(t-2)u(t-2)\)， 求导时用
  \(t\delta(t-t_0)=t_0\delta(t-t_0)\) 消掉四个冲激， 得
  \(f'(t)=u(t+2)-u(t+1)-u(t-1)+u(t-2)\)，再求导得四个冲激， 最终
  \(F(\mathrm{j}\omega)=\frac{2[\cos\omega-\cos 2\omega]}{\omega^{2}}\)。
\end{itemize}

\subsubsection{题型 4-3
由频率响应求冲激响应并求零状态响应}\label{ux9898ux578b-4-3-ux7531ux9891ux7387ux54cdux5e94ux6c42ux51b2ux6fc0ux54cdux5e94ux5e76ux6c42ux96f6ux72b6ux6001ux54cdux5e94}

\textbf{题型特征}：给 \(H(\mathrm{j}\omega)\)
的有理分式（或给微分方程）， 要求 \(h(t)\)，再给一个指数型激励求
\(y_{zs}(t)\)。

\textbf{解题步骤}：

\begin{enumerate}
\def\labelenumi{\arabic{enumi}.}
\tightlist
\item
  把 \(\mathrm{j}\omega\) 看作 \(s\)，对 \(H\) 做部分分式展开，每项化成
  \(\frac{A}{s+\alpha}\) 的标准形；
\item
  逆变换：\(\frac{A}{s+\alpha}\leftrightarrow Ae^{-\alpha t}u(t)\)
  （要求 \(\alpha>0\)，否则系统不稳定、\(h(t)\) 不存在）；
\item
  求激励频谱：\(e^{-\alpha t}u(t)\leftrightarrow\frac{1}{\alpha+\mathrm{j}\omega}\)；
\item
  相乘再展开：\(R_{zs}=H\cdot E\) 的部分分式，注意\textbf{不能}只保留
  一部分项，凡分母出现的因式都要给出对应项；
\item
  若某个留数为零，说明该项确实为零（例如激励极点被 \(H\) 的零点抵消），
  要点明这是巧合而不是写法错误。
\end{enumerate}

\textbf{例题}（【附加作业4 第7题】）：已知

\[
H(\mathrm{j}\omega)=\frac{2\mathrm{j}\omega+3}{-\omega^{2}+\mathrm{j}3\omega+2}
\]

求冲激响应，并用频域法求 \(e(t)=e^{-1.5t}u(t)\) 时的 \(y_{zs}(t)\)。

解：记 \(s=\mathrm{j}\omega\)，则 \(-\omega^{2}=s^{2}\)，故

\[
H(s)=\frac{2s+3}{s^{2}+3s+2}=\frac{2s+3}{(s+1)(s+2)}
\]

部分分式：\(H(s)=\frac{A}{s+1}+\frac{B}{s+2}\)，
\(A=\frac{2s+3}{s+2}\Big|_{s=-1}=\frac{1}{1}=1\)，
\(B=\frac{2s+3}{s+1}\Big|_{s=-2}=\frac{-1}{-1}=1\)。所以

\[
H(\mathrm{j}\omega)=\frac{1}{\mathrm{j}\omega+1}+\frac{1}{\mathrm{j}\omega+2}
\]

逆变换得冲激响应

\[
h(t)=\left(e^{-t}+e^{-2t}\right)u(t)
\]

激励频谱 \(E(\mathrm{j}\omega)=\frac{1}{\mathrm{j}\omega+1.5}\)，于是

\[
R_{zs}(\mathrm{j}\omega)=H(\mathrm{j}\omega)E(\mathrm{j}\omega)
=\frac{2s+3}{(s+1)(s+2)(s+1.5)}
\]

三个极点 \(-1\)、\(-2\)、\(-1.5\)，逐项求留数：

\[
A_{1}=\frac{2s+3}{(s+2)(s+1.5)}\Big|_{s=-1}=\frac{1}{1\times 0.5}=2
\]

\[
A_{2}=\frac{2s+3}{(s+1)(s+1.5)}\Big|_{s=-2}=\frac{-1}{(-1)\times(-0.5)}=-2
\]

\[
A_{3}=\frac{2s+3}{(s+1)(s+2)}\Big|_{s=-1.5}=\frac{0}{(-0.5)\times 0.5}=0
\]

第三项之所以为零，是因为分子 \(2s+3\) 恰在 \(s=-1.5\) 处有零点，
把激励极点抵消了------这是本题的巧合，不是漏项。故

\[
R_{zs}(\mathrm{j}\omega)=\frac{2}{\mathrm{j}\omega+1}
-\frac{2}{\mathrm{j}\omega+2}
\]

逆变换得

\[
y_{zs}(t)=2\left(e^{-t}-e^{-2t}\right)u(t)
\]

校验：\(t\to\infty\) 时 \(y_{zs}\to 0\)，与系统稳定、激励有阻尼一致； 又
\(y_{zs}(0^{+})=0\)，说明无跳变。

\textbf{同类题}：

\begin{itemize}
\tightlist
\item
  【附加作业4 第8题】由 \(r''+3r'+2r=e\) 直接写
  \(H(\mathrm{j}\omega)=\frac{1}{(\mathrm{j}\omega)^{2}+3\mathrm{j}\omega+2}\)，
  \(E=\frac{1}{\mathrm{j}\omega+3}\)，三个极点 \(-1,-2,-3\) 的留数分别为
  \(0.5\)、\(-1\)、\(0.5\)，得
  \(y_{zs}(t)=\left(0.5e^{-t}-e^{-2t}+0.5e^{-3t}\right)u(t)\)。
\end{itemize}

\subsubsection{题型 4-4
调制加理想滤波求系统输出}\label{ux9898ux578b-4-4-ux8c03ux5236ux52a0ux7406ux60f3ux6ee4ux6ce2ux6c42ux7cfbux7edfux8f93ux51fa}

\textbf{题型特征}：框图先乘
\(\cos\omega_0t\)（调制），再经理想带通/低通滤波， 给滤波器的幅频特性与
\(\varphi(\omega)=0\)，求输出。

\textbf{解题步骤}：

\begin{enumerate}
\def\labelenumi{\arabic{enumi}.}
\item
  求输入频谱 \(E(\mathrm{j}\omega)\)（常是门函数或 \(\mathrm{Sa}\)
  型）；
\item
  用频移性质或相乘性质求乘法器输出，频谱左右各搬移 \(\pm\omega_0\)
  后减半：

  \[
  E_1(\mathrm{j}\omega)=\frac{1}{2}\bigl[E(\mathrm{j}(\omega+\omega_0))
  +E(\mathrm{j}(\omega-\omega_0))\bigr]
  \]
\item
  画出 \(E_1\) 的频谱，与 \(|H(\mathrm{j}\omega)|\) 的通带位置对照，
  把落在通带内的部分''切出来''；
\item
  逆变换时可用对称性：频域里宽为 \(\tau\) 的门函数，对应时域
  \(\frac{\tau}{2\pi}\mathrm{Sa}\left(\frac{\tau t}{2}\right)\)；
  逐块逆变换后再用欧拉公式合成实信号。
\end{enumerate}

\textbf{例题}（【附加作业4 第9题】）：
\(e(t)=\frac{\sin 2\pi t}{2\pi t}\)，先与 \(\cos 1000t\) 相乘，
再通过理想带通滤波器。带通幅频在 \(999<|\omega|<1001\) 内为 \(1\)、
其余为 \(0\)，相位特性 \(\varphi(\omega)=0\)。求输出 \(y(t)\)。

解：第一步，求输入频谱。由题型 4-1 同类题的结论

\[
E(\mathrm{j}\omega)=\frac{1}{2}G_{4\pi}(\omega)
\]

即幅度 \(\frac{1}{2}\)、占据 \(|\omega|<2\pi\) 的矩形谱。

第二步，调制。由
\(\mathcal{F}[\cos 1000t]=\pi[\delta(\omega+1000)+\delta(\omega-1000)]\)
与相乘性质：

\[
E_1(\mathrm{j}\omega)=\frac{1}{2\pi}E(\mathrm{j}\omega)
*\pi\bigl[\delta(\omega+1000)+\delta(\omega-1000)\bigr]
\]

利用 \(\delta\) 的筛选性质
\(F(\mathrm{j}\omega)*\delta(\omega-\omega_0)=F(\mathrm{j}(\omega-\omega_0))\)，得

\[
E_1(\mathrm{j}\omega)=\frac{1}{4}\bigl[G_{4\pi}(\omega+1000)
+G_{4\pi}(\omega-1000)\bigr]
\]

即原矩形谱被搬到 \(\pm1000\) 处，每块宽度 \(4\pi\)，幅度
\(\frac{1}{4}\)。

第三步，滤波。通带 \(999<|\omega|<1001\) 宽度为 \(2\)， 而搬移后的每块宽
\(4\pi\approx12.57\)，故每块只被切出中间宽 \(2\) 的一段：

\[
R(\mathrm{j}\omega)=E_1(\mathrm{j}\omega)H(\mathrm{j}\omega)
=\frac{1}{4}G_2(\omega+1000)+\frac{1}{4}G_2(\omega-1000)
\]

第四步，逆变换。由 \(G_2(\omega)\) 的逆变换为
\(\frac{\sin t}{\pi t}=\frac{1}{\pi}\mathrm{Sa}(t)\)，再用频移性质
\(F(\mathrm{j}(\omega-\omega_0))\leftrightarrow f(t)e^{\mathrm{j}\omega_0 t}\)：

\[
y(t)=\frac{1}{4}\cdot\frac{\sin t}{\pi t}e^{-\mathrm{j}1000t}
+\frac{1}{4}\cdot\frac{\sin t}{\pi t}e^{\mathrm{j}1000t}
\]

用欧拉公式合并：

\[
y(t)=\frac{\sin t}{2\pi t}\cos 1000t
=\frac{1}{2\pi}\mathrm{Sa}(t)\cos 1000t,
\qquad -\infty<t<+\infty
\]

物理意义：输出是一个包络为 \(\mathrm{Sa}(t)\) 的
\(1000\ \mathrm{rad/s}\)
窄带信号，即把低频信号''搬''到了载频附近，这正是调制的效果。

\textbf{同类题}：

\begin{itemize}
\tightlist
\item
  【附加作业4 第2题】用频域卷积定理导出 \(\cos\omega_ct\cdot u(t)\)：
  \(\cos\omega_ct\leftrightarrow\pi[\delta(\omega+\omega_c)+\delta(\omega-\omega_c)]\)、
  \(u(t)\leftrightarrow\frac{1}{\mathrm{j}\omega}+\pi\delta(\omega)\)，
  卷积后得
  \(\frac{\pi}{2}[\delta(\omega+\omega_c)+\delta(\omega-\omega_c)]\)
  加上 \(\frac{\mathrm{j}\omega}{\omega_c^{2}-\omega^{2}}\)； 关键是用
  \(\frac{1}{\mathrm{j}\omega}\) 与 \(\delta\) 卷积即得频移项。
\end{itemize}

\subsubsection{题型 4-5
利用对称性与频域卷积定理导出新变换对}\label{ux9898ux578b-4-5-ux5229ux7528ux5bf9ux79f0ux6027ux4e0eux9891ux57dfux5377ux79efux5b9aux7406ux5bfcux51faux65b0ux53d8ux6362ux5bf9}

\textbf{题型特征}：题目要求''由已知变换对导出另一对的变换''，
或''求奇偶分量的频谱''，本质考性质的对偶使用。

\textbf{解题步骤}：

\begin{enumerate}
\def\labelenumi{\arabic{enumi}.}
\tightlist
\item
  认出已知对中谁像谁：\(F(\mathrm{j}\omega)\) 与 \(f(t)\)
  交换位置用对称性；
\item
  套对称性 \(F(\mathrm{j}t)\leftrightarrow 2\pi f(-\omega)\)， 若 \(f\)
  为偶函数则 \(f(-\omega)=f(\omega)\)，结果最简洁；
\item
  时域乘积改用频域卷积（带 \(\frac{1}{2\pi}\)）， 频域乘积改用时域卷积；
\item
  用奇偶分解式 \(f_e(t)=\frac{f(t)+f(-t)}{2}\)、
  \(f_o(t)=\frac{f(t)-f(-t)}{2}\) 与
  \(F(-\mathrm{j}\omega)=R(\omega)-\mathrm{j}X(\omega)\) 对应。
\end{enumerate}

\textbf{例题}：已知
\(e^{-\alpha|t|}\leftrightarrow\frac{2\alpha}{\alpha^{2}+\omega^{2}}\)
（\(\alpha>0\)），用对称性求 \(f(t)=\frac{1}{1+t^{2}}\) 的傅里叶变换。

解：先整理已知对，使其形如 \(F(\mathrm{j}\omega)\) 为实偶函数：

\[
F(\mathrm{j}\omega)=\frac{2\alpha}{\alpha^{2}+\omega^{2}}
\;\longleftrightarrow\;
f(t)=e^{-\alpha|t|}
\]

由对称性 \(F(\mathrm{j}t)\leftrightarrow 2\pi f(-\omega)\)，取
\(\alpha=1\) 得

\[
\frac{2}{1+t^{2}}\;\longleftrightarrow\;2\pi e^{-|\omega|}
\]

两边同除 \(2\)：

\[
\frac{1}{1+t^{2}}\;\longleftrightarrow\;\pi e^{-|\omega|}
\]

结果校验：\(\frac{1}{1+t^{2}}\) 是实偶函数，其频谱必为实偶函数，
\(\pi e^{-|\omega|}\) 满足；且 \(F(0)=\pi\)，而
\(\int_{-\infty}^{+\infty}\frac{\mathrm{d}t}{1+t^{2}}=\pi\)，一致。

\textbf{同类题}：

\begin{itemize}
\tightlist
\item
  【附加作业4 第2题】频域卷积定理的具体用法见题型 4-4 的同类题。
\item
  【附加作业4 第3题】\(f(t)\) 为实函数、
  \(F(\mathrm{j}\omega)=R(\omega)+\mathrm{j}X(\omega)\) 时，用
  \(F(-\mathrm{j}\omega)=R(\omega)-\mathrm{j}X(\omega)\)
  代入奇偶分解即得
  \(F_e(\mathrm{j}\omega)=R(\omega)\)、\(F_o(\mathrm{j}\omega)=\mathrm{j}X(\omega)\)；
  结论要记牢：\textbf{实偶分量对应频谱的实部，实奇分量对应频谱的虚部}。
\end{itemize}

\subsection{易错点}\label{ux6613ux9519ux70b9-3}

\begin{enumerate}
\def\labelenumi{\arabic{enumi}.}
\tightlist
\item
  \textbf{反变换漏掉 \(\frac{1}{2\pi}\)} → 写成
  \(f(t)=\frac{1}{2\pi}\int F(\mathrm{j}\omega)e^{\mathrm{j}\omega t}\mathrm{d}\omega\)
  → 正变换积分里没有这个因子，很容易顺手漏掉。
\item
  \textbf{门函数变换漏掉 \(\frac{\tau}{2}\)} → 用
  \(G_{\tau}(t)\leftrightarrow\tau\mathrm{Sa}\left(\frac{\omega\tau}{2}\right)\)
  → \(\mathrm{Sa}\) 的自变量是半宽乘 \(\omega\)，不是全宽乘 \(\omega\)。
\item
  \textbf{阶跃函数变换漏掉 \(\pi\delta(\omega)\)} → 必须写
  \(u(t)\leftrightarrow\frac{1}{\mathrm{j}\omega}+\pi\delta(\omega)\) →
  少了这一项，凡含阶跃的题目（如作业4第2题）都错一项。
\item
  \textbf{频移符号搞反} → \(f(t)e^{\mathrm{j}\omega_c t}\) 对应
  \(F(\mathrm{j}(\omega-\omega_c))\) → ``时域乘指数、频域右移''，可用
  \(\delta(\omega-\omega_c)\) 对照记忆。
\item
  \textbf{卷积定理与相乘性质混用} → 时域相乘配频域卷积，且要
  \(\frac{1}{2\pi}\) 因子 → 两个变换的因子不对称，只能靠记。
\item
  \textbf{微积分性质漏掉跳变产生的冲激} → \(f(t)\) 有跳变时先写出
  \(f'(t)\) 的冲激项再变换 → 逐段求导并保留 \(t_0\delta(t-t_0)\)
  这类项。
\item
  \textbf{部分分式只保留一部分项} → 分母的每个因式都要给出一项，
  留数为零也要说明为零的原因 → 漏项在时域就是丢掉一个响应分量。
\item
  \textbf{把 \(H(p)\) 与 \(H(\mathrm{j}\omega)\) 混为一谈} →
  只有稳定系统 才能把 \(p\) 换成 \(\mathrm{j}\omega\) 得到频率响应 →
  \(H(\mathrm{j}\omega)\) 是 \(H(p)\) 在虚轴上的取值。
\item
  \textbf{判断无失真只看幅频} → 必须同时满足 \(|H|=K\) 常数与
  \(\varphi(\omega)=-\omega t_d\) 线性 →
  幅度不失真但相位非线性仍会畸变。
\item
  \textbf{抽样定理中 \(f_s\) 与 \(\omega_s\)、\(f_m\) 与 \(\omega_m\)
  混用} → \(f_s\ge 2f_m\) 或 \(\omega_s\ge 2\omega_m\)，两式因子都为
  \(2\)，但单位不同 → 带限信号可用矩形脉冲的''第一个零点''估计 \(f_m\)。
\item
  \textbf{认为理想低通可以物理实现} → 用佩利-维纳准则判定，
  \(|H(\mathrm{j}\omega)|\) 不能在一段频带内恒为零 → 理想低通的 \(h(t)\)
  在 \(t<0\) 已有非零值，非因果。
\end{enumerate}

\section{连续时间系统的复频域分析}\label{ux8fdeux7eedux65f6ux95f4ux7cfbux7edfux7684ux590dux9891ux57dfux5206ux6790}

\subsection{考点清单}\label{ux8003ux70b9ux6e05ux5355-4}

\begin{longtable}[]{@{}
  >{\raggedright\arraybackslash}p{(\linewidth - 4\tabcolsep) * \real{0.2500}}
  >{\raggedright\arraybackslash}p{(\linewidth - 4\tabcolsep) * \real{0.4167}}
  >{\centering\arraybackslash}p{(\linewidth - 4\tabcolsep) * \real{0.3333}}@{}}
\toprule\noalign{}
\begin{minipage}[b]{\linewidth}\raggedright
考点
\end{minipage} & \begin{minipage}[b]{\linewidth}\raggedright
常见考法
\end{minipage} & \begin{minipage}[b]{\linewidth}\centering
重要度
\end{minipage} \\
\midrule\noalign{}
\endhead
\bottomrule\noalign{}
\endlastfoot
拉普拉斯变换的定义 & 写出单边与双边 \(F(s)\)，说明积分下限与各自适用范围
& ★★★ \\
收敛域与信号类型的关系 & 给信号判收敛域；给收敛域判左边、右边或双边信号
& ★★★ \\
常用变换对 & 查表直接写出 \(F(s)\) 或 \(f(t)\) & ★★★ \\
变换性质 & 时移、频移、尺度、时域微分、时域积分、卷积的综合运用 & ★★★ \\
初值与终值定理 & 求
\(f(0^{+})\)、\(f(+\infty)\)（如求阶跃输入下的稳态输出） & ★★ \\
部分分式展开法 & 真分式、共轭复根、重根、假分式四类反变换 & ★★★ \\
\(s\) 域求解微分方程 & 给微分方程与初始状态，求全响应并拆零输入与零状态
& ★★★ \\
自然响应与受迫响应 & 按响应项的极点来源分类 & ★★ \\
电路的 \(s\) 域模型 & 画运算电路，求系统函数或某支路响应 & ★★ \\
线性系统的模拟 & 由微分方程或 \(H(s)\) 画直接模拟框图 & ★★ \\
\end{longtable}

\subsection{核心概念与公式}\label{ux6838ux5fc3ux6982ux5ff5ux4e0eux516cux5f0f-4}

\subsubsection{为什么要把频率推广成复频率}\label{ux4e3aux4ec0ux4e48ux8981ux628aux9891ux7387ux63a8ux5e7fux6210ux590dux9891ux7387}

傅里叶变换要求
\(\int_{-\infty}^{+\infty}|f(t)|\,\mathrm{d}t<\infty\)，像
\(\mathrm{e}^{2t}u(t)\)、\(u(t)\)
这类信号不满足，因此必须先解决''积分不收敛''。

做法是给信号乘一个收敛因子 \(\mathrm{e}^{-\sigma t}\)，使乘积绝对可积：

\[
F(\sigma+\mathrm{j}\omega)=\int_{-\infty}^{+\infty}f(t)\mathrm{e}^{-\sigma t}
\mathrm{e}^{-\mathrm{j}\omega t}\,\mathrm{d}t
=\int_{-\infty}^{+\infty}f(t)\mathrm{e}^{-st}\,\mathrm{d}t ,
\qquad s=\sigma+\mathrm{j}\omega .
\]

\begin{itemize}
\tightlist
\item
  \(\sigma\) 控制包络的收敛与发散（响应快慢），\(\omega\)
  控制振荡；\(\sigma=0\) 时即退化为傅里叶变换。
\item
  用途：初始状态不为零的动态系统、指数增长信号、把微分方程化为代数方程。
\end{itemize}

\subsubsection{拉普拉斯变换的定义（单边与双边必须分清）}\label{ux62c9ux666eux62c9ux65afux53d8ux6362ux7684ux5b9aux4e49ux5355ux8fb9ux4e0eux53ccux8fb9ux5fc5ux987bux5206ux6e05}

双边拉普拉斯变换：

\[
F(s)=\mathcal{L}_{d}[f(t)]=\int_{-\infty}^{+\infty}f(t)\mathrm{e}^{-st}\,\mathrm{d}t ,
\] \[
f(t)=\frac{1}{2\pi \mathrm{j}}\int_{\sigma-\mathrm{j}\infty}^{\sigma+\mathrm{j}\infty}
F(s)\mathrm{e}^{st}\,\mathrm{d}s .
\]

单边拉普拉斯变换：

\[
F(s)=\mathcal{L}[f(t)]=\int_{0^{-}}^{+\infty}f(t)\mathrm{e}^{-st}\,\mathrm{d}t ,
\qquad
f(t)=\left[\frac{1}{2\pi \mathrm{j}}\int_{\sigma-\mathrm{j}\infty}^{\sigma+\mathrm{j}\infty}
F(s)\mathrm{e}^{st}\,\mathrm{d}s\right]u(t).
\]

\begin{longtable}[]{@{}
  >{\raggedright\arraybackslash}p{(\linewidth - 4\tabcolsep) * \real{0.3333}}
  >{\raggedright\arraybackslash}p{(\linewidth - 4\tabcolsep) * \real{0.3333}}
  >{\raggedright\arraybackslash}p{(\linewidth - 4\tabcolsep) * \real{0.3333}}@{}}
\toprule\noalign{}
\begin{minipage}[b]{\linewidth}\raggedright
项目
\end{minipage} & \begin{minipage}[b]{\linewidth}\raggedright
单边
\end{minipage} & \begin{minipage}[b]{\linewidth}\raggedright
双边
\end{minipage} \\
\midrule\noalign{}
\endhead
\bottomrule\noalign{}
\endlastfoot
积分下限 & \(0^{-}\) & \(-\infty\) \\
反映的信号区间 & 仅 \(t\ge 0\) 部分 & 全部 \(t\) \\
收敛域形状 & 半平面 \(\mathrm{Re}(s)>\sigma_{0}\) & 条带
\(\sigma_{1}<\mathrm{Re}(s)<\sigma_{2}\) \\
唯一性 & \(F(s)\) 唯一对应 \(t\ge 0\) 的波形 & 必须同时给 \(F(s)\)
与收敛域 \\
\end{longtable}

成立条件：积分在收敛域内绝对收敛。下限取 \(0^{-}\) 是为了把 \(t=0\)
处的冲激以及 \(t=0^{-}\)
的初始状态一并计入，这样微分性质里的初始项才与电路''动作发生前''的状态对齐。

\subsubsection{收敛域及其与信号类型的关系}\label{ux6536ux655bux57dfux53caux5176ux4e0eux4fe1ux53f7ux7c7bux578bux7684ux5173ux7cfb}

收敛域（ROC）是使
\(\int |f(t)\mathrm{e}^{-\sigma t}|\,\mathrm{d}t<\infty\) 的 \(s\)
的集合， 边界 \(\sigma_{0}\) 称收敛横坐标。对有理 \(F(s)\)，ROC
的边界一定穿过极点，且 ROC 内不含极点。

\begin{longtable}[]{@{}
  >{\raggedright\arraybackslash}p{(\linewidth - 4\tabcolsep) * \real{0.3333}}
  >{\raggedright\arraybackslash}p{(\linewidth - 4\tabcolsep) * \real{0.3333}}
  >{\raggedright\arraybackslash}p{(\linewidth - 4\tabcolsep) * \real{0.3333}}@{}}
\toprule\noalign{}
\begin{minipage}[b]{\linewidth}\raggedright
信号类型
\end{minipage} & \begin{minipage}[b]{\linewidth}\raggedright
ROC 形状
\end{minipage} & \begin{minipage}[b]{\linewidth}\raggedright
例
\end{minipage} \\
\midrule\noalign{}
\endhead
\bottomrule\noalign{}
\endlastfoot
右边（有始、因果） & \(\mathrm{Re}(s)>\sigma_{0}\) &
\(\mathrm{e}^{-2t}u(t)\)：\(\mathrm{Re}(s)>-2\) \\
左边（有终） & \(\mathrm{Re}(s)<\sigma_{0}\) &
\(-\mathrm{e}^{-2t}u(-t)\)：\(\mathrm{Re}(s)<-2\) \\
双边 & \(\sigma_{1}<\mathrm{Re}(s)<\sigma_{2}\) & 两项相加，条带形 \\
时限且绝对可积 & 整个 \(s\) 平面 & 单个矩形脉冲 \\
单边变换 & 必为 \(\mathrm{Re}(s)>\sigma_{0}\) & 与 \(t\ge 0\)
部分有关 \\
\end{longtable}

三条常用推论：

\begin{enumerate}
\def\labelenumi{\arabic{enumi}.}
\tightlist
\item
  单边变换中 \(\sigma_{0}\) 取 \(F(s)\) 各极点实部的最大值（右边信号的
  ROC 在最右极点之右）。
\item
  若 ROC 包含虚轴，则傅里叶变换存在且
  \(F(\mathrm{j}\omega)=F(s)\big|_{s=\mathrm{j}\omega}\)。
\item
  有理 \(F(s)\) 与不同 ROC
  搭配会给出不同原信号，这是双边反变换''一表多值''的根源。
\end{enumerate}

\subsubsection{常用变换对}\label{ux5e38ux7528ux53d8ux6362ux5bf9}

下表均为单边变换，\(u(t)\) 表示 \(t\ge 0\) 部分。

\begin{longtable}[]{@{}
  >{\raggedright\arraybackslash}p{(\linewidth - 4\tabcolsep) * \real{0.3333}}
  >{\raggedright\arraybackslash}p{(\linewidth - 4\tabcolsep) * \real{0.3333}}
  >{\raggedright\arraybackslash}p{(\linewidth - 4\tabcolsep) * \real{0.3333}}@{}}
\toprule\noalign{}
\begin{minipage}[b]{\linewidth}\raggedright
\(f(t)\)
\end{minipage} & \begin{minipage}[b]{\linewidth}\raggedright
\(F(s)\)
\end{minipage} & \begin{minipage}[b]{\linewidth}\raggedright
收敛域
\end{minipage} \\
\midrule\noalign{}
\endhead
\bottomrule\noalign{}
\endlastfoot
\(\delta(t)\) & \(1\) & 整个 \(s\) 平面 \\
\(\delta'(t)\) & \(s\) & 整个 \(s\) 平面 \\
\(u(t)\) & \(\dfrac{1}{s}\) & \(\mathrm{Re}(s)>0\) \\
\(t^{n}u(t)\) & \(\dfrac{n!}{s^{n+1}}\) & \(\mathrm{Re}(s)>0\) \\
\(\mathrm{e}^{-at}u(t)\) & \(\dfrac{1}{s+a}\) & \(\mathrm{Re}(s)>-a\) \\
\(t\,\mathrm{e}^{-at}u(t)\) & \(\dfrac{1}{(s+a)^{2}}\) &
\(\mathrm{Re}(s)>-a\) \\
\(\sin\omega_{0}t\,u(t)\) & \(\dfrac{\omega_{0}}{s^{2}+\omega_{0}^{2}}\)
& \(\mathrm{Re}(s)>0\) \\
\(\cos\omega_{0}t\,u(t)\) & \(\dfrac{s}{s^{2}+\omega_{0}^{2}}\) &
\(\mathrm{Re}(s)>0\) \\
\(\mathrm{e}^{-at}\sin\omega_{0}t\,u(t)\) &
\(\dfrac{\omega_{0}}{(s+a)^{2}+\omega_{0}^{2}}\) &
\(\mathrm{Re}(s)>-a\) \\
\(\mathrm{e}^{-at}\cos\omega_{0}t\,u(t)\) &
\(\dfrac{s+a}{(s+a)^{2}+\omega_{0}^{2}}\) & \(\mathrm{Re}(s)>-a\) \\
\end{longtable}

记法提示：把 \(s\) 换成 \(s+a\) 就是''乘
\(\mathrm{e}^{-at}\)``（频移）， 把分母 \(s^{2}+\omega_{0}^{2}\)
因式分解为
\((s-\mathrm{j}\omega_{0})(s+\mathrm{j}\omega_{0})\)，就能与极点
\(\pm \mathrm{j}\omega_{0}\) 对应。

\subsubsection{拉普拉斯变换的性质}\label{ux62c9ux666eux62c9ux65afux53d8ux6362ux7684ux6027ux8d28}

\begin{longtable}[]{@{}
  >{\raggedright\arraybackslash}p{(\linewidth - 6\tabcolsep) * \real{0.2500}}
  >{\raggedright\arraybackslash}p{(\linewidth - 6\tabcolsep) * \real{0.2500}}
  >{\raggedright\arraybackslash}p{(\linewidth - 6\tabcolsep) * \real{0.2500}}
  >{\raggedright\arraybackslash}p{(\linewidth - 6\tabcolsep) * \real{0.2500}}@{}}
\toprule\noalign{}
\begin{minipage}[b]{\linewidth}\raggedright
性质
\end{minipage} & \begin{minipage}[b]{\linewidth}\raggedright
时域
\end{minipage} & \begin{minipage}[b]{\linewidth}\raggedright
复频域
\end{minipage} & \begin{minipage}[b]{\linewidth}\raggedright
条件或收敛域
\end{minipage} \\
\midrule\noalign{}
\endhead
\bottomrule\noalign{}
\endlastfoot
线性 & \(a_{1}f_{1}+a_{2}f_{2}\) & \(a_{1}F_{1}+a_{2}F_{2}\) & 至少为两
ROC 的交 \\
尺度变换 & \(f(at),\ a>0\) & \(\frac{1}{a}F\!\left(\frac{s}{a}\right)\)
& ROC 按 \(a\) 缩放 \\
时移 & \(f(t-t_{0})u(t-t_{0})\) & \(\mathrm{e}^{-st_{0}}F(s)\) &
单边；ROC 不变 \\
频移 & \(\mathrm{e}^{s_{0}t}f(t)\) & \(F(s-s_{0})\) & ROC 平移
\(\mathrm{Re}(s_{0})\) \\
时域微分 & \(f'(t)\) & \(sF(s)-f(0^{-})\) & 单边，\(f\) 可导 \\
时域积分 & \(\int_{0^{-}}^{t}f(\tau)\,\mathrm{d}\tau\) &
\(\dfrac{F(s)}{s}\) & 单边 \\
频域微分 & \(t f(t)\) & \(-\dfrac{\mathrm{d}F(s)}{\mathrm{d}s}\) & ROC
不变 \\
卷积定理 & \(f_{1}(t)*f_{2}(t)\) & \(F_{1}(s)F_{2}(s)\) & 至少为两 ROC
的交 \\
\end{longtable}

需要写全的两条：

\begin{enumerate}
\def\labelenumi{\arabic{enumi}.}
\tightlist
\item
  高阶微分（\(0^{-}\) 系统，初始项不能丢）
\end{enumerate}

\[
\mathcal{L}\left[f^{(n)}(t)\right]
=s^{n}F(s)-\sum_{k=0}^{n-1}s^{n-1-k}f^{(k)}(0^{-}),
\]

二阶即 \(s^{2}F(s)-sf(0^{-})-f'(0^{-})\)，三阶即
\(s^{3}F(s)-s^{2}f(0^{-})-sf'(0^{-})-f''(0^{-})\)。

\begin{enumerate}
\def\labelenumi{\arabic{enumi}.}
\setcounter{enumi}{1}
\tightlist
\item
  积分性质的两种下限（不可混为一式）
\end{enumerate}

\[
\mathcal{L}\left[\int_{0^{-}}^{t}f(\tau)\,\mathrm{d}\tau\right]=\frac{F(s)}{s},
\qquad
\mathcal{L}\left[\int_{-\infty}^{t}f(\tau)\,\mathrm{d}\tau\right]
=\frac{F(s)}{s}+\frac{1}{s}\int_{-\infty}^{0^{-}}f(\tau)\,\mathrm{d}\tau .
\]

单边拉普拉斯变换的积分下限默认取 \(0^{-}\)，此时才有 \(\dfrac{F(s)}{s}\)
这个简洁结果； 下限放到 \(-\infty\) 时必须补上 \(t<0\) 段的累积量
\(\dfrac{1}{s}\int_{-\infty}^{0^{-}}f(\tau)\,\mathrm{d}\tau\)。 只有
\(f\) 为因果信号（或 \(f\) 在 \(t<0\)
段的积分为零）时该常数项为零，两式才相同。

初值定理（成立条件：\(f(t)\) 与 \(f'(t)\) 存在且拉普拉斯变换存在）：

\[
f(0^{+})=\lim_{t\to 0^{+}}f(t)=\lim_{s\to \infty}sF(s).
\]

若 \(f(t)\) 或其导数含冲激（即 \(F(s)\)
为假分式、\(\lim_{s\to\infty}F(s)\) 不存在）， 先用长除把 \(F(s)\)
分成多项式与严格真分式 \(F_{p}(s)\) 之和，再取
\(f(0^{+})=\lim_{s\to\infty}sF_{p}(s)\)。

终值定理（成立条件：\(F(s)\) 的极点全在左半平面，\(s=0\)
处至多一阶极点， 即 \(sF(s)\) 的极点全在左半平面）：

\[
f(+\infty)=\lim_{t\to +\infty}f(t)=\lim_{s\to 0}sF(s).
\]

终值定理是求''阶跃输入下的稳态输出''最快的路子：输入 \(u(t)\) 时
\(Y(s)=H(s)/s\)， 故 \(y(+\infty)=H(0)\)。

\subsubsection{单边拉普拉斯反变换：部分分式展开法}\label{ux5355ux8fb9ux62c9ux666eux62c9ux65afux53cdux53d8ux6362ux90e8ux5206ux5206ux5f0fux5c55ux5f00ux6cd5}

思路：把有理 \(F(s)\) 拆成若干个 \(\dfrac{K_{i}}{s-s_{i}}\) 之和，
每一项对应一个基本时间模式
\(K_{i}\mathrm{e}^{s_{i}t}u(t)\)，再线性叠加。 设
\(F(s)=\dfrac{N(s)}{D(s)}\)。

\textbf{情形 1：\(m<n\) 且 \(D(s)=0\) 只有单实根}

\[
F(s)=\sum_{i=1}^{n}\frac{K_{i}}{s-s_{i}},
\qquad
K_{i}=(s-s_{i})F(s)\Big|_{s=s_{i}}=\frac{N(s_{i})}{D'(s_{i})},
\] \[
f(t)=\sum_{i=1}^{n}K_{i}\mathrm{e}^{s_{i}t}u(t).
\]

\textbf{情形 2：共轭复根
\(s_{1,2}=\sigma_{1}\pm \mathrm{j}\omega_{0}\)}（\(D(s)\)
实系数，复根必成对）

系数共轭 \(K_{2}=K_{1}^{*}\)，令
\(K_{1}=|K_{1}|\mathrm{e}^{\mathrm{j}\phi_{1}}\)，两项合并为实信号：

\[
K_{1}\mathrm{e}^{s_{1}t}+K_{2}\mathrm{e}^{s_{2}t}
=2|K_{1}|\mathrm{e}^{\sigma_{1}t}\cos(\omega_{0}t+\phi_{1}).
\]

实部 \(\sigma_{1}\) 控制包络（衰减或增长），虚部 \(\omega_{0}\)
控制振荡频率。 这里的 \(\omega_{0}\)
就是这对共轭根的虚部，与``常用变换对''一节中正弦信号的角频率是同一个量，
不另设符号。

\textbf{情形 3：\(r\) 阶重根 \(p\)}

\[
F(s)=\sum_{k=1}^{r}\frac{K_{1k}}{(s-p)^{k}},
\qquad
K_{1k}=\frac{1}{(r-k)!}\frac{\mathrm{d}^{r-k}}{\mathrm{d}s^{r-k}}
\left[(s-p)^{r}F(s)\right]_{s=p},
\] 对应时间项 \(\dfrac{K_{1k}t^{r-k}}{(r-k)!}\mathrm{e}^{pt}\)。

\textbf{情形 4：假分式 \(m\ge n\)}

先长除，化为多项式与真分式之和：

\[
F(s)=\frac{N(s)}{D(s)}=M(s)+\frac{N_{1}(s)}{D(s)},
\]

再对真分式部分用情形 1～3。多项式部分对应冲激及其导数：

\[
\mathcal{L}^{-1}\left[s^{n}\right]=\delta^{(n)}(t),
\qquad \mathcal{L}^{-1}[1]=\delta(t).
\]

\subsubsection{用拉普拉斯变换求解微分方程（零输入加零状态）}\label{ux7528ux62c9ux666eux62c9ux65afux53d8ux6362ux6c42ux89e3ux5faeux5206ux65b9ux7a0bux96f6ux8f93ux5165ux52a0ux96f6ux72b6ux6001}

设系统方程与初始状态

\[
\sum_{i=0}^{n}a_{i}y^{(i)}(t)=\sum_{j=0}^{m}b_{j}f^{(j)}(t),
\]

并设初始状态 \(y(0^{-}),y'(0^{-}),\dots,y^{(n-1)}(0^{-})\) 全部已知。

两边取单边变换，记 \(A(s)=\sum a_{i}s^{i}\)、\(B(s)=\sum b_{j}s^{j}\)：

\[
A(s)Y(s)-I(s)=B(s)F(s),
\] \[
I(s)=\sum_{i=1}^{n}a_{i}\sum_{k=0}^{i-1}s^{i-1-k}y^{(k)}(0^{-}).
\]

于是

\[
Y(s)=\underbrace{\frac{I(s)}{A(s)}}_{Y_{zi}(s)}
+\underbrace{\frac{B(s)}{A(s)}F(s)}_{Y_{zs}(s)},
\qquad
H(s)=\frac{B(s)}{A(s)},\quad Y_{zs}(s)=H(s)F(s).
\]

四步流程：① 逐项取变换得代数方程；② 解出 \(Y(s)\)； ③
按极点做部分分式（分母 \(A(s)\) 的根给零输入项，\(F(s)\)
的根给受迫项）； ④ 反变换得 \(y(t)\)。

响应成分的划分（按极点来源）：

\begin{itemize}
\tightlist
\item
  由 \(A(s)=0\) 的根（即 \(H(s)\)
  的极点、系统的特征根）产生的项，是自然响应 \(y_{n}(t)\)；
\item
  由 \(F(s)\) 的极点产生的项，是受迫响应 \(y_{f}(t)\)；
\item
  零输入响应只含自然响应成分，零状态响应同时含自然与受迫两种成分。
\end{itemize}

\subsubsection{\texorpdfstring{电路的 \(s\)
域模型}{电路的 s 域模型}}\label{ux7535ux8defux7684-s-ux57dfux6a21ux578b}

把每个元件换成''运算阻抗 +
初始状态等效源''，其余步骤与电阻电路完全相同（KCL、KVL、
分压分流、节点法）。

\begin{longtable}[]{@{}
  >{\raggedright\arraybackslash}p{(\linewidth - 6\tabcolsep) * \real{0.2500}}
  >{\raggedright\arraybackslash}p{(\linewidth - 6\tabcolsep) * \real{0.2500}}
  >{\raggedright\arraybackslash}p{(\linewidth - 6\tabcolsep) * \real{0.2500}}
  >{\raggedright\arraybackslash}p{(\linewidth - 6\tabcolsep) * \real{0.2500}}@{}}
\toprule\noalign{}
\begin{minipage}[b]{\linewidth}\raggedright
元件
\end{minipage} & \begin{minipage}[b]{\linewidth}\raggedright
时域关系
\end{minipage} & \begin{minipage}[b]{\linewidth}\raggedright
运算阻抗
\end{minipage} & \begin{minipage}[b]{\linewidth}\raggedright
串联等效电压源
\end{minipage} \\
\midrule\noalign{}
\endhead
\bottomrule\noalign{}
\endlastfoot
\(R\) & \(u=Ri\) & \(R\) & 无 \\
\(L\) & \(u=L\dfrac{\mathrm{d}i}{\mathrm{d}t}\) & \(sL\) &
\(Li_{L}(0^{-})\) \\
\(C\) & \(i=C\dfrac{\mathrm{d}u}{\mathrm{d}t}\) & \(\dfrac{1}{sC}\) &
\(\dfrac{u_{C}(0^{-})}{s}\) \\
\end{longtable}

相应方程为

\[
U_{L}(s)=sLI_{L}(s)-Li_{L}(0^{-}),
\qquad
U_{C}(s)=\frac{I_{C}(s)}{sC}+\frac{u_{C}(0^{-})}{s}.
\]

并联（诺顿）形式：\(L\) 等效为 \(\dfrac{1}{sL}\) 并联电流源
\(\dfrac{i_{L}(0^{-})}{s}\)； \(C\) 等效为 \(sC\) 并联电流源
\(Cu_{C}(0^{-})\)。

求系统函数时把全部初始状态置零（电压源短路、电流源开路），只留阻抗网络，
这样得到的 \(H(s)\) 才只描述系统本身。

\subsubsection{线性系统的模拟}\label{ux7ebfux6027ux7cfbux7edfux7684ux6a21ux62df}

同一系统的四种表示：微分方程、系统函数、框图或流图、状态变量。框图的基本运算单元是
加法器、数乘器与积分器 \(\dfrac{1}{s}\)（不用微分器，避免放大噪声）。

直接型（\(n\) 阶）画法，设

\[
H(s)=\frac{b_{m}s^{m}+\cdots+b_{1}s+b_{0}}{s^{n}+a_{n-1}s^{n-1}+\cdots+a_{1}s+a_{0}} :
\]

\begin{enumerate}
\def\labelenumi{\arabic{enumi}.}
\tightlist
\item
  取 \(n\) 个积分器串成链，从最高阶开始，令
  \(q^{(n)}=x-a_{n-1}q^{(n-1)}-\cdots-a_{1}q'-a_{0}q\)， 即从
  \(q^{(n-1)},\dots,q\) 各引一条增益为 \(-a_{n-1},\dots,-a_{0}\)
  的反馈支路到最前端加法器；
\item
  前向加和：\(y=b_{m}q^{(m)}+\cdots+b_{1}q'+b_{0}q\)，
  即从各级积分器输出再引增益 \(b_{j}\) 的前馈支路到末端加法器。
\end{enumerate}

组合规律：子系统串联时 \(H(s)=\prod H_{i}(s)\)，并联时
\(H(s)=\sum H_{i}(s)\)。 串联对应代数乘法，并联对应代数加法。

\subsection{典型题型与解法}\label{ux5178ux578bux9898ux578bux4e0eux89e3ux6cd5-4}

\subsubsection{\texorpdfstring{题型 5-1 由信号表达式标出 \(s\)
平面上的复频率}{题型 5-1 由信号表达式标出 s 平面上的复频率}}\label{ux9898ux578b-5-1-ux7531ux4fe1ux53f7ux8868ux8fbeux5f0fux6807ux51fa-s-ux5e73ux9762ux4e0aux7684ux590dux9891ux7387}

\textbf{题型特征}：给出 \(f(t)\) 的解析式（实指数、\(t^{n}\)
乘指数、正弦余弦），要求写出它在 \(s\)
平面上对应的复频率（极点位置），或反过来由极点写波形。

\textbf{解题步骤}： 1. 把 \(f(t)\) 化成
\(\sum A_{i}t^{n_{i}}\mathrm{e}^{s_{i}t}\) 的形式，\(\sin\)、\(\cos\)
先用欧拉公式拆开。 2. 每个指数项的指数 \(s_{i}\) 就是复频率；因子
\(t^{n}\) 表示该复频率是 \(n+1\) 重。 3. 在 \(s\)
平面标点：实指数落在实轴，振荡成对落在关于实轴对称的位置。

\textbf{例题}（自编例）：求 \(f(t)=\mathrm{e}^{-3t}\cos 4t\,u(t)\)
的复频率与收敛域。

解：由欧拉公式

\[
\mathrm{e}^{-3t}\cos 4t=\frac{1}{2}\mathrm{e}^{(-3+\mathrm{j}4)t}
+\frac{1}{2}\mathrm{e}^{(-3-\mathrm{j}4)t},
\]

故复频率为一对共轭值 \(s_{1,2}=-3\pm \mathrm{j}4\)， 在 \(s\)
平面上是左半平面内关于实轴对称的一对点，实部 \(-3\) 决定衰减快慢， 虚部
\(4\ \mathrm{rad/s}\) 决定振荡频率。收敛域 \(\mathrm{Re}(s)>-3\)。

\textbf{同类题}： - 【附加作业6 第1题(1)(2)】\(\mathrm{e}^{2t}\)
是单个实指数，复频率 \(s_{1}=2\)； \(t\mathrm{e}^{-t}\) 含一阶因子
\(t\)，故 \(s_{1,2}=-1\) 为二重复频率。 - 【附加作业6
第1题(3)(4)】\(\cos 2t\) 拆成
\(\frac{1}{2}\mathrm{e}^{\pm \mathrm{j}2t}\)， 得
\(s_{1,2}=\pm \mathrm{j}2\)（虚轴上的一对，等幅振荡）；
\(\mathrm{e}^{-t}\sin(-5t)\) 中负号只改变幅度符号，复频率仍为
\(s_{1,2}=-1\pm \mathrm{j}5\)。

\subsubsection{题型 5-2
真分式的部分分式展开与反变换}\label{ux9898ux578b-5-2-ux771fux5206ux5f0fux7684ux90e8ux5206ux5206ux5f0fux5c55ux5f00ux4e0eux53cdux53d8ux6362}

\textbf{题型特征}：给 \(F(s)\)（分子阶次低于分母），求
\(f(t)\)；根为单实根，或可先约分化简。

\textbf{解题步骤}： 1. 检查是否真分式（\(m<n\)），不是就先长除（见题型
5-3）。 2. 求 \(D(s)=0\) 的根，判断是单实根、共轭复根还是重根。 3.
单实根用 \(K_{i}=(s-s_{i})F(s)|_{s=s_{i}}\) 或 \(N(s_{i})/D'(s_{i})\)
求系数。 4. 逐项反变换并乘 \(u(t)\)；用初值定理
\(\lim_{s\to\infty}sF(s)\) 回代校验。

\textbf{例题}（【附加作业6 第2题(1)】）：求
\(F(s)=\dfrac{s}{(s+1)(s+4)}\) 的原函数。

解：\(m=1<n=2\)，为真分式，\(D(s)=0\) 的根为
\(s_{1}=-1\)、\(s_{2}=-4\)，均为单实根。设

\[
F(s)=\frac{K_{1}}{s+1}+\frac{K_{2}}{s+4},
\]

用掩模法（各项乘以对应因子后取 \(s=s_{i}\)）：

\[
K_{1}=(s+1)F(s)\Big|_{s=-1}=\frac{-1}{-1+4}=-\frac{1}{3},
\qquad
K_{2}=(s+4)F(s)\Big|_{s=-4}=\frac{-4}{-4+1}=\frac{4}{3}.
\]

于是

\[
F(s)=-\frac{1}{3}\cdot\frac{1}{s+1}+\frac{4}{3}\cdot\frac{1}{s+4},
\]

查变换对
\(\mathcal{L}^{-1}\left[\dfrac{1}{s+a}\right]=\mathrm{e}^{-at}u(t)\)，得

\[
f(t)=\left(-\frac{1}{3}\mathrm{e}^{-t}+\frac{4}{3}\mathrm{e}^{-4t}\right)u(t).
\]

校验：\(f(0^{+})=-\frac{1}{3}+\frac{4}{3}=1\)，而
\(\lim_{s\to\infty}sF(s)=1\)，一致。

\textbf{同类题}： - 【附加作业6 第2题(2)】\(\dfrac{s+1}{s^{2}-1}\)
先因式分解约分：\(s^{2}-1=(s-1)(s+1)\)， 极零相消后得
\(\dfrac{1}{s-1}\)，收敛域 \(\mathrm{Re}(s)>1\)，故
\(f(t)=\mathrm{e}^{t}u(t)\)； 注意相消后只剩右半平面极点
\(s=1\)，这是增长型（不稳定）信号。

\subsubsection{题型 5-3
假分式的反变换（多项式项对应冲激及其导数）}\label{ux9898ux578b-5-3-ux5047ux5206ux5f0fux7684ux53cdux53d8ux6362ux591aux9879ux5f0fux9879ux5bf9ux5e94ux51b2ux6fc0ux53caux5176ux5bfcux6570}

\textbf{题型特征}：分子阶次不低于分母（\(m\ge n\)），必须先长除；答案是冲激及其导数加指数项。

\textbf{解题步骤}： 1.
长除：\(F(s)=M(s)+\dfrac{N_{1}(s)}{D(s)}\)，\(M(s)\) 是 \(s\) 的多项式。
2. \(M(s)\) 逐项用 \(\mathcal{L}^{-1}[s^{n}]=\delta^{(n)}(t)\) 反变换。
3. 真分式部分按题型 5-2 展开并反变换。 4. 结果中一定含 \(\delta(t)\) 或
\(\delta'(t)\) 这类奇异项，不能漏。

\textbf{例题}（【附加作业6 第2题(3)】）：求
\(F(s)=\dfrac{s^{3}+6s^{2}+6s}{s^{2}+6s+8}\) 的原函数。

解：\(m=3>n=2\)，为假分式。先做多项式除法：

\[
s^{3}+6s^{2}+6s=\left(s^{2}+6s+8\right)s+\left(-2s\right),
\]

故

\[
F(s)=s+\frac{-2s}{s^{2}+6s+8}=s+\frac{-2s}{(s+2)(s+4)}.
\]

对真分式部分设
\(\dfrac{-2s}{(s+2)(s+4)}=\dfrac{K_{1}}{s+2}+\dfrac{K_{2}}{s+4}\)，则

\[
K_{1}=\frac{-2(-2)}{-2+4}=2,
\qquad
K_{2}=\frac{-2(-4)}{-4+2}=-4 .
\]

于是

\[
F(s)=s+\frac{2}{s+2}-\frac{4}{s+4}.
\]

由 \(\mathcal{L}^{-1}[s]=\delta'(t)\) 与查表得

\[
f(t)=\delta'(t)+\left(2\mathrm{e}^{-2t}-4\mathrm{e}^{-4t}\right)u(t).
\]

\textbf{同类题}： - 【附加作业6 第2题(1)】若把 \(\dfrac{s}{(s+1)(s+4)}\)
误当假分式去长除，会白算一轮；
先比阶次再决定要不要长除，是本类题的第一动作。

\subsubsection{题型 5-4
用拉普拉斯变换求全响应并区分自然响应与受迫响应}\label{ux9898ux578b-5-4-ux7528ux62c9ux666eux62c9ux65afux53d8ux6362ux6c42ux5168ux54cdux5e94ux5e76ux533aux5206ux81eaux7136ux54cdux5e94ux4e0eux53d7ux8febux54cdux5e94}

\textbf{题型特征}：给二阶（或更高阶）微分方程、输入表达式与 \(0^{-}\)
初始状态，
要求全响应，并把响应拆成自然响应与受迫响应，或拆成零输入与零状态。

\textbf{解题步骤}： 1.
对等式两边逐项取单边变换，套高阶微分性质，把初始状态写成常数项移到右边。
2. 解出 \(Y(s)\)，写成''输入项除特征多项式 + 初始项除特征多项式''两半。
3. 分别对两半做部分分式展开（分母 \(A(s)\) 的根与 \(F(s)\)
的根一起展开）。 4. 反变换得全响应；按极点来源标注自然项与受迫项。 5. 用
\(y(0^{+})=y(0^{-})\) 与 \(\lim_{s\to\infty}sY(s)\) 校验数值。

\textbf{例题}（【附加作业6 第3题】）：已知因果 LTI 系统
\(y''(t)+3y'(t)+2y(t)=3e(t)\)，\(e(t)=\mathrm{e}^{-4t}u(t)\)，
\(y(0^{-})=1\)，\(y'(0^{-})=2\)。求 1）用拉普拉斯变换法求全响应；
2）分析响应成分，写出自然响应与受迫响应。

解 1）两边取单边变换，记
\(Y(s)=\mathcal{L}[y(t)]\)、\(E(s)=\mathcal{L}[e(t)]\)， 由二阶微分性质
\(y''\to s^{2}Y(s)-sy(0^{-})-y'(0^{-})\)、\(y'\to sY(s)-y(0^{-})\)：

\[
s^{2}Y(s)-s\cdot 1-2+3\left[sY(s)-1\right]+2Y(s)=3E(s).
\]

整理得

\[
\left(s^{2}+3s+2\right)Y(s)=\left(s+5\right)+3E(s),
\qquad E(s)=\frac{1}{s+4},
\]

即

\[
Y(s)=\frac{3}{(s+1)(s+2)}\cdot\frac{1}{s+4}
+\frac{s+5}{(s+1)(s+2)} .
\]

第一项为零状态响应，第二项为零输入响应。分别展开：

零状态部分，三个单实根 \(s=-1,-2,-4\)：

\[
\frac{3}{(s+1)(s+2)(s+4)}
=\frac{1}{s+1}-\frac{1.5}{s+2}+\frac{0.5}{s+4},
\]

其中 \(K_{-1}=\frac{3}{(1)(3)}=1\)，\(K_{-2}=\frac{3}{(-1)(2)}=-1.5\)，
\(K_{-4}=\frac{3}{(-3)(-2)}=0.5\)。

零输入部分：

\[
\frac{s+5}{(s+1)(s+2)}=\frac{4}{s+1}-\frac{3}{s+2},
\]

其中 \(K_{-1}=\frac{4}{1}=4\)，\(K_{-2}=\frac{3}{-1}=-3\)。

合并同类项（同一极点的系数相加）：

\[
Y(s)=\frac{5}{s+1}-\frac{4.5}{s+2}+\frac{0.5}{s+4},
\]

反变换得全响应

\[
y(t)=\left(5\mathrm{e}^{-t}-4.5\mathrm{e}^{-2t}+0.5\mathrm{e}^{-4t}\right)u(t).
\]

校验：\(y(0^{+})=5-4.5+0.5=1=y(0^{-})\)；且
\(\lim_{s\to\infty}sY(s)=\lim s\cdot\frac{5}{s}=5\)， 与
\(y'(0^{+})=-5+9-2=2\) 不矛盾（初值定理给的是 \(y(0^{+})\) 而非导数）。

解 2）按极点来源分类：\(s=-1\)、\(s=-2\) 是特征方程 \(s^{2}+3s+2=0\)
的根，属系统自身极点，对应自然响应； \(s=-4\) 是输入
\(e(t)=\mathrm{e}^{-4t}u(t)\) 带来的极点，对应受迫响应。故

自然响应：\(\left(5\mathrm{e}^{-t}-4.5\mathrm{e}^{-2t}\right)u(t)\)；
受迫响应：\(0.5\mathrm{e}^{-4t}u(t)\)。

注意零输入响应 \(\left(4\mathrm{e}^{-t}-3\mathrm{e}^{-2t}\right)u(t)\)
只含自然成分， 零状态响应
\(\left(\mathrm{e}^{-t}-1.5\mathrm{e}^{-2t}+0.5\mathrm{e}^{-4t}\right)u(t)\)
同时含自然与受迫成分。

\textbf{同类题}： - 【附加作业6 第9题】只需求 \(G(s)\)
而非时域波形，思路是''先由微分方程读 \(H(s)\)，
再用积分性质与频移性质按部就班变换''，见题型 6-6。

\subsubsection{题型 5-5
由系统函数画直接模拟框图}\label{ux9898ux578b-5-5-ux7531ux7cfbux7edfux51fdux6570ux753bux76f4ux63a5ux6a21ux62dfux6846ux56fe}

\textbf{题型特征}：给微分方程或
\(H(s)\)，要求画出只用加法器、数乘器、积分器的直接模拟框图。

\textbf{解题步骤}： 1. 由方程写出
\(H(s)=B(s)/A(s)\)（或直接把算子方程看成 \(A(p)y=B(p)x\)）。 2.
定积分器个数：等于分母最高次幂 \(n\)；积分器链的最高阶量写成
\(q^{(n)}=x-a_{n-1}q^{(n-1)}-\cdots-a_{0}q\)，从加法器输出接第一个积分器。
3. 反馈支路系数取 \(-a_{i}\)，从第 \(i\) 个积分器之后引出。 4.
前向加和：从各级输出引 \(b_{j}\) 到末端加法器得 \(y\)。

\textbf{例题}（自编例）：画 \(H(s)=\dfrac{1}{s^{2}+3s+2}\)
的直接模拟框图。

解：分母二阶，用 2 个积分器。取第二个积分器的输入为 \(q''\)，则

\[
q''=x-3q'-2q,
\qquad y=q .
\]

框图结构（文字描述）：输入 \(x\)
进入前端加法器，加法器输出接第一个积分器得 \(q'\)， 再经第二个积分器得
\(q\)；从 \(q'\) 引增益 \(-3\) 的反馈支路、从 \(q\) 引增益 \(-2\)
的反馈支路， 都回到前端加法器；输出 \(y\) 直接取自 \(q\)（前馈系数
\(b_{1}=b_{2}=0\)，\(b_{0}=1\)）。

\textbf{同类题}： - 【附加作业6 第4题(1)】\((p^{3}+3p+2)y(t)=x(t)\)，即
\(H(s)=\dfrac{1}{s^{3}+3s+2}\)： 3 个积分器串成链，最高阶 \(y'''\)
由加法器给出，需 \(y'''=x-3y'-2y\)， 故从 \(y'\) 引增益 \(-3\)、从 \(y\)
引增益 \(-2\) 的反馈回加法器，中间一项 \(y''\) 无反馈支路。 -
【附加作业6 第4题(2)】\((p^{3}+3p^{2}+3p+2)y(t)=(p^{2}+2p)x(t)\)，即
\(H(s)=\dfrac{s^{2}+2s}{s^{3}+3s^{2}+3s+2}\)：内部变量取 \(q\)，反馈三条
\(-3\)（来自 \(q''\)）、\(-3\)（来自 \(q'\)）、\(-2\)（来自 \(q\)），
前向再引 \(1\)（来自 \(q''\)）、\(2\)（来自 \(q'\)）到输出加法器得
\(y=q''+2q'\)。

\subsubsection{\texorpdfstring{题型 5-6 电路的 \(s\)
域模型与系统函数}{题型 5-6 电路的 s 域模型与系统函数}}\label{ux9898ux578b-5-6-ux7535ux8defux7684-s-ux57dfux6a21ux578bux4e0eux7cfbux7edfux51fdux6570}

\textbf{题型特征}：给含 \(R\)、\(L\)、\(C\)
与初始状态的电路，要求画运算电路、求 \(H(s)\) 或某响应。

\textbf{解题步骤}： 1. 逐元件换成运算阻抗，并把
\(i_{L}(0^{-})\)、\(u_{C}(0^{-})\) 画成串联电压源。 2. 求 \(H(s)\)
时把初始状态等效源置零，只留阻抗网络。 3. 用分压、分流或节点法列 \(s\)
域代数方程，解出 \(U_{2}(s)/U_{1}(s)\)。 4. 需要时域响应时再对
\(H(s)F(s)\) 做部分分式反变换。

\textbf{例题}（自编例）：\(R\) 与 \(C\) 串联，以 \(u_{1}(t)\)
为输入、\(u_{C}(t)\) 为输出， 电容初始电压为零，求 \(H(s)\) 与
\(h(t)\)。

解：\(s\) 域分压（电容阻抗 \(\dfrac{1}{sC}\)，电阻 \(R\)）：

\[
H(s)=\frac{U_{C}(s)}{U_{1}(s)}
=\frac{\frac{1}{sC}}{R+\frac{1}{sC}}
=\frac{1}{1+sRC}=\frac{1}{RC}\cdot\frac{1}{s+\frac{1}{RC}} .
\]

反变换得 \(h(t)=\dfrac{1}{RC}\mathrm{e}^{-\frac{t}{RC}}u(t)\)， 极点
\(s=-\dfrac{1}{RC}\) 在左半平面，时间常数 \(\tau=RC\)
同时决定上升速度与拖尾长度。

\textbf{同类题}： - 【附加作业6 第5题】极零图反求 \(H(s)\)
后写微分方程，属于第 6 章题型 6-1； 其中
\(H(0)=\int_{0}^{\infty}h(t)\,\mathrm{d}t\) 的用法与终值定理同源。

\subsection{易错点}\label{ux6613ux9519ux70b9-4}

\begin{enumerate}
\def\labelenumi{\arabic{enumi}.}
\tightlist
\item
  \textbf{只写 \(F(s)\) 不写收敛域} → 必须同时写
  \(\mathrm{Re}(s)>\sigma_{0}\) → 双边变换中 \(F(s)\)
  与收敛域共同才能唯一确定原信号。
\item
  \textbf{把单边变换当双边用} → 单边只反映 \(t\ge 0\) 部分 → 单边下限是
  \(0^{-}\)， 含 \(0^{-}\) 初始状态信息，双边没有这层含义。
\item
  \textbf{时移性质漏掉 \(u(t-t_{0})\)} → 写成
  \(\mathcal{L}[f(t-t_{0})u(t-t_{0})]=\mathrm{e}^{-st_{0}}F(s)\) →
  单边变换要求 \(t<t_{0}\) 段确实为零，否则要补 \(t<0\) 部分的贡献。
\item
  \textbf{微分性质丢初始项} → \(y''\to s^{2}Y(s)-sy(0^{-})-y'(0^{-})\) →
  初始项就是''全响应一次算出''的来源，丢了等于把系统当零状态。
\item
  \textbf{对假分式直接做部分分式} → 先长除再展开 → 多项式部分对应
  \(\delta(t)\) 及其导数（第2题(3) 漏 \(\delta'(t)\) 是典型失分）。
\item
  \textbf{共轭复根写成复指数} → 合并成
  \(2|K_{1}|\mathrm{e}^{\sigma_{1}t}\cos(\omega_{0}t+\phi_{1})\) →
  实系数有理式的复根必共轭成对，实信号里不能留 \(\mathrm{j}\)。
\item
  \textbf{不查条件就用终值定理} → 先看 \(sF(s)\) 的极点是否全在左半平面
  → 虚轴上有极点时 \(f(+\infty)\) 不存在，算出来的是假值。
\item
  \textbf{积分性质上下限混用} → \(\int_{0^{-}}^{t}\) 与
  \(\int_{-\infty}^{t}\) 要分清 → 后者必须补
  \(\dfrac{1}{s}\int_{-\infty}^{0^{-}}f(\tau)\,\mathrm{d}\tau\)，
  只有因果信号该常数项才为零、两式才相同。
\item
  \textbf{用微分器画模拟框图} → 一律用积分器链 → 微分器会放大高频噪声，
  直接型的标准结构是积分器链加反馈。
\item
  \textbf{求 \(H(s)\) 时不置零初始状态} → 先把等效源置零 →
  否则算出的不是系统函数， 而是某个特定初始状态下的响应。
\end{enumerate}

\section{连续时间系统的系统函数}\label{ux8fdeux7eedux65f6ux95f4ux7cfbux7edfux7684ux7cfbux7edfux51fdux6570}

\subsection{考点清单}\label{ux8003ux70b9ux6e05ux5355-5}

\begin{longtable}[]{@{}
  >{\raggedright\arraybackslash}p{(\linewidth - 4\tabcolsep) * \real{0.2500}}
  >{\raggedright\arraybackslash}p{(\linewidth - 4\tabcolsep) * \real{0.4167}}
  >{\centering\arraybackslash}p{(\linewidth - 4\tabcolsep) * \real{0.3333}}@{}}
\toprule\noalign{}
\begin{minipage}[b]{\linewidth}\raggedright
考点
\end{minipage} & \begin{minipage}[b]{\linewidth}\raggedright
常见考法
\end{minipage} & \begin{minipage}[b]{\linewidth}\centering
重要度
\end{minipage} \\
\midrule\noalign{}
\endhead
\bottomrule\noalign{}
\endlastfoot
系统函数的定义与求法 & 由微分方程、框图、电路求 \(H(s)\) & ★★★ \\
极零图 & 由 \(H(s)\) 画极零图；由极零图反求 \(H(s)\) 并用 \(H(0)\) 定标
& ★★★ \\
零极点分布与 \(h(t)\) & 判断响应模式、衰减快慢、是否振荡 & ★★★ \\
零极点与频率特性 & 由极零图粗略画幅频、相频特性曲线 & ★★ \\
全通与最小相位系统 & 判断类型；说明幅度恒定或相移最小 & ★★★ \\
BIBO 稳定性判据 & 判稳定、临界稳定、不稳定 & ★★★ \\
罗斯-霍尔维茨判据 & 列罗斯表判右半平面根个数；首列零、全零行 & ★★★ \\
稳定性参数范围 & 求使系统稳定的参数 \(K\) 范围 & ★★★ \\
反馈系统的系统函数 & 由 \(1+GH=0\) 得特征方程再判稳 & ★★ \\
\end{longtable}

\subsection{核心概念与公式}\label{ux6838ux5fc3ux6982ux5ff5ux4e0eux516cux5f0f-5}

\subsubsection{系统函数的定义与求法}\label{ux7cfbux7edfux51fdux6570ux7684ux5b9aux4e49ux4e0eux6c42ux6cd5}

定义（零状态下响应的拉普拉斯变换与激励的拉普拉斯变换之比）：

\[
H(s)=\frac{Y_{zs}(s)}{F(s)}=\frac{B(s)}{A(s)}
=\frac{b_{m}s^{m}+b_{m-1}s^{m-1}+\cdots+b_{1}s+b_{0}}
{a_{n}s^{n}+a_{n-1}s^{n-1}+\cdots+a_{1}s+a_{0}} .
\]

\(H(s)\) 只由系统本身（结构与参数）决定，与输入、初始状态无关；
\(Y_{zs}(s)=H(s)F(s)\)，\(h(t)=\mathcal{L}^{-1}[H(s)]\)。

三条求法：

\begin{enumerate}
\def\labelenumi{\arabic{enumi}.}
\tightlist
\item
  由微分方程直接读系数：\(A(s)\) 取左边（输出）系数，\(B(s)\)
  取右边（输入）系数。
\item
  由冲激响应：\(H(s)=\mathcal{L}[h(t)]\)。
\item
  由电路或框图：元件及支路全换成 \(s\) 域量，初始状态置零，解代数方程。
\end{enumerate}

按端口分四类：策动点（输入）函数含输入阻抗 \(Z_{11}(s)\)、输入导纳
\(Y_{11}(s)\)；
转移（传输）函数含转移阻抗、转移导纳、电压传输函数、电流传输函数。

\subsubsection{极零图}\label{ux6781ux96f6ux56fe}

把 \(H(s)\) 因式分解：

\[
H(s)=\frac{N(s)}{D(s)}
=H_{0}\frac{\prod_{j=1}^{m}(s-z_{j})}{\prod_{i=1}^{n}(s-p_{i})}.
\]

\(D(s)=0\) 的根 \(p_{1},\dots,p_{n}\) 为极点，\(N(s)=0\) 的根
\(z_{1},\dots,z_{m}\) 为零点， 把它们画在 \(s\)
平面上（极点画叉、零点画圈）就是极零图。

要点：

\begin{enumerate}
\def\labelenumi{\arabic{enumi}.}
\tightlist
\item
  实际系统 \(H(s)\)
  是实系数有理式，故极零点是实数或共轭复数，极零图关于实轴对称。
\item
  极点决定系统的自然响应模式，零点影响哪些频率被抑制。
\item
  极点个数 \(n\) 与零点个数 \(m\) 之差决定 \(h(t)\) 中是否含
  \(\delta(t)\) 及其导数 （\(m\ge n\) 时含）。
\item
  收敛域以最右极点实部为界，故极零图加 \(H_{0}\) 就基本确定了 \(H(s)\)。
\end{enumerate}

\subsubsection{零极点分布与冲激响应波形}\label{ux96f6ux6781ux70b9ux5206ux5e03ux4e0eux51b2ux6fc0ux54cdux5e94ux6ce2ux5f62}

部分分式展开 \(H(s)=\sum_{i}\dfrac{k_{i}}{s-p_{i}}\)，得

\[
h(t)=\mathcal{L}^{-1}[H(s)]=\sum_{i=1}^{n}k_{i}\mathrm{e}^{p_{i}t}u(t),
\]

写成 \(p_{i}=\sigma_{i}+\mathrm{j}\omega_{i}\) 后每一项为
\(k_{i}\mathrm{e}^{\sigma_{i}t}
\mathrm{e}^{\mathrm{j}\omega_{i}t}\)：

\begin{longtable}[]{@{}lll@{}}
\toprule\noalign{}
极点位置 & \(\sigma_{i}\)、\(\omega_{i}\) & \(h(t)\) 波形 \\
\midrule\noalign{}
\endhead
\bottomrule\noalign{}
\endlastfoot
负实轴 & \(\sigma_{i}<0\)，\(\omega_{i}=0\) & 单调衰减 \\
正实轴 & \(\sigma_{i}>0\)，\(\omega_{i}=0\) & 单调增长 \\
左半平面共轭对 & \(\sigma_{i}<0\)，\(\omega_{i}\ne 0\) & 减幅振荡 \\
虚轴共轭对 & \(\sigma_{i}=0\)，\(\omega_{i}\ne 0\) & 等幅振荡 \\
右半平面共轭对 & \(\sigma_{i}>0\)，\(\omega_{i}\ne 0\) & 增幅振荡 \\
\end{longtable}

两个定量结论：极点离虚轴越远（\(|\sigma_{i}|\)
越大），衰减越快，时间常数约为
\(\tau=1/|\sigma_{i}|\)；虚部越大，振荡越快，振荡周期为
\(2\pi/|\omega_{i}|\)。 高阶极点 \(p\)（\(r\) 重）带来
\(t^{r-1}\mathrm{e}^{pt}\) 因子，波形形状不变但多一个多项式包络。

\subsubsection{零极点与频率特性的几何关系}\label{ux96f6ux6781ux70b9ux4e0eux9891ux7387ux7279ux6027ux7684ux51e0ux4f55ux5173ux7cfb}

只在虚轴上观察 \(s=\mathrm{j}\omega\)：

\[
H(\mathrm{j}\omega)=H_{0}\frac{\prod_{j=1}^{m}(\mathrm{j}\omega-z_{j})}
{\prod_{i=1}^{n}(\mathrm{j}\omega-p_{i})},
\] \[
|H(\mathrm{j}\omega)|=H_{0}\frac{B_{1}B_{2}\cdots B_{m}}{A_{1}A_{2}\cdots A_{n}},
\qquad
\varphi(\omega)=\sum_{j=1}^{m}\beta_{j}-\sum_{i=1}^{n}\alpha_{i},
\]

其中
\(B_{j}=|\mathrm{j}\omega-z_{j}|\)、\(A_{i}=|\mathrm{j}\omega-p_{i}|\)
是虚轴上的动点到零点、 极点的距离，\(\beta_{j}\)、\(\alpha_{i}\)
是相应矢量与正实轴的夹角。

读图规则：\(\mathrm{j}\omega\) 靠近极点时 \(A_{i}\)
变小，幅频出现波峰；靠近零点时 \(B_{j}\) 变小，
幅频出现波谷；极零点越靠近虚轴，峰谷越尖锐；极点正好落在虚轴上时幅频趋于无穷大
（零点落在虚轴上则幅频为零）。幅频是 \(\omega\)
的偶函数，相频是奇函数，故只画 \(\omega\ge 0\) 即可。

\subsubsection{全通函数与最小相位函数}\label{ux5168ux901aux51fdux6570ux4e0eux6700ux5c0fux76f8ux4f4dux51fdux6570}

\textbf{全通函数}：零点与极点关于虚轴镜像对称的系统，满足

\[
|H(\mathrm{j}\omega)|=H_{0},
\]

特点：幅度不随频率变化，只改变相位，因此常用来做相位校正而不引入幅度失真。
稳定系统的极点只能在左半平面，故稳定全通系统的零点只能全在右半平面。

\textbf{最小相位函数}：极零点全部位于虚轴以左（左半平面）的系统。
``最小相位''的含义是：在 \(|H(\mathrm{j}\omega)|\)
完全相同的系统中，它的相移 \(\varphi(\omega)\) 最小，
因为左半平面零点提供的相位滞后比右半平面镜像零点小。

判据速记：零点全在左半平面即最小相位；极点与零点关于虚轴镜像即全通。

两者的关系：除平凡情形（常数增益 \(H(s)=K\)，此时既无零点也无极点）外，
全通函数与最小相位函数\textbf{互斥}。理由是稳定全通系统的极点必须在左半平面，
与极点关于虚轴镜像的零点就必然落在右半平面，于是它不可能满足''极零点全在虚轴以左''，
即稳定全通系统一定不是最小相位系统。

\subsubsection{BIBO
稳定性与极点位置判据}\label{bibo-ux7a33ux5b9aux6027ux4e0eux6781ux70b9ux4f4dux7f6eux5224ux636e}

定义：对任意有界激励 \(|f(t)|\le M_{f}\)，响应都有界
\(|y(t)|\le M_{y}\)，则系统 BIBO 稳定。

充要条件（因果系统）：冲激响应绝对可积，

\[
\int_{0}^{+\infty}|h(t)|\,\mathrm{d}t<\infty .
\]

极点判据（因果有理系统）：

\begin{longtable}[]{@{}ll@{}}
\toprule\noalign{}
极点情况 & 结论 \\
\midrule\noalign{}
\endhead
\bottomrule\noalign{}
\endlastfoot
极点全部严格位于左半平面 & 稳定 \\
虚轴上只有一阶极点，其余在左半平面 & 临界稳定 \\
有极点在右半平面，或虚轴上有高阶极点 & 不稳定 \\
\end{longtable}

从零输入响应看同一结论：一阶极点
\(p_{i}=\sigma_{i}+\mathrm{j}\omega_{i}\) 要求
\(\sigma_{i}\le 0\)；\(r\) 阶极点要求 \(\sigma_{i}<0\)，因为
\(t^{r-1}\mathrm{e}^{\sigma_{i}t}\) 在 \(\sigma_{i}=0\) 时无界。

临界稳定的工程含义：参数稍变就可能失稳，且激励恰在虚轴同一位置有一阶极点时
响应会出现重极点而发散，所以不少教材直接把临界稳定归入不稳定。

快速必要条件：\(D(s)\) 的系数全部非零且同号（无缺项、无符号变化）。
这只是必要条件，满足后必须再用罗斯判据。

\subsubsection{罗斯-霍尔维茨判据}\label{ux7f57ux65af-ux970dux5c14ux7ef4ux8328ux5224ux636e}

目的：不求解根，数出 \(D(s)=0\) 在右半平面（实部大于零）的根的个数。

造表规则。设
\(D(s)=a_{n}s^{n}+a_{n-1}s^{n-1}+\cdots+a_{0}\)，把系数按下面方式填前两行：

\[
\begin{array}{c|cccc}
s^{n} & a_{n} & a_{n-2} & a_{n-4} & \cdots \\
s^{n-1} & a_{n-1} & a_{n-3} & a_{n-5} & \cdots
\end{array}
\]

其后每一行由上面两行递推，记上两行为 \(b_{1},b_{2},\dots\) 与
\(c_{1},c_{2},\dots\)，则

\[
d_{1}=\frac{b_{1}c_{2}-b_{2}c_{1}}{b_{1}},
\qquad
d_{2}=\frac{b_{1}c_{3}-b_{2}c_{2}}{b_{1}},
\qquad \cdots
\]

即''左上乘右下减右上乘左下，再除以本行首列元素''。

判据：

\begin{enumerate}
\def\labelenumi{\arabic{enumi}.}
\tightlist
\item
  第一列元素符号不变的次数为零且无特殊情况 →
  全部根在左半平面，系统稳定。
\item
  第一列元素的符号改变次数等于右半平面根的个数，也等于不稳定根的个数。
\item
  首列出现零：用无穷小正量 \(\varepsilon\) 代替该零继续算（令
  \(\varepsilon\to 0\) 判符号）， 或先把 \(D(s)\) 乘以 \((s+1)\)
  再重新造表。
\item
  出现全零行：说明 \(s\)
  平面上存在关于原点对称的根（虚轴极点或实轴上的对称根）。
  取全零行上一行的系数构造辅助多项式 \(P(s)\)，用
  \(\dfrac{\mathrm{d}P(s)}{\mathrm{d}s}\)
  的系数替代全零行继续造表，判有无右半平面根；再解 \(P(s)=0\)
  判断虚轴上极点的阶数： 单极点则临界稳定，重极点则不稳定。
\end{enumerate}

\subsubsection{反馈系统的系统函数与稳定性}\label{ux53cdux9988ux7cfbux7edfux7684ux7cfbux7edfux51fdux6570ux4e0eux7a33ux5b9aux6027}

前向通道 \(G(s)\)、反馈通道 \(K(s)\)、负反馈时

\[
H(s)=\frac{G(s)}{1+G(s)K(s)} .
\]

闭环极点由特征方程

\[
1+G(s)K(s)=0
\]

的根决定，因此反馈系统的稳定性判据与前面完全一致：把 \(1+G(s)K(s)\)
通分后的分子 当作新的
\(D(s)\)，再上罗斯表。负反馈把极点往左推（提高稳定度），
但增益过大时也可能把极点推入右半平面，这就是''求 \(K\)
的稳定范围''这类题的物理含义。

\subsection{典型题型与解法}\label{ux5178ux578bux9898ux578bux4e0eux89e3ux6cd5-5}

\subsubsection{题型 6-1
由极零图反求系统函数并判定系统类型}\label{ux9898ux578b-6-1-ux7531ux6781ux96f6ux56feux53cdux6c42ux7cfbux7edfux51fdux6570ux5e76ux5224ux5b9aux7cfbux7edfux7c7bux578b}

\textbf{题型特征}：给极零图（或文字给出极点零点位置），另给一个定标条件
（如
\(H(0)\)、\(\int_{0}^{\infty}h(t)\,\mathrm{d}t\)、某个频率点的幅度），
要求写出 \(H(s)\)、判断全通或最小相位、写出微分方程。

\textbf{解题步骤}： 1. 由图读全部零点 \(z_{j}\) 与极点
\(p_{i}\)（共轭对合写成 \((s-\sigma)^{2}+\omega^{2}\)）。 2. 写成
\(H(s)=H_{0}\dfrac{\prod(s-z_{j})}{\prod(s-p_{i})}\)，\(H_{0}\) 待定。
3. 用定标条件求 \(H_{0}\)：因果系统的
\(\int_{0}^{\infty}h(t)\,\mathrm{d}t=H(0)\)， 即令 \(s=0\) 代入。 4.
判类型：零点全在左半平面即最小相位；极零关于虚轴镜像即全通。 5. 由
\(\sum a_{i}s^{i}Y=\sum b_{j}s^{j}X\)
逐项对应写出微分方程，注意两边阶次不同。

\textbf{例题}（【附加作业6 第5题】）：实因果 LTI 系统的极零图如图 1
所示： 零点位于 \(s=-3\)，且在 \(s=-1\) 处有二重零点；极点位于
\(s=-2\)， 以及 \(s=-1\pm \mathrm{j}2\) 一对共轭极点。又知
\(\int_{0}^{\infty}h(t)\,\mathrm{d}t=1.5\)。 求
1）系统类型（全通或最小相位）与 \(H(s)\)；2）系统的微分方程。

解 1）先写因式形式。零点：\((s+3)\) 与 \((s+1)^{2}\)； 极点：\((s+2)\)
与 \((s+1-\mathrm{j}2)(s+1+\mathrm{j}2)=(s+1)^{2}+4\)。故

\[
H(s)=H_{0}\frac{(s+1)^{2}(s+3)}{\left[(s+1)^{2}+4\right](s+2)},
\qquad \mathrm{Re}(s)>-1 .
\]

收敛域取最右极点 \(s=-1\) 之右。

判类型：零点
\(s=-3\)、\(s=-1\)（二重）全部位于左半平面，故为最小相位系统；
又零点与极点不关于虚轴镜像，故不是全通系统。

定标：\(h(t)\) 为因果，故

\[
\int_{0}^{\infty}h(t)\,\mathrm{d}t=H(0)
=H_{0}\frac{(0+1)^{2}(0+3)}{\left[(0+1)^{2}+4\right](0+2)}
=H_{0}\frac{3}{5\cdot 2}=H_{0}\frac{3}{10}=1.5,
\]

解得 \(H_{0}=5\)。展开分子分母：

\[
(s+1)^{2}(s+3)=s^{3}+5s^{2}+7s+3,
\] \[
\left[(s+1)^{2}+4\right](s+2)=(s^{2}+2s+5)(s+2)=s^{3}+4s^{2}+9s+10,
\]

故

\[
H(s)=\frac{5\left(s^{3}+5s^{2}+7s+3\right)}{s^{3}+4s^{2}+9s+10},
\qquad \mathrm{Re}(s)>-1 .
\]

解 2）由 \(H(s)=\dfrac{Y(s)}{X(s)}\) 交叉相乘：

\[
\left(s^{3}+4s^{2}+9s+10\right)Y(s)
=5\left(s^{3}+5s^{2}+7s+3\right)X(s),
\]

逐项对应
\(s^{k}\to\dfrac{\mathrm{d}^{k}}{\mathrm{d}t^{k}}\)（零状态下无反变换初始项），得

\[
y'''+4y''+9y'+10y=5x'''+25x''+35x'+15x .
\]

\textbf{同类题}： - 【附加作业6 第6题】由模拟框图求 \(H(s)\)
时，思路是''围绕相加器逐级写回路方程， 再用末级求和得
\(y\)``：两路一阶环并联，\(H(s)=\dfrac{1}{s+2}+\dfrac{1}{s+a}
  =\dfrac{2s+a+2}{(s+2)(s+a)}\)，最后由阶跃输入的稳态值
\(y(+\infty)=H(0)=\dfrac{3}{4}\) 解出 \(a=4\)，得
\(H(s)=\dfrac{2s+6}{(s+2)(s+4)}\)。 该题把''框图求
\(H\)``与''终值定参数''两件事合在一起，是本章综合题的常见结构。

\subsubsection{题型 6-2
由零极点分布判断冲激响应波形与系统稳定性}\label{ux9898ux578b-6-2-ux7531ux96f6ux6781ux70b9ux5206ux5e03ux5224ux65adux51b2ux6fc0ux54cdux5e94ux6ce2ux5f62ux4e0eux7cfbux7edfux7a33ux5b9aux6027}

\textbf{题型特征}：给 \(H(s)\) 或只给极点位置，要求定性说明 \(h(t)\)
由哪些模式组成、 衰减多快、是否振荡、系统是否稳定。

\textbf{解题步骤}： 1. 求 \(D(s)=0\)
的根，判断实根与共轭复根各几对、重数多少。 2. 每个一阶极点给一个
\(\mathrm{e}^{pt}\) 项；共轭对给
\(\mathrm{e}^{\sigma t}\cos(\omega t+\phi)\) 项； \(r\) 阶极点给
\(t^{r-1}\mathrm{e}^{pt}\) 项。 3. 按 \(\sigma\) 判衰减或增长、按
\(\omega\) 判振荡快慢，用文字描述波形。 4.
极点全部严格在左半平面则稳定，虚轴上有单极点则临界稳定，否则不稳定。

\textbf{例题}（自编例）：设
\(H(s)=\dfrac{s+4}{(s+1)\left[(s+2)^{2}+9\right]}\)，说明 \(h(t)\)
的波形与稳定性。

解：极点
\(p_{1}=-1\)（负实轴）、\(p_{2,3}=-2\pm \mathrm{j}3\)（左半平面共轭对），
零点 \(z_{1}=-4\)（左半平面）。部分分式给出

\[
h(t)=\left[k_{1}\mathrm{e}^{-t}+2|k_{2}|\mathrm{e}^{-2t}\cos(3t+\phi)\right]u(t),
\]

波形是''单调慢衰减项叠加减幅振荡项''：\(\mathrm{e}^{-t}\) 的时间常数
\(\tau=1\)， 衰减慢但无振荡；\(\mathrm{e}^{-2t}\cos 3t\) 的包络时间常数
\(\tau=1/2\)，衰减快， 振荡角频率 \(3\ \mathrm{rad/s}\)，周期
\(\dfrac{2\pi}{3}\ \mathrm{s}\)。 极点全部严格位于左半平面，故系统稳定。

\textbf{同类题}： - 【附加作业6 第5题】极点为 \(-2\) 与
\(-1\pm \mathrm{j}2\)，故 \(h(t)\) 含 \(k_{1}\mathrm{e}^{-2t}\) 与
\(2|k_{2}|\mathrm{e}^{-t}\cos(2t+\phi)\) 两类模式，衰减型、稳定。

\subsubsection{题型 6-3
罗斯判据：常规情况与首列出现零}\label{ux9898ux578b-6-3-ux7f57ux65afux5224ux636eux5e38ux89c4ux60c5ux51b5ux4e0eux9996ux5217ux51faux73b0ux96f6}

\textbf{题型特征}：给高阶特征多项式，要求判断稳定性并数出正实部、负实部根的个数；
或表中首列出现零。

\textbf{解题步骤}： 1.
先查必要条件：系数是否全非零且同号（不同号直接判不稳定）。 2.
按奇偶次幂系数的交错方式填前两行，逐行按''左上乘右下减右上乘左下再除首列''递推。
3. 数第一列符号改变次数 =
右半平面根个数；总根数减去它即为左半平面根个数。 4. 首列出现零：用
\(\varepsilon\)（无穷小正数）代替后继续算，令 \(\varepsilon\to 0\)
判符号； 或先把 \(D(s)\) 乘以 \((s+1)\) 再重新造表。

\textbf{例题}（常规情形，【附加作业6 第7题(1)】）：
\(D(s)=s^{4}+7s^{3}+17s^{2}+17s+6=0\)。

解：系数为 \(1,7,17,17,6\)，全正，满足必要条件。造表：

\[
\begin{array}{c|ccc}
s^{4} & 1 & 17 & 6 \\
s^{3} & 7 & 17 & 0 \\
s^{2} & \frac{102}{7} & 6 & 0 \\
s^{1} & \frac{240}{17} & 0 & 0 \\
s^{0} & 6 & 0 & 0
\end{array}
\]

各元素来历：

\[
\frac{7\times 17-1\times 17}{7}=\frac{102}{7},
\qquad
\frac{7\times 6-1\times 0}{7}=6,
\] \[
\frac{\frac{102}{7}\times 17-7\times 6}{\frac{102}{7}}
=\frac{1734-294}{102}=\frac{1440}{102}=\frac{240}{17},
\] \[
\frac{\frac{240}{17}\times 6-\frac{102}{7}\times 0}{\frac{240}{17}}=6 .
\]

递推的统一口径是''左上乘右下减右上乘左下，再除以紧邻上一行的首列元素''；
某一行右侧不足两列时补零，补零不影响首列符号的判断。

第一列 \(1,\ 7,\ \frac{102}{7},\ \frac{240}{17},\ 6\)
全为正，符号无改变， 故系统稳定，有 \(4\) 个负实部特征根、\(0\)
个正实部特征根。 （验算：\(D(s)=(s+1)^{2}(s+2)(s+3)\)，四个根
\(-1,-1,-2,-3\) 全为负实数，与结论一致。）

\textbf{例题}（首列出现零，取自 PPT 第 40
页例）：\(D(s)=s^{4}+s^{3}+2s^{2}+2s+3\)。

解：前三行

\[
\begin{array}{c|ccc}
s^{4} & 1 & 2 & 3 \\
s^{3} & 1 & 2 & 0 \\
s^{2} & 0 & 3 & 0
\end{array}
\]

\(s^{2}\) 行首列为零，无法作除数。用 \(\varepsilon>0\) 代替该零继续：

\[
\begin{array}{c|ccc}
s^{4} & 1 & 2 & 3 \\
s^{3} & 1 & 2 & 0 \\
s^{2} & \varepsilon & 3 & 0 \\
s^{1} & 2-\frac{3}{\varepsilon} & 0 & 0 \\
s^{0} & 3 & 0 & 0
\end{array}
\]

其中 \(s^{1}\) 行首列
\(=\dfrac{\varepsilon\times 2-1\times 3}{\varepsilon}=2-\dfrac{3}{\varepsilon}\)。
令 \(\varepsilon\to 0^{+}\)，该元素趋于 \(-\infty\)。第一列
\(1,\ 1,\ \varepsilon,\ 2-\frac{3}{\varepsilon},\ 3\) 的符号为
\(+,\ +,\ +,\ -,\ +\)， 改变两次，故有 \(2\) 个正实部根，系统不稳定。

\textbf{同类题}： - 【附加作业6 第7题(2)】\(s^{4}+5s^{3}+2s+10=0\)（缺
\(s^{2}\) 项）： 表为
\(s^{4}:1,0,10\)；\(s^{3}:5,2,0\)；\(s^{2}:-\frac{2}{5},10\)；\(s^{1}:127,0\)；\(s^{0}:10\)，
第一列 \(1,5,-\frac{2}{5},127,10\) 变号两次，2 个正实部根、2
个负实部根，不稳定。 - 【附加作业6
第7题(5)】\(s^{5}+2s^{4}+2s^{3}+4s^{2}+11s+10=0\)： 造到 \(s^{3}\)
行时首列为零（\(\frac{2\times 2-1\times 4}{2}=0\)），改用''乘
\((s+1)\)``： 新方程
\(s^{6}+3s^{5}+4s^{4}+6s^{3}+15s^{2}+21s+10=0\)，表为
\(s^{6}:1,4,15,10\)；\(s^{5}:3,6,21\)；\(s^{4}:2,8,10\)；\(s^{3}:-6,6\)；
\(s^{2}:10,10\)；\(s^{1}:12\)；\(s^{0}:10\)，第一列
\(1,3,2,-6,10,12,10\) 变号两次， 2 个正实部根、3 个负实部根，不稳定。

\subsubsection{题型 6-4
罗斯判据：全零行与虚轴极点}\label{ux9898ux578b-6-4-ux7f57ux65afux5224ux636eux5168ux96f6ux884cux4e0eux865aux8f74ux6781ux70b9}

\textbf{题型特征}：造表过程中某一整行全为零；结论往往是''临界稳定''或''不稳定''，
必须判断虚轴极点的阶数。

\textbf{解题步骤}： 1.
出现全零行时，取上一行的系数（对应偶数次幂）写辅助多项式 \(P(s)\)。 2.
求导
\(\dfrac{\mathrm{d}P(s)}{\mathrm{d}s}\)，用其系数填全零行，继续造表。 3.
数第一列变号次数，判断有无右半平面根。 4. 解 \(P(s)=0\)
得虚轴（或对称）根：单根则临界稳定，重根则不稳定。

\textbf{例题}（【附加作业6 第7题(3)】）：
\(D(s)=s^{4}+2s^{3}+7s^{2}+10s+10=0\)。

解：造表

\[
\begin{array}{c|ccc}
s^{4} & 1 & 7 & 10 \\
s^{3} & 2 & 10 & 0 \\
s^{2} & 2 & 10 & 0 \\
s^{1} & 0 & 0 & 0
\end{array}
\]

\(s^{2}\) 行的两个元素：\(\dfrac{2\times 7-1\times 10}{2}=2\)，
\(\dfrac{2\times 10-1\times 0}{2}=10\)； \(s^{1}\) 行本应为
\(\dfrac{2\times 10-2\times 10}{2}=0\)，整行为零。

取全零行上一行（\(s^{2}\) 行）构造辅助多项式

\[
P(s)=2s^{2}+10,
\qquad
\frac{\mathrm{d}P(s)}{\mathrm{d}s}=4s,
\]

用 \(4\) 填 \(s^{1}\) 行；\(s^{0}\) 行再由 \(s^{1}\) 行与 \(s^{2}\)
行同法算出： \(\dfrac{4\times 10-2\times 0}{4}=10\)。

\[
\begin{array}{c|ccc}
s^{4} & 1 & 7 & 10 \\
s^{3} & 2 & 10 & 0 \\
s^{2} & 2 & 10 & 0 \\
s^{1} & 4 & 0 & 0 \\
s^{0} & 10 & 0 & 0
\end{array}
\]

第一列 \(1,\ 2,\ 2,\ 4,\ 10\) 全正，符号无改变，故右半平面无极点。 再解
\(P(s)=2s^{2}+10=0\) 得 \(s_{1,2}=\pm \mathrm{j}\sqrt{5}\)，
这是虚轴上的一对单极点，故系统临界稳定（有 \(2\) 个纯虚根、\(2\)
个负实部根）。

验证：\(D(s)=\left(s^{2}+5\right)\left(s^{2}+2s+2\right)\)，
后一因子根为 \(-1\pm \mathrm{j}1\)，确为 2 个负实部根。

\textbf{同类题}： - 【附加作业6
第7题(4)】\(4s^{5}+6s^{4}+2s^{3}+4s^{2}+11s+10=0\)： 表为
\(s^{5}:4,2,11\)；\(s^{4}:6,4,10\)；\(s^{3}:-\frac{2}{3},\frac{13}{3}\)；
\(s^{2}:43,10\)；\(s^{1}:\frac{193}{43},0\)；\(s^{0}:10\)。 第一列
\(4,6,-\frac{2}{3},43,\frac{193}{43},10\) 变号两次， 故 2 个正实部根、3
个负实部根，不稳定（无全零行，属常规情形）。

\subsubsection{题型 6-5
用罗斯判据求稳定参数范围}\label{ux9898ux578b-6-5-ux7528ux7f57ux65afux5224ux636eux6c42ux7a33ux5b9aux53c2ux6570ux8303ux56f4}

\textbf{题型特征}：特征方程的某个系数含待定参数 \(K\)（常是增益），
要求系统稳定时 \(K\) 的取值范围。

\textbf{解题步骤}： 1. 把含 \(K\) 的特征多项式写成
\(D(s)\)，必要时先通分（反馈系统用 \(1+GK=0\)）。 2. 造罗斯表，\(K\)
当作常数参与运算。 3. 令第一列每个元素都大于零，得到若干关于 \(K\)
的不等式。 4. 取交集，得 \(K\)
的稳定范围；边界值对应临界稳定（虚轴极点）。

\textbf{例题}（【附加作业6
第8题】）：\(s^{3}+4s^{2}+4s+K=0\)，求系统稳定的 \(K\) 值范围。

解：造表

\[
\begin{array}{c|cc}
s^{3} & 1 & 4 \\
s^{2} & 4 & K \\
s^{1} & \frac{16-K}{4} & 0 \\
s^{0} & K & 0
\end{array}
\]

其中 \(s^{1}\) 行首列
\(=\dfrac{4\times 4-1\times K}{4}=\dfrac{16-K}{4}\)。

稳定要求第一列全部为正：

\[
1>0,\qquad 4>0,\qquad \frac{16-K}{4}>0,\qquad K>0 .
\]

前两条自动满足，后两条给出

\[
K<16,
\qquad
K>0 .
\]

即稳定范围为 \(0<K<16\)。

边界含义：\(K=16\) 时 \(s^{1}\) 行为零，由 \(s^{2}\) 行得辅助多项式
\(P(s)=4s^{2}+16=0\)，根为 \(\pm \mathrm{j}2\)，系统临界稳定； \(K=0\)
时 \(s^{0}\) 行为零，原点处有极点，同样临界稳定。

\textbf{同类题}： - 【附加作业6
第6题】同类思想的反用：已知稳态值定参数，再用极点判稳。 该题
\(H(s)=\dfrac{2s+6}{(s+2)(s+4)}\) 的极点为
\(-2\)、\(-4\)，全部在左半平面， 故对阶跃输入可以放心用终值定理。

\subsubsection{题型 6-6
由系统函数与信号运算求像函数（积分、频移）}\label{ux9898ux578b-6-6-ux7531ux7cfbux7edfux51fdux6570ux4e0eux4fe1ux53f7ux8fd0ux7b97ux6c42ux50cfux51fdux6570ux79efux5206ux9891ux79fb}

\textbf{题型特征}：先由微分方程写出 \(H(s)\)，再把 \(h(t)\)
经积分、乘指数、微分等运算构成 新信号 \(g(t)\)，求
\(G(s)\)。考的是性质运用的准确性而非时域积分。

\textbf{解题步骤}： 1. 由微分方程读 \(A(s)\)、\(B(s)\)，得
\(H(s)=B(s)/A(s)\)。 2.
逐条套性质：\(\int_{-\infty}^{t}h(\tau)\,\mathrm{d}\tau\leftrightarrow
   \dfrac{H(s)}{s}+\dfrac{1}{s}\int_{-\infty}^{0^{-}}h(\tau)\,\mathrm{d}\tau\)，
因果系统该附加项为零，才简写为 \(\dfrac{H(s)}{s}\)； 另有
\(\mathrm{e}^{-at}h(t)\leftrightarrow H(s+a)\)，\(h'(t)\leftrightarrow sH(s)\)。
3. 频移与乘 \(\dfrac{1}{s}\) 的顺序不能颠倒：先做积分再乘指数， 就是把
\(\dfrac{H(s)}{s}\) 中的 \(s\) 整体换成 \(s+a\)。 4. 代简并写出收敛域。

\textbf{例题}（【附加作业6 第9题】）：因果 LTI 系统
\(y''(t)+5ay'(t)+6a^{2}y(t)=x'(t)-ax(t)\)，\(a>0\)，\(h(t)\)
为单位冲激响应。 若
\(g(t)=\mathrm{e}^{-at}\int_{-\infty}^{t}h(\tau)\,\mathrm{d}\tau\)，求
\(G(s)\)。

解：先求
\(H(s)\)。零状态下两边取变换，\(s^{2}Y+5asY+6a^{2}Y=(s-a)X\)，故

\[
H(s)=\frac{s-a}{s^{2}+5as+6a^{2}}=\frac{s-a}{(s+2a)(s+3a)},
\qquad \mathrm{Re}(s)>-2a .
\]

再算积分项。因为 \(h(t)\) 因果，

\[
\int_{-\infty}^{t}h(\tau)\,\mathrm{d}\tau=h(t)*u(t)
\ \longleftrightarrow\ \frac{H(s)}{s},
\]

于是

\[
g(t)=\mathrm{e}^{-at}\left[\int_{-\infty}^{t}h(\tau)\,\mathrm{d}\tau\right]
\ \longleftrightarrow\
G(s)=\frac{H(s+a)}{s+a},
\]

这里用的是频移性质：时域乘 \(\mathrm{e}^{-at}\) 等价于把 \(s\) 换成
\(s+a\)， 同时分母里的 \(s\) 也要一起换。代入

\[
H(s+a)=\frac{(s+a)-a}{(s+a+2a)(s+a+3a)}=\frac{s}{(s+3a)(s+4a)},
\]

故

\[
G(s)=\frac{1}{s+a}\cdot\frac{s}{(s+3a)(s+4a)}
=\frac{s}{(s+a)(s+3a)(s+4a)},
\qquad \mathrm{Re}(s)>-a .
\]

\textbf{同类题}： - 【附加作业6 第3题】同一套代数的时域版本：先取变换得
\(Y(s)\)， 再按极点拆分（见题型
5-4）；两题的区别只在于本题不要求反变换。

\subsection{易错点}\label{ux6613ux9519ux70b9-5}

\begin{enumerate}
\def\labelenumi{\arabic{enumi}.}
\tightlist
\item
  \textbf{只查系数同号就下''稳定''结论} →
  同号只是必要条件，必须再列罗斯表 → 反例 \(s^{3}+s^{2}+2s+8\)
  系数全正，但 \(s^{1}\) 行 \(=\frac{1\times 2-1\times 8}{1}=-6<0\)，
  变号两次，不稳定。
\item
  \textbf{罗斯表首列出现零就硬除} → 用 \(\varepsilon\) 代替或乘
  \((s+1)\) 重建 → 除零无意义，两种修补法的结论必须一致。
\item
  \textbf{出现全零行就判''稳定''或''不稳定''} → 取上一行写辅助多项式再判
  → 全零行只说明有对称根，虚轴上单极点仍属临界稳定。
\item
  \textbf{变号次数数错} → 只逐行比较第一列相邻两元素的正负 →
  符号变化次数等于右半平面根个数，跳行比较会漏。
\item
  \textbf{极点个数与根个数混用} → 变号次数对应 \(D(s)\)
  的根总数（含重数）→ 重极点要按重数计数。
\item
  \textbf{用 \(H(0)\) 定标时忘了因果前提} → 只有 \(h(t)\) 因果时
  \(\int_{0}^{\infty}h=H(0)\) → 非因果信号该积分不收敛，不能用。
\item
  \textbf{全通与最小相位判反} → 全通看''零点与极点关于虚轴镜像''，
  最小相位看''极零点全在左半平面'' →
  稳定全通的零点必在右半平面，故它一定不是
  最小相位系统；除平凡情形（常数增益
  \(H(s)=K\)）外，两者互斥，不能同时成立。
\item
  \textbf{把虚轴一阶极点当稳定} → 判为临界稳定 → BIBO
  要求极点严格在左半平面。
\item
  \textbf{由 \(H(s)\) 写微分方程时左右搞反} → 分母系数配输出
  \(y\)，分子系数配输入 \(x\) → 移项后全部归到等号左边，\(x\)
  的导数项阶次与分子一致。
\item
  \textbf{只画幅频不说明偶奇性} → 幅频是 \(\omega\)
  的偶函数、相频是奇函数 → 因此只需画
  \(\omega\ge 0\)，且负频率部分是镜像（或反号）。
\end{enumerate}

\section*{附录
A　专题横向对比}\label{ux9644ux5f55-a-ux4e13ux9898ux6a2aux5411ux5bf9ux6bd4}
\addcontentsline{toc}{section}{附录 A　专题横向对比}

正文按章回答''这一章怎么考''，本附录回答另一半问题：同一件事在六个域里各放在哪里、彼此怎么换算。
三个附录都不重述推导，只做横向对照与考前压缩；需要完整推导与例题时见相应章节。

\subsection*{A.1
求零状态响应的三条路径}\label{a.1-ux6c42ux96f6ux72b6ux6001ux54cdux5e94ux7684ux4e09ux6761ux8defux5f84}
\addcontentsline{toc}{subsection}{A.1 求零状态响应的三条路径}

求 \(y_{zs}(t)\)
是本课程出现频率最高的计算，三条路径给出的结果必然相同，差别只在
``题目给了什么、要求以什么形式交付''。

\begin{longtable}[]{@{}
  >{\raggedright\arraybackslash}p{(\linewidth - 8\tabcolsep) * \real{0.1304}}
  >{\raggedright\arraybackslash}p{(\linewidth - 8\tabcolsep) * \real{0.2174}}
  >{\raggedright\arraybackslash}p{(\linewidth - 8\tabcolsep) * \real{0.2174}}
  >{\raggedright\arraybackslash}p{(\linewidth - 8\tabcolsep) * \real{0.1304}}
  >{\raggedright\arraybackslash}p{(\linewidth - 8\tabcolsep) * \real{0.3043}}@{}}
\toprule\noalign{}
\begin{minipage}[b]{\linewidth}\raggedright
路径
\end{minipage} & \begin{minipage}[b]{\linewidth}\raggedright
适用前提
\end{minipage} & \begin{minipage}[b]{\linewidth}\raggedright
计算步骤
\end{minipage} & \begin{minipage}[b]{\linewidth}\raggedright
优势
\end{minipage} & \begin{minipage}[b]{\linewidth}\raggedright
典型适用信号
\end{minipage} \\
\midrule\noalign{}
\endhead
\bottomrule\noalign{}
\endlastfoot
时域：卷积 & 已有 \(h(t)\)；两函数图形简单或可分段积分 &
定支撑区，排四个分界点，逐段积 & 不用查表，分段点定准后只算定积分 &
矩形、三角、分段直线 \\
频域：\(H(\mathrm{j}\omega)\) & \(f(t)\)
绝对可积，或取其广义变换；系统稳定 & 写 \(H(\mathrm{j}\omega)\)，乘
\(F\)，反变换 & 调制、滤波、抽样类题一步到位 &
\(\cos\)、\(\sin\)、门函数、带限信号 \\
复频域：\(H(s)\) & \(H(s)\) 为有理函数，\(F(s)\) 存在 & 乘得
\(Y_{zs}(s)\)，部分分式，查表 & 反变换表最全，同时带 \(u(t)\) 收口 &
指数、阶跃、\(t\) 乘指数类 \\
\end{longtable}

三条路径的收口方式不同，这是丢分的常见来源：时域路径的答案本来是分段函数，写完全域积分就
错；频域路径的答案里常留 \(\delta(\omega)\) 与
\(\mathrm{Sa}\)，需要靠对称性倒回来；复频域路径的 答案每一项都必须带
\(u(t)\)。

\textbf{什么时候只能用（或必须优先用）拉普拉斯变换}：

\begin{enumerate}
\def\labelenumi{\arabic{enumi}.}
\tightlist
\item
  \textbf{不满足绝对可积}。\(\mathrm{e}^{2t}u(t)\)、\(t^{2}u(t)\)、\(u(t)\)
  这类信号的傅里叶变换只在广 义函数意义下存在，频谱里会出现 \(\delta\)
  及其导数，反变换反而更繁；直接走 \(s\) 域，收敛域
  \(\mathrm{Re}(s)>\sigma_{0}\) 把''为什么不收敛''一并交代清楚。
\item
  \textbf{含初始条件}。傅里叶变换是零状态工具，初始储能在频域没有位置，无法交出
  \(y_{zi}(t)\)； 拉氏变换的微分性质自带 \(f^{(k)}(0^{-})\)
  项，全响应一次算完，再按极点归属拆成 \(y_{zi}\) 与 \(y_{zs}\)。
\item
  \textbf{需要看零输入响应、判稳定性、求稳态值}。零输入响应的极点必须靠
  \(A(s)=0\) 的根给出； 因果与稳定要靠收敛域与极点位置判定；终值定理只在
  \(s\) 域有。
\end{enumerate}

反过来，若题目只问某个频率的正弦稳态输出，用相量法或频率响应最快，此时不必做任何反变换。

\subsection*{A.2
三大变换性质对照表}\label{a.2-ux4e09ux5927ux53d8ux6362ux6027ux8d28ux5bf9ux7167ux8868}
\addcontentsline{toc}{subsection}{A.2 三大变换性质对照表}

下表左右两列是同一性质的两种写法。看这张表的目的不是背，而是记牢三处''不对称''。

\begin{longtable}[]{@{}
  >{\raggedright\arraybackslash}p{(\linewidth - 6\tabcolsep) * \real{0.1304}}
  >{\raggedright\arraybackslash}p{(\linewidth - 6\tabcolsep) * \real{0.2609}}
  >{\raggedright\arraybackslash}p{(\linewidth - 6\tabcolsep) * \real{0.3913}}
  >{\raggedright\arraybackslash}p{(\linewidth - 6\tabcolsep) * \real{0.2174}}@{}}
\toprule\noalign{}
\begin{minipage}[b]{\linewidth}\raggedright
性质
\end{minipage} & \begin{minipage}[b]{\linewidth}\raggedright
傅里叶变换
\end{minipage} & \begin{minipage}[b]{\linewidth}\raggedright
单边拉普拉斯变换
\end{minipage} & \begin{minipage}[b]{\linewidth}\raggedright
关键差异
\end{minipage} \\
\midrule\noalign{}
\endhead
\bottomrule\noalign{}
\endlastfoot
线性 & \(aF_{1}+bF_{2}\) & \(aF_{1}(s)+bF_{2}(s)\) & 无差异 \\
时移 & \(F(\mathrm{j}\omega)\mathrm{e}^{-\mathrm{j}\omega t_{0}}\) &
\(f(t-t_{0})u(t-t_{0})\to F(s)\mathrm{e}^{-st_{0}}\) & 单边必须补
\(u(t-t_{0})\) \\
频移 & \(F\bigl(\mathrm{j}(\omega-\omega_{0})\bigr)\) & \(F(s-s_{0})\) &
单边要求 \(\mathrm{Re}(s_{0})\) 平移收敛域 \\
尺度变换 &
\(\dfrac{1}{\lvert a\rvert}F\left(\mathrm{j}\dfrac{\omega}{a}\right)\) &
\(\dfrac{1}{a}F\left(\dfrac{s}{a}\right)\)，\(a>0\) & 傅里叶允许 \(a<0\)
并产生反褶 \\
时域卷积 & \(F_{1}F_{2}\) & \(F_{1}(s)F_{2}(s)\) & 单边默认两信号因果 \\
频域卷积 & \(\dfrac{1}{2\pi}F_{1}*F_{2}\) & 复卷积，极少考 & 因子
\(\frac{1}{2\pi}\) 与卷积定理不对称 \\
时域微分 & \((\mathrm{j}\omega)^{n}F(\mathrm{j}\omega)\) &
\(s^{n}F(s)-\sum s^{n-1-k}f^{(k)}(0^{-})\) &
单边带初始项，这是最大差异 \\
时域积分 &
\(\dfrac{F(\mathrm{j}\omega)}{\mathrm{j}\omega}+\pi F(0)\delta(\omega)\)
& \(\dfrac{F(s)}{s}\) & 下限 \(0^{-}\) 与 \(-\infty\) 的区别 \\
对称性 & \(F(\mathrm{j}t)\leftrightarrow 2\pi f(-\omega)\) & 无对应关系
& 单边定义破坏了对称性 \\
\end{longtable}

\textbf{三处必须记住的差异}：

\begin{enumerate}
\def\labelenumi{\arabic{enumi}.}
\tightlist
\item
  \textbf{单边拉氏的微分性质带初始条件项}。二阶形式是
  \(s^{2}F(s)-sf(0^{-})-f'(0^{-})\)，初始项不是记忆负担而是复频域方法的价值所在：正因为有
  它，一个微分方程加一组初始条件才能一步化成代数方程。
\item
  \textbf{单边拉氏的积分性质下限是 \(0^{-}\)，不是 \(-\infty\)}。
  \(\mathcal{L}\left[\int_{0^{-}}^{t}f(\tau)\mathrm{d}\tau\right]=\dfrac{F(s)}{s}\)；
  若写成 \(\int_{-\infty}^{t}\)，非因果信号会多出
  \(\dfrac{1}{s}\int_{-\infty}^{0^{-}}f(\tau)\mathrm{d}\tau\)
  这一项，因果信号下该项为零， 所以平时看不出差别，换题就错。
\item
  \textbf{傅里叶变换的尺度变换对 \(a<0\) 有反褶}。因子是
  \(\frac{1}{\lvert a\rvert}\)， \(a<0\) 时
  \(F\left(\mathrm{j}\frac{\omega}{a}\right)\)
  自带一次频率反褶；单边拉氏的尺度性质 只对 \(a>0\) 成立，因为 \(t<0\)
  段已被定义域丢掉，写 \(a<0\) 没有意义。
\end{enumerate}

\subsection*{A.3
常用变换对三列对照表}\label{a.3-ux5e38ux7528ux53d8ux6362ux5bf9ux4e09ux5217ux5bf9ux7167ux8868}
\addcontentsline{toc}{subsection}{A.3 常用变换对三列对照表}

下表把正文第 4 章、第 5
章的两张变换对表并排放，方便看清''同一个信号在两个域里长什么样''。
\(F(s)\) 列一律指单边拉普拉斯变换，故只反映 \(t\ge 0\)
部分；收敛域一列缺省写''---``。

\begin{longtable}[]{@{}
  >{\raggedright\arraybackslash}p{(\linewidth - 6\tabcolsep) * \real{0.1702}}
  >{\raggedright\arraybackslash}p{(\linewidth - 6\tabcolsep) * \real{0.4894}}
  >{\raggedright\arraybackslash}p{(\linewidth - 6\tabcolsep) * \real{0.1702}}
  >{\raggedright\arraybackslash}p{(\linewidth - 6\tabcolsep) * \real{0.1702}}@{}}
\toprule\noalign{}
\begin{minipage}[b]{\linewidth}\raggedright
\(f(t)\)
\end{minipage} & \begin{minipage}[b]{\linewidth}\raggedright
\(F(\mathrm{j}\omega)\)
\end{minipage} & \begin{minipage}[b]{\linewidth}\raggedright
\(F(s)\)
\end{minipage} & \begin{minipage}[b]{\linewidth}\raggedright
收敛域
\end{minipage} \\
\midrule\noalign{}
\endhead
\bottomrule\noalign{}
\endlastfoot
\(\delta(t)\) & \(1\) & \(1\) & 整个 \(s\) 平面 \\
\(\delta(t-t_{0})\) & \(\mathrm{e}^{-\mathrm{j}\omega t_{0}}\) &
\(\mathrm{e}^{-st_{0}}\) & 整个 \(s\) 平面 \\
\(u(t)\) & \(\dfrac{1}{\mathrm{j}\omega}+\pi\delta(\omega)\) &
\(\dfrac{1}{s}\) & \(\mathrm{Re}(s)>0\) \\
\(1\) & \(2\pi\delta(\omega)\) & \(\dfrac{1}{s}\) &
\(\mathrm{Re}(s)>0\) \\
\(\mathrm{sgn}(t)\) & \(\dfrac{2}{\mathrm{j}\omega}\) & \(\dfrac{1}{s}\)
& \(\mathrm{Re}(s)>0\) \\
\(\mathrm{e}^{-at}u(t)\) & \(\dfrac{1}{a+\mathrm{j}\omega}\) &
\(\dfrac{1}{s+a}\) & \(\mathrm{Re}(s)>-a\) \\
\(\mathrm{e}^{-a\lvert t\rvert}\) & \(\dfrac{2a}{a^{2}+\omega^{2}}\) &
--- & --- \\
\(t\,u(t)\) & \(\mathrm{j}\pi\delta'(\omega)-\dfrac{1}{\omega^{2}}\) &
\(\dfrac{1}{s^{2}}\) & \(\mathrm{Re}(s)>0\) \\
\(t^{n}\mathrm{e}^{-at}u(t)\) &
\(\dfrac{n!}{(a+\mathrm{j}\omega)^{n+1}}\) & \(\dfrac{n!}{(s+a)^{n+1}}\)
& \(\mathrm{Re}(s)>-a\) \\
\(\cos\omega_{0}t\,u(t)\) &
\(\dfrac{\pi}{2}\bigl[\delta(\omega-\omega_{0})+\delta(\omega+\omega_{0})\bigr]+\dfrac{\mathrm{j}\omega}{\omega_{0}^{2}-\omega^{2}}\)
& \(\dfrac{s}{s^{2}+\omega_{0}^{2}}\) & \(\mathrm{Re}(s)>0\) \\
\(\sin\omega_{0}t\,u(t)\) &
\(\dfrac{\pi}{2\mathrm{j}}\bigl[\delta(\omega-\omega_{0})-\delta(\omega+\omega_{0})\bigr]+\dfrac{\omega_{0}}{\omega_{0}^{2}-\omega^{2}}\)
& \(\dfrac{\omega_{0}}{s^{2}+\omega_{0}^{2}}\) & \(\mathrm{Re}(s)>0\) \\
\(\mathrm{e}^{-at}\cos\omega_{0}t\,u(t)\) & 取 \(F(s)\) 令
\(s=\mathrm{j}\omega\) & \(\dfrac{s+a}{(s+a)^{2}+\omega_{0}^{2}}\) &
\(\mathrm{Re}(s)>-a\) \\
\(\mathrm{e}^{-at}\sin\omega_{0}t\,u(t)\) & 取 \(F(s)\) 令
\(s=\mathrm{j}\omega\) &
\(\dfrac{\omega_{0}}{(s+a)^{2}+\omega_{0}^{2}}\) &
\(\mathrm{Re}(s)>-a\) \\
\(u(t)-u(t-\tau)\) &
\(\tau\mathrm{Sa}\left(\dfrac{\omega\tau}{2}\right)\mathrm{e}^{-\frac{\mathrm{j}\omega\tau}{2}}\)
& \(\dfrac{1-\mathrm{e}^{-s\tau}}{s}\) & 整个 \(s\) 平面 \\
\(\mathrm{Sa}(\omega_{c}t)\) &
\(\dfrac{\pi}{\omega_{c}}\)，\(\lvert\omega\rvert<\omega_{c}\)，其余为
\(0\) & --- & --- \\
\end{longtable}

表中十三类信号覆盖了附加作业与真题里出现的全部基本波形。四点读表说明：

\begin{enumerate}
\def\labelenumi{\arabic{enumi}.}
\tightlist
\item
  \textbf{矩形脉冲一行取的是因果门}
  \(u(t)-u(t-\tau)\)。若改用关于原点对称的门函数 \(G_{\tau}(t)\)，则
  \(G_{\tau}(t)\leftrightarrow\tau\mathrm{Sa}\left(\frac{\omega\tau}{2}\right)\)，
  二者只差平移因子
  \(\mathrm{e}^{-\frac{\mathrm{j}\omega\tau}{2}}\)，见``常用信号的傅里叶变换对''一节。
\item
  \textbf{\(1\)、\(\mathrm{sgn}(t)\)、\(u(t)\) 三行的 \(F(s)\) 都是
  \(1/s\)}，这不是巧合：单边变换只积 \(t\ge 0\) 段，三个信号在 \(t>0\)
  上完全相同。若题目给的是双边变换，这三行必须分开写。
\item
  \textbf{\(\mathrm{j}\pi\delta'(\omega)\) 项的来历}：由频域微分性质
  \(t f(t)\leftrightarrow \mathrm{j}\dfrac{\mathrm{d}F}{\mathrm{d}\omega}\)
  作用在 \(u(t)\) 上 得到，考场上直接推比背表更稳。
\item
  \textbf{``取 \(F(s)\) 令
  \(s=\mathrm{j}\omega\)''成立的前提是收敛域含虚轴}，即 \(a>0\)。\(a<0\)
  时该信号 无傅里叶变换，不能这样替换（见 A.4 的稳定判据）。
\end{enumerate}

\subsection*{A.4
系统特性判据对照}\label{a.4-ux7cfbux7edfux7279ux6027ux5224ux636eux5bf9ux7167}
\addcontentsline{toc}{subsection}{A.4 系统特性判据对照}

同一个性质往往有两条判据：拿 \(h(t)\)
看是时域判据，拿收敛域与极点看是变换域判据。
考场上一句话用哪条，取决于题目给的是波形还是 \(H(s)\)。

\begin{longtable}[]{@{}
  >{\raggedright\arraybackslash}p{(\linewidth - 6\tabcolsep) * \real{0.1579}}
  >{\raggedright\arraybackslash}p{(\linewidth - 6\tabcolsep) * \real{0.2632}}
  >{\raggedright\arraybackslash}p{(\linewidth - 6\tabcolsep) * \real{0.3158}}
  >{\raggedright\arraybackslash}p{(\linewidth - 6\tabcolsep) * \real{0.2632}}@{}}
\toprule\noalign{}
\begin{minipage}[b]{\linewidth}\raggedright
特性
\end{minipage} & \begin{minipage}[b]{\linewidth}\raggedright
时域判据
\end{minipage} & \begin{minipage}[b]{\linewidth}\raggedright
变换域判据
\end{minipage} & \begin{minipage}[b]{\linewidth}\raggedright
易混提示
\end{minipage} \\
\midrule\noalign{}
\endhead
\bottomrule\noalign{}
\endlastfoot
线性 & 齐次性与叠加性同时成立 & \(H(s)\) 与输入幅度、初态无关；\(Y=HF\)
恒成立 & 含非零常数项只算增量线性 \\
时不变 & \(f(t-t_{0})\to y(t-t_{0})\) &
\(H(s)\)、\(H(\mathrm{j}\omega)\) 的系数全为常数 & 系数显含 \(t\)
即变时 \\
因果 & \(h(t)=0\)，对一切 \(t<0\) & \(H(s)\) 的收敛域是某右半平面
\(\mathrm{Re}(s)>\sigma_{0}\) & 等价于 \(\sigma_{0}\) 取最右极点实部 \\
稳定 &
\(\displaystyle\int_{-\infty}^{+\infty}\lvert h(t)\rvert\mathrm{d}t<\infty\)
& 收敛域含虚轴；有理系统等价于极点全在左半平面 & 收敛域含虚轴才可令
\(s=\mathrm{j}\omega\) \\
\end{longtable}

\textbf{因果与稳定的联合判据}（考试最爱考的一句）：因果系统的收敛域必为右半平面，稳定要求收敛域
含虚轴，两条同时成立的充要条件是\textbf{极点全部严格位于左半平面}。这也解释了为什么''因果且稳定''
只需看极点：收敛域与极点的相对位置已被因果性唯一确定。

\textbf{BIBO 稳定、渐近稳定、临界稳定三者不能混用}：

\begin{itemize}
\tightlist
\item
  \textbf{BIBO 稳定}只问''有界输入是否必有有界输出''，等价于 \(h(t)\)
  绝对可积；
\item
  \textbf{渐近稳定}更强，要求
  \(h(t)\to 0\)（\(t\to+\infty\)），等价于极点全部在左半平面；
\item
  \textbf{临界稳定}指虚轴上允许一阶极点：此时 \(h(t)\)
  有等幅振荡或非零常值，绝对可积被破坏、 \(h(t)\)
  也不趋于零，所以\textbf{临界稳定不是 BIBO
  稳定}，只是''不发散''。不少教材直接把临界稳定
  归入不稳定，答题时把''临界稳定（不满足 BIBO）``写全最安全。
\item
  判据的顺序也要记牢：系数同号只是必要条件，必须再列罗斯表才能下结论（见``罗斯-霍尔维茨判据''一节）。
\end{itemize}

\subsection*{A.5
一图流判断流程}\label{a.5-ux4e00ux56feux6d41ux5224ux65adux6d41ux7a0b}
\addcontentsline{toc}{subsection}{A.5 一图流判断流程}

拿到一道题，先花三十秒回答三个问题，再动手。

\begin{itemize}
\tightlist
\item
  \textbf{第一问：题目要的交付物是什么？}

  \begin{itemize}
  \tightlist
  \item
    要 \(y(t)\)（响应）→ 走下面第二条主线。
  \item
    要 \(h(t)\)、\(H(s)\)、极零图 → 部分分式或长除法，见第 2、6 章。
  \item
    要''是否稳定''\,``参数范围''→ 罗斯判据，跳到第四条主线。
  \item
    要''频谱''\,``带宽''``功率''→ 第 3、4
    章，做级数或变换即可，不涉及求响应。
  \end{itemize}
\item
  \textbf{第二问：题目给了哪些已知量？}

  \begin{itemize}
  \tightlist
  \item
    给了 \(h(t)\) 与 \(f(t)\) 的波形 →
    时域卷积（图形法），四个分界点定区间。
  \item
    给了微分方程与初始条件 → 拉氏变换一次算全响应，再拆
    \(y_{zi}\)、\(y_{zs}\)。
  \item
    给了 \(H(\mathrm{j}\omega)\) 或电路 → 频域相乘，反变换用部分分式。
  \item
    给了 \(H(s)\) 与 \(F(s)\) → 直接乘，部分分式查表，最快。
  \item
    只给了 \(h(t)\) 的表达式，没给 \(f(t)\) → 多半是求 \(g(t)=\int h\)
    或判稳定。
  \end{itemize}
\item
  \textbf{第三问：激励是什么类型的信号？}

  \begin{itemize}
  \tightlist
  \item
    直流、阶跃、指数 → 拉氏或终值定理（阶跃输入下
    \(y(+\infty)=H(0)\)）。
  \item
    单频正弦或余弦 → 频率响应：幅度乘
    \(\lvert H(\mathrm{j}\omega)\rvert\)，相位加 \(\varphi(\omega)\)。
  \item
    周期非正弦 → 傅里叶级数逐项加权，不要做积分。
  \item
    有限长脉冲、门函数 → 傅里叶变换 +
    卷积定理；或时域卷积，看哪个支撑区简单。
  \item
    抽样后的信号 → 频域复制与混叠，先算 \(f_{s}\) 与 \(2f_{m}\)
    的大小关系。
  \end{itemize}
\item
  \textbf{第四条主线（判稳定性，独立于求响应）}

  \begin{itemize}
  \tightlist
  \item
    先写 \(D(s)\)，检查系数是否同号无缺项：不同号直接判不稳定。
  \item
    再列罗斯表，数首列变号次数，得到右半平面根的个数。
  \item
    首列出现零：用 \(\varepsilon\) 代替，或乘 \((s+1)\) 重建。
  \item
    出现全零行：取上一行写辅助多项式，对其求导替代全零行，再解
    \(P(s)=0\) 看虚轴极点阶数。
  \item
    反馈系统先由 \(1+G(s)K(s)=0\) 求特征方程，再回到''写
    \(D(s)\)``这一步。
  \end{itemize}
\end{itemize}

三条三十秒口诀：

\begin{enumerate}
\def\labelenumi{\arabic{enumi}.}
\tightlist
\item
  \textbf{有初态找拉氏，无初态看波形}：初态一出现，频域法就交不出答案。
\item
  \textbf{要稳态用 \(H(\mathrm{j}\omega)\)，要瞬态用
  \(H(s)\)}：频率响应天然只描述零状态稳态。
\item
  \textbf{图形简单走时域，表达式简单走变换域}：分段的横线图形卷积只需算三角形面积，而指数乘
  多项式硬积会出错。
\end{enumerate}

\section*{附录
B　公式速查}\label{ux9644ux5f55-b-ux516cux5f0fux901fux67e5}
\addcontentsline{toc}{section}{附录 B　公式速查}

本附录只列结论与成立条件，不写推导。带条件的公式在考场上要把条件抄在结果旁边，条件本身
就是采分点。

\subsection*{B.1
信号与基本运算}\label{b.1-ux4fe1ux53f7ux4e0eux57faux672cux8fd0ux7b97}
\addcontentsline{toc}{subsection}{B.1 信号与基本运算}

\begin{longtable}[]{@{}
  >{\raggedright\arraybackslash}p{(\linewidth - 4\tabcolsep) * \real{0.2727}}
  >{\raggedright\arraybackslash}p{(\linewidth - 4\tabcolsep) * \real{0.2727}}
  >{\raggedright\arraybackslash}p{(\linewidth - 4\tabcolsep) * \real{0.4545}}@{}}
\toprule\noalign{}
\begin{minipage}[b]{\linewidth}\raggedright
项目
\end{minipage} & \begin{minipage}[b]{\linewidth}\raggedright
结论
\end{minipage} & \begin{minipage}[b]{\linewidth}\raggedright
成立条件
\end{minipage} \\
\midrule\noalign{}
\endhead
\bottomrule\noalign{}
\endlastfoot
冲激与阶跃 &
\(\delta(t)=\dfrac{\mathrm{d}u(t)}{\mathrm{d}t}\)，\(u(t)=\displaystyle\int_{-\infty}^{t}\delta(\tau)\mathrm{d}\tau\)
& 广义函数意义下 \\
取样 &
\(\displaystyle\int_{\alpha}^{\beta}f(t)\delta(t-t_{0})\mathrm{d}t=f(t_{0})\)
& \(\alpha<t_{0}<\beta\)，否则为 \(0\) \\
筛选 & \(f(t)\delta(t-t_{0})=f(t_{0})\delta(t-t_{0})\) & 恒成立 \\
冲激尺度 & \(\delta(at)=\dfrac{1}{\lvert a\rvert}\delta(t)\) &
\(a\ne 0\) \\
奇偶分解 &
\(f_{e}(t)=\dfrac{f(t)+f(-t)}{2}\)，\(f_{o}(t)=\dfrac{f(t)-f(-t)}{2}\) &
定义在对称区间 \\
波形变换 &
\(f(at+b)\)：\((t_{k},f_{k})\to\left(\dfrac{t_{k}-b}{a},\ f_{k}\right)\)
& 先定 \(a\) 的符号再搬点 \\
能量与功率 &
\(E=\displaystyle\int_{-\infty}^{+\infty}\lvert f(t)\rvert^{2}\mathrm{d}t\)，\(P=\displaystyle\lim_{T\to\infty}\frac{1}{2T}\int_{-T}^{T}\lvert f(t)\rvert^{2}\mathrm{d}t\)
& 能量信号 \(0<E<\infty\)；功率信号 \(0<P<\infty\) \\
\end{longtable}

波形变换一行的两种写法等价：把括号内化成 \(a(t-t_{0})\)
的标准形后，关键点映射为
\((t_{k},f_{k})\to\left(t_{0}+t_{k}/a,\ Af_{k}\right)\)，其中 \(A\)
为该步的幅度增益； 用 \(f(at+b)\)
的原始形式算，就是上表那一行。作图顺序按''尺度、反褶、平移、标量相乘''，
\textbf{先平移后尺度时平移量要除以
\(a\)}，这是最容易错的一步，见``信号的简单处理''一节。

\subsection*{B.2 卷积}\label{b.2-ux5377ux79ef}
\addcontentsline{toc}{subsection}{B.2 卷积}

定义与因果形式：

\[
f_{1}(t)*f_{2}(t)=\int_{-\infty}^{+\infty}f_{1}(\tau)f_{2}(t-\tau)\,\mathrm{d}\tau .
\]

若 \(f_{1}\)、\(f_{2}\) 都是因果信号，则上下限改为 \(0\) 到 \(t\)，即
\(f_{1}(t)*f_{2}(t)=\displaystyle\int_{0}^{t}f_{1}(\tau)f_{2}(t-\tau)\mathrm{d}\tau\)。

\begin{longtable}[]{@{}
  >{\raggedright\arraybackslash}p{(\linewidth - 4\tabcolsep) * \real{0.3333}}
  >{\raggedright\arraybackslash}p{(\linewidth - 4\tabcolsep) * \real{0.3333}}
  >{\raggedright\arraybackslash}p{(\linewidth - 4\tabcolsep) * \real{0.3333}}@{}}
\toprule\noalign{}
\begin{minipage}[b]{\linewidth}\raggedright
性质
\end{minipage} & \begin{minipage}[b]{\linewidth}\raggedright
结论
\end{minipage} & \begin{minipage}[b]{\linewidth}\raggedright
条件
\end{minipage} \\
\midrule\noalign{}
\endhead
\bottomrule\noalign{}
\endlastfoot
交换、分配、结合 &
\(f_{1}*f_{2}=f_{2}*f_{1}\)，\(f_{1}*(f_{2}+f_{3})=f_{1}*f_{2}+f_{1}*f_{3}\)
& 恒成立 \\
时移 & \(f_{1}(t-t_{1})*f_{2}(t-t_{2})=y(t-t_{1}-t_{2})\) & 恒成立 \\
微积分 &
\(\left(f_{1}*f_{2}\right)'=f_{1}'*f_{2}=f_{1}*f_{2}'\)，积分可任取一个因子做
& 可导 \\
与冲激 &
\(f(t)*\delta(t-t_{0})=f(t-t_{0})\)，\(f(t)*\delta^{(n)}(t-t_{0})=f^{(n)}(t-t_{0})\)
& 恒成立 \\
与阶跃 &
\(f(t)*u(t)=\displaystyle\int_{-\infty}^{t}f(\tau)\mathrm{d}\tau\) &
恒成立 \\
\end{longtable}

常用卷积结果（\(a\) 为任意实数，\(E\) 为幅度，\(\tau\) 为宽度）：

\begin{longtable}[]{@{}
  >{\raggedright\arraybackslash}p{(\linewidth - 2\tabcolsep) * \real{0.5000}}
  >{\raggedright\arraybackslash}p{(\linewidth - 2\tabcolsep) * \real{0.5000}}@{}}
\toprule\noalign{}
\begin{minipage}[b]{\linewidth}\raggedright
卷积
\end{minipage} & \begin{minipage}[b]{\linewidth}\raggedright
结果
\end{minipage} \\
\midrule\noalign{}
\endhead
\bottomrule\noalign{}
\endlastfoot
\(u(t)*u(t)\) & \(t\,u(t)\) \\
\(t\,u(t)*u(t)\) & \(\dfrac{1}{2}t^{2}u(t)\) \\
\(\mathrm{e}^{at}u(t)*\mathrm{e}^{at}u(t)\) &
\(t\,\mathrm{e}^{at}u(t)\)（谐振来源） \\
两个同宽同高矩形 & 三角形，底 \(2\tau\)，高 \(E^{2}\tau\) \\
两个不同宽矩形 & 梯形，总宽 \(\tau_{1}+\tau_{2}\)，平顶宽
\(\lvert\tau_{1}-\tau_{2}\rvert\)，高
\(E^{2}\min(\tau_{1},\tau_{2})\) \\
\end{longtable}

\textbf{卷积积分限的三条判据}（设 \(f_{1}\) 只在 \([a,b]\)
非零，\(f_{2}\) 只在 \([c,d]\) 非零）：

\begin{enumerate}
\def\labelenumi{\arabic{enumi}.}
\tightlist
\item
  重叠区间为
  \(\left[\max(a,t-d),\ \min(b,t-c)\right]\)，区间非空才积，否则结果写
  \(0\)；
\item
  四个分界点 \(a+c\)、\(a+d\)、\(b+c\)、\(b+d\)
  排序后即为分段点，互不相同的个数等于分段段数；
\item
  结果支撑区起点 \(a+c\)、长度 \((b-a)+(d-c)\)，用它回验分段是否写全。
\end{enumerate}

\subsection*{B.3 傅里叶级数}\label{b.3-ux5085ux91ccux53f6ux7ea7ux6570}
\addcontentsline{toc}{subsection}{B.3 傅里叶级数}

本附录统一记基频 \(\omega_{0}=\dfrac{2\pi}{T}\)（第 3 章正文记作
\(\Omega\)，是同一个量）。

三角形式与系数（周期为 \(T\)）：

\[
f(t)=\frac{a_{0}}{2}+\sum_{n=1}^{+\infty}\left(a_{n}\cos n\omega_{0}t+b_{n}\sin n\omega_{0}t\right),
\] \[
a_{n}=\frac{2}{T}\int_{t_{1}}^{t_{1}+T}f(t)\cos n\omega_{0}t\,\mathrm{d}t,
\qquad
b_{n}=\frac{2}{T}\int_{t_{1}}^{t_{1}+T}f(t)\sin n\omega_{0}t\,\mathrm{d}t .
\]

指数形式与换算：

\[
f(t)=\sum_{n=-\infty}^{+\infty}c_{n}\mathrm{e}^{\mathrm{j}n\omega_{0}t},
\qquad
c_{n}=\frac{1}{T}\int_{t_{1}}^{t_{1}+T}f(t)\mathrm{e}^{-\mathrm{j}n\omega_{0}t}\mathrm{d}t .
\]

两种形式的换算（实信号，\(n\ge 1\)）：

\[
c_{0}=\frac{a_{0}}{2}=\frac{A_{0}}{2},
\qquad
\lvert c_{n}\rvert=\frac{A_{n}}{2}=\frac{1}{2}\sqrt{a_{n}^{2}+b_{n}^{2}},
\qquad
\varphi_{n}=-\arctan\frac{b_{n}}{a_{n}} .
\]

\(\varphi_{n}\) 的象限由 \((a_{n},-b_{n})\)
所在象限决定，不能只取反正切主值；双边谱必须成对 画 \(c_{n}\) 与
\(c_{-n}\)。

周期信号的平均功率与帕塞瓦尔定理（\(f\) 在一周期内平方可积）：

\[
P=\frac{1}{T}\int_{t_{1}}^{t_{1}+T}f^{2}(t)\,\mathrm{d}t
=\frac{a_{0}^{2}}{4}+\sum_{n=1}^{+\infty}\frac{a_{n}^{2}+b_{n}^{2}}{2}
=\sum_{n=-\infty}^{+\infty}\lvert c_{n}\rvert^{2},
\qquad
f_{\mathrm{rms}}=\sqrt{P} .
\]

功率可以按频率逐项相加，有效值不行------这是帕塞瓦尔定理最常被用错的地方。

成立条件：满足狄利克雷条件（一周期内只有有限个第一类间断点与有限个极值点，且绝对可积）；
此时级数在连续点等于 \(f(t)\)，在间断点收敛到
\(\dfrac{f(t_{0}^{+})+f(t_{0}^{-})}{2}\)。

必备模型（周期矩形脉冲，幅度 \(E\)、脉宽 \(\tau\)、周期 \(T\)）：

\[
c_{n}=\frac{E\tau}{T}\,\mathrm{Sa}\left(\frac{n\omega_{0}\tau}{2}\right),
\qquad
\mathrm{Sa}(0)=1 .
\]

频带宽度三种定义，题目必须认准是哪一种：以最大幅度的 \(\frac{1}{10}\)
为限；以幅度谱第一个 过零点为限（最常用）；以包含 \(90\%\)
功率为限。按第二种定义，矩形脉冲
\(B_{f}=\dfrac{1}{\tau}\)、\(B_{\omega}=\dfrac{2\pi}{\tau}\)，即脉宽与带宽成反比。

\subsection*{B.4 傅里叶变换}\label{b.4-ux5085ux91ccux53f6ux53d8ux6362}
\addcontentsline{toc}{subsection}{B.4 傅里叶变换}

\[
F(\mathrm{j}\omega)=\int_{-\infty}^{+\infty}f(t)\mathrm{e}^{-\mathrm{j}\omega t}\mathrm{d}t,
\qquad
f(t)=\frac{1}{2\pi}\int_{-\infty}^{+\infty}F(\mathrm{j}\omega)\mathrm{e}^{\mathrm{j}\omega t}\mathrm{d}\omega .
\]

反变换前的 \(\dfrac{1}{2\pi}\) 只属于反变换，正变换没有。

存在条件：\(f(t)\)
在整条实轴上绝对可积，且只有有限个第一类间断点与有限个极值点
（狄利克雷条件，充分条件）。不满足时，借助冲激函数可定义广义傅里叶变换。

\begin{longtable}[]{@{}
  >{\raggedright\arraybackslash}p{(\linewidth - 4\tabcolsep) * \real{0.3333}}
  >{\raggedright\arraybackslash}p{(\linewidth - 4\tabcolsep) * \real{0.3333}}
  >{\raggedright\arraybackslash}p{(\linewidth - 4\tabcolsep) * \real{0.3333}}@{}}
\toprule\noalign{}
\begin{minipage}[b]{\linewidth}\raggedright
性质
\end{minipage} & \begin{minipage}[b]{\linewidth}\raggedright
结论
\end{minipage} & \begin{minipage}[b]{\linewidth}\raggedright
条件
\end{minipage} \\
\midrule\noalign{}
\endhead
\bottomrule\noalign{}
\endlastfoot
时移与频移 &
\(f(t-t_{0})\leftrightarrow F\mathrm{e}^{-\mathrm{j}\omega t_{0}}\)；\(f(t)\mathrm{e}^{\mathrm{j}\omega_{0}t}\leftrightarrow F\bigl(\mathrm{j}(\omega-\omega_{0})\bigr)\)
& 恒成立 \\
尺度 &
\(f(at)\leftrightarrow\dfrac{1}{\lvert a\rvert}F\left(\mathrm{j}\dfrac{\omega}{a}\right)\)
& \(a\ne 0\)，\(a<0\) 带反褶 \\
对称 & \(F(\mathrm{j}t)\leftrightarrow 2\pi f(-\omega)\) & 恒成立 \\
微分 &
\(f^{(n)}(t)\leftrightarrow(\mathrm{j}\omega)^{n}F(\mathrm{j}\omega)\) &
可用跳变逐项写 \\
积分 &
\(\displaystyle\int_{-\infty}^{t}f\leftrightarrow\frac{F(\mathrm{j}\omega)}{\mathrm{j}\omega}+\pi F(0)\delta(\omega)\)
& \(F(0)\ne 0\) 时冲激项必写 \\
卷积定理 & \(f_{1}*f_{2}\leftrightarrow F_{1}F_{2}\) & 恒成立 \\
相乘定理 & \(f_{1}f_{2}\leftrightarrow\dfrac{1}{2\pi}F_{1}*F_{2}\) &
因子不能丢 \\
\end{longtable}

调制与抽样两条最常用的推论：

\[
f(t)\cos\omega_{0}t\ \longleftrightarrow\
\frac{1}{2}\left[F\bigl(\mathrm{j}(\omega+\omega_{0})\bigr)+F\bigl(\mathrm{j}(\omega-\omega_{0})\bigr)\right],
\]

\[
F_{s}(\mathrm{j}\omega)=\frac{1}{T_{s}}\sum_{k=-\infty}^{+\infty}F\bigl(\mathrm{j}(\omega-k\omega_{s})\bigr).
\]

恢复条件两条：抽样角频率
\(\omega_{s}\ge 2\omega_{m}\)（奈奎斯特），恢复低通的截止频率
\(\omega_{m}<\omega_{c}<\omega_{s}-\omega_{m}\)。

周期信号的傅里叶变换（冲激串）：

\[
F(\mathrm{j}\omega)=2\pi\sum_{n=-\infty}^{+\infty}c_{n}\delta(\omega-n\omega_{0}),
\qquad
\sum_{k=-\infty}^{+\infty}\delta(t-kT)\ \longleftrightarrow\
\omega_{0}\sum_{n=-\infty}^{+\infty}\delta(\omega-n\omega_{0}).
\]

单脉冲与周期化的桥：\(c_{n}=\dfrac{1}{T}F_{0}(\mathrm{j}n\omega_{0})\)，\(F_{0}\)
是取一个周期、
其余补零所得信号的变换。能量式（帕塞瓦尔）：\(E=\displaystyle\int\lvert f\rvert^{2}\mathrm{d}t
=\frac{1}{2\pi}\int\lvert F(\mathrm{j}\omega)\rvert^{2}\mathrm{d}\omega\)，要求
\(f\) 能量有限。

完整变换对见 A.3；判决''能否用频域法''的依据是上面这条存在条件。

\subsection*{B.5
拉普拉斯变换}\label{b.5-ux62c9ux666eux62c9ux65afux53d8ux6362}
\addcontentsline{toc}{subsection}{B.5 拉普拉斯变换}

单边与双边（必须分清，混用是硬错）：

单边拉普拉斯变换：

\[
F(s)=\int_{0^{-}}^{+\infty}f(t)\mathrm{e}^{-st}\mathrm{d}t .
\]

双边拉普拉斯变换：

\[
F(s)=\int_{-\infty}^{+\infty}f(t)\mathrm{e}^{-st}\mathrm{d}t .
\]

下限取 \(0^{-}\) 是为了把 \(t=0\) 处的冲激与 \(t=0^{-}\)
的初始状态一并计入，这样微分性质的初始
项才与电路''动作发生前''的状态对齐。单边变换只反映 \(t\ge 0\) 部分，因此
\(F(s)\) 唯一确定 \(t\ge 0\) 的波形；双边变换必须 \(F(s)\)
与收敛域一起给才唯一。

收敛域与信号类型的对应：

\begin{longtable}[]{@{}
  >{\raggedright\arraybackslash}p{(\linewidth - 4\tabcolsep) * \real{0.3846}}
  >{\raggedright\arraybackslash}p{(\linewidth - 4\tabcolsep) * \real{0.4615}}
  >{\raggedright\arraybackslash}p{(\linewidth - 4\tabcolsep) * \real{0.1538}}@{}}
\toprule\noalign{}
\begin{minipage}[b]{\linewidth}\raggedright
信号类型
\end{minipage} & \begin{minipage}[b]{\linewidth}\raggedright
收敛域形状
\end{minipage} & \begin{minipage}[b]{\linewidth}\raggedright
例
\end{minipage} \\
\midrule\noalign{}
\endhead
\bottomrule\noalign{}
\endlastfoot
右边（因果） & \(\mathrm{Re}(s)>\sigma_{0}\) &
\(\mathrm{e}^{-2t}u(t)\)：\(\mathrm{Re}(s)>-2\) \\
左边（反因果） & \(\mathrm{Re}(s)<\sigma_{0}\) &
\(-\mathrm{e}^{-2t}u(-t)\)：\(\mathrm{Re}(s)<-2\) \\
双边 & \(\sigma_{1}<\mathrm{Re}(s)<\sigma_{2}\) & 两项之和，条带形 \\
时限且绝对可积 & 整个 \(s\) 平面 & 因果门函数 \\
\end{longtable}

三条推论：收敛域边界必穿过极点且内部不含极点；收敛域含虚轴时
\(F(\mathrm{j}\omega)=\left.F(s)\right\rvert_{s=\mathrm{j}\omega}\)；同一个有理式配不同收敛域
给出不同原信号。

单边拉普拉斯变换的性质（与傅里叶变换逐条对照见 A.2）：

\begin{longtable}[]{@{}
  >{\raggedright\arraybackslash}p{(\linewidth - 4\tabcolsep) * \real{0.3000}}
  >{\raggedright\arraybackslash}p{(\linewidth - 4\tabcolsep) * \real{0.4000}}
  >{\raggedright\arraybackslash}p{(\linewidth - 4\tabcolsep) * \real{0.3000}}@{}}
\toprule\noalign{}
\begin{minipage}[b]{\linewidth}\raggedright
性质
\end{minipage} & \begin{minipage}[b]{\linewidth}\raggedright
复频域
\end{minipage} & \begin{minipage}[b]{\linewidth}\raggedright
条件
\end{minipage} \\
\midrule\noalign{}
\endhead
\bottomrule\noalign{}
\endlastfoot
线性 & \(a_{1}F_{1}(s)+a_{2}F_{2}(s)\) & 至少为两收敛域的交 \\
尺度变换 & \(\dfrac{1}{a}F\left(\dfrac{s}{a}\right)\) & \(a>0\) \\
时移 & \(\mathrm{e}^{-st_{0}}F(s)\) & 必须写 \(f(t-t_{0})u(t-t_{0})\) \\
频移 & \(F(s-s_{0})\) & 收敛域平移 \(\mathrm{Re}(s_{0})\) \\
时域微分 & \(s^{n}F(s)\) 减去初始项 & 初始项不可省，见下式 \\
时域积分 & \(\dfrac{F(s)}{s}\) & 下限取 \(0^{-}\) \\
频域微分 & \(-\dfrac{\mathrm{d}F(s)}{\mathrm{d}s}\) 对应 \(tf(t)\) &
收敛域不变 \\
卷积定理 & \(F_{1}(s)F_{2}(s)\) & 至少为两收敛域的交 \\
\end{longtable}

\[
\mathcal{L}\left[f^{(n)}(t)\right]
=s^{n}F(s)-\sum_{k=0}^{n-1}s^{n-1-k}f^{(k)}(0^{-}).
\]

常用变换对见 A.3，这里只补与初值、终值有关的结论。

初值定理（\(f\) 与 \(f'\)
存在且变换存在）：\(f(0^{+})=\lim\limits_{s\to\infty}sF(s)\)。若
\(F(s)\) 是假分式（含冲激项），先长除分出多项式，再对真分式部分取极限。

终值定理（成立条件：\(sF(s)\)
的极点全在左半平面）：\(f(+\infty)=\lim\limits_{s\to 0}sF(s)\)。
虚轴上有极点时终值不存在，算出来的是假值。阶跃输入下
\(Y(s)=\dfrac{H(s)}{s}\)，故 \(y(+\infty)=H(0)\)。

\textbf{部分分式展开的全部系数公式}（设
\(F(s)=\dfrac{N(s)}{D(s)}\)，\(m<n\)）：

\begin{longtable}[]{@{}
  >{\raggedright\arraybackslash}p{(\linewidth - 4\tabcolsep) * \real{0.2500}}
  >{\raggedright\arraybackslash}p{(\linewidth - 4\tabcolsep) * \real{0.4167}}
  >{\raggedright\arraybackslash}p{(\linewidth - 4\tabcolsep) * \real{0.3333}}@{}}
\toprule\noalign{}
\begin{minipage}[b]{\linewidth}\raggedright
情形
\end{minipage} & \begin{minipage}[b]{\linewidth}\raggedright
系数公式
\end{minipage} & \begin{minipage}[b]{\linewidth}\raggedright
时间项
\end{minipage} \\
\midrule\noalign{}
\endhead
\bottomrule\noalign{}
\endlastfoot
单实根 &
\(K_{i}=(s-s_{i})F(s)\big\rvert_{s=s_{i}}=\dfrac{N(s_{i})}{D'(s_{i})}\)
& \(K_{i}\mathrm{e}^{s_{i}t}u(t)\) \\
共轭复根 &
\(K_{2}=K_{1}^{*}\)，\(K_{1}=\lvert K_{1}\rvert\mathrm{e}^{\mathrm{j}\phi_{1}}\)
&
\(2\lvert K_{1}\rvert\mathrm{e}^{\sigma_{1}t}\cos(\omega_{0}t+\phi_{1})\) \\
\(r\) 重根 \(p\) &
\(K_{1k}=\dfrac{1}{(r-k)!}\left.\dfrac{\mathrm{d}^{r-k}}{\mathrm{d}s^{r-k}}\left[(s-p)^{r}F(s)\right]\right\rvert_{s=p}\)
& \(\dfrac{K_{1k}t^{r-k}}{(r-k)!}\mathrm{e}^{pt}u(t)\) \\
假分式 & 先长除：\(F(s)=M(s)+\dfrac{N_{1}(s)}{D(s)}\) & \(M(s)\) 给
\(\delta(t)\) 及其导数 \\
\end{longtable}

假分式一项务必单独记：\(\mathcal{L}^{-1}[1]=\delta(t)\)，\(\mathcal{L}^{-1}[s^{n}]=\delta^{(n)}(t)\)。
重根系数一定要逐阶求导，\(r=2\) 时两个系数分别是
\(\left[(s-p)^{2}F(s)\right]_{s=p}\) 与它的一阶导数在 \(s=p\) 的值。

\subsection*{B.6
系统函数与稳定性}\label{b.6-ux7cfbux7edfux51fdux6570ux4e0eux7a33ux5b9aux6027}
\addcontentsline{toc}{subsection}{B.6 系统函数与稳定性}

\[
H(s)=\frac{Y_{zs}(s)}{F(s)}=\frac{B(s)}{A(s)},
\qquad
H(\mathrm{j}\omega)=\left.H(s)\right\rvert_{s=\mathrm{j}\omega}.
\]

后一个等号只在收敛域含虚轴（即系统稳定）时成立，这是
\(H(\mathrm{j}\omega)\) 与 \(H(s)\) 唯一 正式的关系。

\(A(s)=0\) 的根为极点 \(p_{i}\)，\(B(s)=0\) 的根为零点
\(z_{i}\)；实系数有理式的极零点只能是实数
或共轭复数，极零图关于实轴对称。极点个数与零点个数之差决定 \(h(t)\)
是否含 \(\delta(t)\) 及其导数（\(m\ge n\) 时含）。

极点位置与 \(h(t)\) 波形（\(p_{i}=\sigma_{i}+\mathrm{j}\omega_{i}\)）：

\begin{longtable}[]{@{}
  >{\raggedright\arraybackslash}p{(\linewidth - 4\tabcolsep) * \real{0.3448}}
  >{\raggedright\arraybackslash}p{(\linewidth - 4\tabcolsep) * \real{0.2069}}
  >{\raggedright\arraybackslash}p{(\linewidth - 4\tabcolsep) * \real{0.4483}}@{}}
\toprule\noalign{}
\begin{minipage}[b]{\linewidth}\raggedright
极点位置
\end{minipage} & \begin{minipage}[b]{\linewidth}\raggedright
参数
\end{minipage} & \begin{minipage}[b]{\linewidth}\raggedright
\(h(t)\) 行为
\end{minipage} \\
\midrule\noalign{}
\endhead
\bottomrule\noalign{}
\endlastfoot
负实轴 & \(\sigma_{i}<0\)，\(\omega_{i}=0\) & 单调衰减，时间常数
\(1/\lvert\sigma_{i}\rvert\) \\
正实轴 & \(\sigma_{i}>0\)，\(\omega_{i}=0\) & 单调增长 \\
左半平面共轭对 & \(\sigma_{i}<0\)，\(\omega_{i}\ne 0\) & 减幅振荡，周期
\(2\pi/\lvert\omega_{i}\rvert\) \\
虚轴共轭对 & \(\sigma_{i}=0\)，\(\omega_{i}\ne 0\) &
等幅振荡，临界稳定 \\
右半平面共轭对 & \(\sigma_{i}>0\)，\(\omega_{i}\ne 0\) & 增幅振荡；\(r\)
重极点为发散 \\
\end{longtable}

稳定判据三级：必要条件是 \(D(s)\)
系数同号且无缺项；充要条件是罗斯表首列同号（等效于极点
全在左半平面）；因果且稳定等价于极点全在左半平面。反馈系统
\(H(s)=\dfrac{G(s)}{1+G(s)K(s)}\)， 闭环极点由特征方程 \(1+G(s)K(s)=0\)
决定。

罗斯表递推（把上两行记作 \(b_{1},b_{2},\dots\) 与
\(c_{1},c_{2},\dots\)）：

\[
d_{1}=\frac{b_{1}c_{2}-b_{2}c_{1}}{b_{1}},
\qquad
d_{2}=\frac{b_{1}c_{3}-b_{2}c_{2}}{b_{1}},
\qquad\cdots
\]

口诀：左上乘右下，减右上乘左下，再除以上一行的首列元素。首列变号次数等于右半平面根个数
（重根按重数计）。

两种特殊情况的处理规则：

\begin{enumerate}
\def\labelenumi{\arabic{enumi}.}
\tightlist
\item
  \textbf{首列出现零}：用无穷小正量 \(\varepsilon\) 代替该零继续造表，令
  \(\varepsilon\to 0^{+}\) 判符号；或先把 \(D(s)\) 乘 \((s+1)\)
  再重建表。两种做法结论必须一致，卷面上写清用的是哪种。
\item
  \textbf{出现全零行}：说明存在关于原点对称的根。取全零行的上一行系数构造辅助多项式
  \(P(s)\)， 用 \(\dfrac{\mathrm{d}P(s)}{\mathrm{d}s}\)
  的系数替代全零行继续造表；再解 \(P(s)=0\)：虚轴上
  单极点判临界稳定，重极点判不稳定。
\end{enumerate}

\subsection*{B.7
常用数学工具}\label{b.7-ux5e38ux7528ux6570ux5b66ux5de5ux5177}
\addcontentsline{toc}{subsection}{B.7 常用数学工具}

积化和差（级数展开与卷积化简都要用）：

\begin{longtable}[]{@{}ll@{}}
\toprule\noalign{}
公式 & 结果 \\
\midrule\noalign{}
\endhead
\bottomrule\noalign{}
\endlastfoot
\(\sin A\cos B\) & \(\dfrac{1}{2}\left[\sin(A+B)+\sin(A-B)\right]\) \\
\(\cos A\cos B\) & \(\dfrac{1}{2}\left[\cos(A+B)+\cos(A-B)\right]\) \\
\(\sin A\sin B\) & \(\dfrac{1}{2}\left[\cos(A-B)-\cos(A+B)\right]\) \\
\end{longtable}

二倍角与降幂：\(\sin 2A=2\sin A\cos A\)，
\(\cos 2A=\cos^{2}A-\sin^{2}A=1-2\sin^{2}A\)，
\(\sin^{2}A=\dfrac{1-\cos 2A}{2}\)，\(\cos^{2}A=\dfrac{1+\cos 2A}{2}\)。

欧拉公式与复数：

\[
\mathrm{e}^{\mathrm{j}\theta}=\cos\theta+\mathrm{j}\sin\theta,
\qquad
\cos\theta=\frac{\mathrm{e}^{\mathrm{j}\theta}+\mathrm{e}^{-\mathrm{j}\theta}}{2},
\qquad
\sin\theta=\frac{\mathrm{e}^{\mathrm{j}\theta}-\mathrm{e}^{-\mathrm{j}\theta}}{2\mathrm{j}} .
\]

\[
a+\mathrm{j}b=\sqrt{a^{2}+b^{2}}\,\mathrm{e}^{\mathrm{j}\varphi},
\qquad
\sum_{n=0}^{N-1}q^{n}=\frac{1-q^{N}}{1-q}\quad (q\ne 1).
\]

其中 \(\varphi=\arctan\dfrac{b}{a}\)，实数部分为负时要按 \((a,b)\)
所在象限加或减 \(\pi\)，不能
只取反正切主值。等比求和在周期信号展开与冲激串求和里反复出现。

常用定积分（\(a>0\)，\(\omega\) 任意）：

\begin{longtable}[]{@{}
  >{\raggedright\arraybackslash}p{(\linewidth - 2\tabcolsep) * \real{0.5000}}
  >{\raggedright\arraybackslash}p{(\linewidth - 2\tabcolsep) * \real{0.5000}}@{}}
\toprule\noalign{}
\begin{minipage}[b]{\linewidth}\raggedright
积分
\end{minipage} & \begin{minipage}[b]{\linewidth}\raggedright
结果
\end{minipage} \\
\midrule\noalign{}
\endhead
\bottomrule\noalign{}
\endlastfoot
\(\displaystyle\int_{0}^{\infty}t\mathrm{e}^{-at}\mathrm{d}t\) &
\(\dfrac{1}{a^{2}}\) \\
\(\displaystyle\int_{0}^{\infty}t^{n}\mathrm{e}^{-at}\mathrm{d}t\) &
\(\dfrac{n!}{a^{n+1}}\) \\
\(\displaystyle\int_{0}^{\infty}\mathrm{e}^{-at}\cos\omega t\,\mathrm{d}t\)
& \(\dfrac{a}{a^{2}+\omega^{2}}\) \\
\(\displaystyle\int_{0}^{\infty}\mathrm{e}^{-at}\sin\omega t\,\mathrm{d}t\)
& \(\dfrac{\omega}{a^{2}+\omega^{2}}\) \\
\(\displaystyle\int_{0}^{\infty}t\mathrm{e}^{-at}\cos\omega t\,\mathrm{d}t\)
& \(\dfrac{a^{2}-\omega^{2}}{(a^{2}+\omega^{2})^{2}}\) \\
\(\displaystyle\int_{0}^{\infty}t\mathrm{e}^{-at}\sin\omega t\,\mathrm{d}t\)
& \(\dfrac{2a\omega}{(a^{2}+\omega^{2})^{2}}\) \\
\end{longtable}

抽样函数的常用值（\(\mathrm{Sa}(x)=\dfrac{\sin x}{x}\)，与
\(\mathrm{sinc}\) 相差一个 \(\pi\)
因子）：\(\mathrm{Sa}(0)=1\)；\(\mathrm{Sa}(n\pi)=0\)（\(n\)
为非零整数）；
\(\displaystyle\int_{-\infty}^{+\infty}\mathrm{Sa}(x)\mathrm{d}x=\pi\)；
\(\displaystyle\int_{0}^{+\infty}\mathrm{Sa}(x)\mathrm{d}x=\frac{\pi}{2}\)；
\(\displaystyle\int_{-\infty}^{+\infty}\mathrm{Sa}^{2}(x)\mathrm{d}x=\pi\)。

\section*{附录
C　易错点与陷阱总汇}\label{ux9644ux5f55-c-ux6613ux9519ux70b9ux4e0eux9677ux9631ux603bux6c47}
\addcontentsline{toc}{section}{附录 C　易错点与陷阱总汇}

六章正文的''易错点''节按章罗列，这里按性质重排成三类。每条都是考场提醒：前半句是错误做法，
后半句是正确做法，括号里是一句话原因。

\subsection*{C.1 概念类}\label{c.1-ux6982ux5ff5ux7c7b}
\addcontentsline{toc}{subsection}{C.1 概念类}

\begin{enumerate}
\def\labelenumi{\arabic{enumi}.}
\tightlist
\item
  错误：把 \(y_{zi}(t)\) 当自然响应。正确：自然响应是 \(y_{zi}\) 加上
  \(y_{zs}\) 中与特征根
  对应的那部分（两者划分依据不同：谁引起对极点归属）。
\item
  错误：把 \(y_{zs}(t)\) 当受迫响应。正确：\(y_{zs}\)
  同时含自然与受迫两种成分（它由全部 极点决定，不只看激励极点）。
\item
  错误：把''稳定''等同于''渐近稳定''。正确：BIBO 稳定只要求 \(h(t)\)
  绝对可积（虚轴一阶极点 属临界稳定，此时 \(h(t)\) 不趋于零）。
\item
  错误：以为 \(H(\mathrm{j}\omega)\) 与 \(H(s)\)
  是一回事。正确：\(H(\mathrm{j}\omega)\) 只是 \(H(s)\)
  在虚轴上的截线，且要求收敛域含虚轴（系统稳定）。
\item
  错误：以为任何系统都能令
  \(s=\mathrm{j}\omega\)。正确：只有稳定系统的频率响应才有意义
  （右半平面极点使 \(h\) 不绝对可积）。
\item
  错误：以为 \(h(t)\) 是系统对任意输入的响应。正确：\(h(t)\) 是激励为
  \(\delta(t)\) 时的零状态 响应（它描述系统本身）。
\item
  错误：以为周期信号没有傅里叶变换。正确：周期信号的频谱是冲激串
  \(2\pi\sum c_{n}\delta(\omega-n\omega_{0})\)（广义变换）。
\item
  错误：认为负频率''不存在所以不必画''。正确：双边谱必须画负频率，实信号的双边幅度谱偶
  对称（漏掉就丢一半幅度）。
\item
  错误：把取样特性的结果当成 \(t\) 的函数。正确：结果是数
  \(f(t_{0})\)（积分已把 \(t\) 消掉）。
\item
  错误：以为子系统串联时 \(H\)
  相加。正确：串联相乘、并联相加（对应卷积与求和）。
\item
  错误：以为全通系统就是无失真系统。正确：全通只保证幅频为常数，相频非线性时仍有相位
  失真（无失真要求两条同时满足）。
\item
  错误：以为负反馈一定稳定。正确：增益过大会把闭环极点推入右半平面（这正是求
  \(K\) 稳定 范围的考点）。
\item
  错误：以为频带越宽响应越慢。正确：\(t_{r}B_{f}\approx 0.5\)，带宽与上升时间成反比。
\item
  错误：以为系数全正就稳定。正确：同号只是必要条件，必须再列罗斯表（见``罗斯-霍尔维茨判据''一节的反例）。
\end{enumerate}

\subsection*{C.2 运算类}\label{c.2-ux8fd0ux7b97ux7c7b}
\addcontentsline{toc}{subsection}{C.2 运算类}

\begin{enumerate}
\def\labelenumi{\arabic{enumi}.}
\tightlist
\item
  错误：卷积一律写 \(\int_{-\infty}^{+\infty}\)
  不分段。正确：先定支撑区，用四个分界点划区间
  （区间外被积函数为零，硬套全域得不到分段结果）。
\item
  错误：反褶后把平移方向写反。正确：\(h(t-\tau)\) 是把 \(h(-\tau)\)
  向右平移 \(t\) （\(t>0\)），写成 \(h(-(\tau+t))\) 就变成左移。
\item
  错误：部分分式重根系数只求一次。正确：按
  \(\dfrac{1}{(r-k)!}\dfrac{\mathrm{d}^{r-k}}{\mathrm{d}s^{r-k}}\)
  逐阶求导（\(r\) 阶重根要 求 \(r\) 次）。
\item
  错误：对假分式直接做部分分式。正确：先长除分出多项式（多项式项对应
  \(\delta(t)\) 及其 导数）。
\item
  错误：微分性质只写 \(s^{2}Y(s)\)。正确：写全
  \(s^{2}Y(s)-sy(0^{-})-y'(0^{-})\)（初始项 是全响应一次算出的根本）。
\item
  错误：时移性质写成 \(f(t-t_{0})\)。正确：单边变换要写
  \(f(t-t_{0})u(t-t_{0})\)（\(t<t_{0}\) 段必须确实为零）。
\item
  错误：反变换漏
  \(\dfrac{1}{2\pi}\)。正确：该因子只属于傅里叶反变换（正变换没有）。
\item
  错误：频域卷积漏
  \(\dfrac{1}{2\pi}\)。正确：相乘定理必须带该因子（与卷积定理不对称）。
\item
  错误：冲激尺度写成
  \(\delta(at)=\delta(t)\)。正确：\(\delta(at)=\dfrac{1}{\lvert a\rvert}
  \delta(t)\)（面积必须守 1）。
\item
  错误：尺度变换顺序做错。正确：先平移后压缩时平移量要除以
  \(a\)（每一步都对变量 \(t\) 操作）。
\item
  错误：罗斯表整行为零时直接跳过。正确：用辅助多项式求导后的系数替代该行（全零行说明
  有对称根，不能直接判稳定或不定）。
\item
  错误：罗斯表首列是零还继续除。正确：用 \(\varepsilon\) 代替或乘
  \((s+1)\) 重建（除零无 意义）。
\item
  错误：变号次数跳行比较。正确：只比较首列相邻两元素的正负（跳行会漏计）。
\item
  错误：单边频谱直接抄双边幅度。正确：\(n\ge 1\) 时
  \(A_{n}=2\lvert c_{n}\rvert\)（单边把 负频率的能量折叠过来）。
\item
  错误：把 \(a_{0}\) 当直流分量。正确：直流分量是
  \(\dfrac{a_{0}}{2}=c_{0}\)（展开式自带 \(\frac{1}{2}\)）。
\item
  错误：求 \(H(s)\)
  时不置零初始状态。正确：先把等效源置零（否则得到的是特定初态下的
  响应）。
\item
  错误：反变换结果不写 \(u(t)\)。正确：每个指数项后面都补
  \(u(t)\)（零状态响应在 \(t<0\) 为 零）。
\item
  错误：抽样定理把 \(f_{s}\) 与 \(\omega_{s}\)
  混用。正确：\(f_{s}\ge 2f_{m}\) 与 \(\omega_{s}\ge 2\omega_{m}\)
  各自成立（两个 \(2\) 都对，单位不同）。
\item
  错误：不查条件就用终值定理。正确：先看 \(sF(s)\)
  的极点是否全在左半平面（虚轴有极点时 终值不存在）。
\item
  错误：卷积性质与乘法律混用。正确：只用
  \(\left(f_{1}*f_{2}\right)'=f_{1}'*f_{2}=f_{1}*f_{2}'\)
  这一条恒等式（不能把导数分别 分配后各自求值再相乘）。
\end{enumerate}

\subsection*{C.3
记号与符号类}\label{c.3-ux8bb0ux53f7ux4e0eux7b26ux53f7ux7c7b}
\addcontentsline{toc}{subsection}{C.3 记号与符号类}

\begin{enumerate}
\def\labelenumi{\arabic{enumi}.}
\tightlist
\item
  错误：\(\mathcal{F}\) 与 \(\mathcal{L}\) 混用。正确：傅里叶变换写
  \(F(\mathrm{j}\omega)=\mathcal{F}[f(t)]\)，拉氏写
  \(F(s)=\mathcal{L}[f(t)]\)（自变量不同）。
\item
  错误：虚数单位写成斜体字母。正确：一律写正体
  \(\mathrm{j}\)（与支路电流 \(i\) 区分）。
\item
  错误：写 \(F(s)\) 不写收敛域。正确：\(F(s)\) 后紧跟
  \(\mathrm{Re}(s)>\sigma_{0}\)（双边变换 尤其必需）。
\item
  错误：单边变换下限写成 \(0\) 或 \(0^{+}\)。正确：一律 \(0^{-}\)（要含
  \(t=0\) 处的冲激与初态）。
\item
  错误：\(\mathrm{Sa}\) 与 \(\mathrm{sinc}\)
  直接换用。正确：\(\mathrm{Sa}(x)=\frac{\sin x}{x}\)，
  \(\mathrm{sinc}(x)=\frac{\sin\pi x}{\pi x}\)（差一个 \(\pi\) 因子）。
\item
  错误：阶跃写成 \(\varepsilon(t)\)。正确：全篇写
  \(u(t)\)（课件写法与本书记号只差字体）。
\item
  错误：卷积写成点乘或 \(\ast\)。正确：写
  \(f_{1}(t)*f_{2}(t)\)（点乘是相乘，不是卷积）。
\item
  错误：自然底写成斜体 \(e\)，或与电荷 \(e\) 混用。正确：写正体
  \(\mathrm{e}^{-at}\) （指数基数与变量要分清）。
\item
  错误：\(\omega_{c}\) 与 \(f_{c}\) 混写。正确：\(\omega\) 的单位是
  \(\mathrm{rad/s}\)，\(f\) 的 单位是 \(\mathrm{Hz}\)，二者差 \(2\pi\)。
\item
  错误：基频一会儿写 \(\omega_{0}\) 一会儿写
  \(\Omega\)。正确：同一次作答只用一种（本书第 3 章用
  \(\Omega\)，其余章用 \(\omega_{0}\)）。
\item
  错误：响应记号 \(r(t)\) 与 \(y(t)\)
  混用。正确：同一题内统一（前言约定响应为 \(y(t)\)， 零输入、零状态写作
  \(y_{zi}\)、\(y_{zs}\)）。
\item
  错误：把''增量线性''直接写成''线性''。正确：写''严格意义非线性，但零状态响应满足线性''
  （含非零常数项时必须这样表述）。
\end{enumerate}

\end{document}
